\begin{savequote}
{\tiny Ein Mann ging von Jerusalem nach Jericho hinab und wurde von
Räubern überfallen. Sie plünderten ihn aus und schlugen ihn nieder; dann gingen sie weg und ließen ihn 
halbtot liegen. Zufällig kam ein Priester denselben Weg
herab; er sah ihn und ging weiter. Auch ein Levit kam zu
der  Stelle; er sah ihn und ging weiter. Dann kam ein Mann
aus  Samarien, der auf der Reise war. Als er ihn sah, hatte
er  Mitleid, ging zu ihm hin, goss Öl und Wein auf seine
Wunden  und verband sie. 

Dann hob er ihn auf sein Reittier,
brachte  ihn zu einer Herberge und sorgte für ihn. Am
anderen Morgen  holte er zwei Denare hervor, gab sie dem
Wirt und sagte:  Sorge für ihn, und wenn du mehr für ihn
brauchst, werde  ich es dir bezahlen, wenn ich wiederkomme.
Was meinst du:  Wer von diesen dreien hat sich als der
Nächste dessen erwiesen,  der von den Räubern überfallen
wurde?  

Der Gesetzeslehrer antwortete: Der, der
barmherzig\index{Barmherzigkeit} an  ihm gehandelt hat. Da
sagte Jesus zu ihm: Dann geh und handle genauso!

Lk 10,25-37}
% \qauthor{{\tiny Lk 10,25-37}}
\end{savequote}

\chapter[{[MED]} Klassische Medizinische Herausforderungen]
{[MED]\\Klassische Medizinische Herausforderungen}
\label{chp:MED}
% \AasRsN{RS VI}

\emph{Maintainer: Sebastian Gabriel}

\minitoc


% \begin{verbatim}
% 
% Changelog:
% 
% 2009-03-14: In LaTeX überführt
% 
% 2009-02-25: Start Changelog. Bis inkl. bakt. Meningitis mit Korr. von
% Koch verglichen und korr.
% \end{verbatim}




\clearpage
\section{Herz-Kreislaufstörungen}

\dictum[Goethe]{Blut ist ein ganz besond'rer Saft.}

\label{sec:herz-kreislauf-stoerungen}
\label{topic:herz-kreislauf-stoerungen}

% \paragraph{Überblick über die wichtigsten
% Herz-Kreislauf-Notfälle}
% 
% \begin{itemize}
%     \item Herz-/Kreislaufstillstand
%     \item Herzinsuffizienz
%     \begin{itemize}
%         \item Linksherz-Versagen
%         \item Rechtsherz-Versagen
%     \end{itemize}
%     \item (Akutes) Koronarsyndrom
%     \begin{itemize}
%         \item Angina pectoris
%         \item Myokardinfarkt (MCI )
%     \end{itemize}
%     \item Herzrhythmusstörungen
%     \item Störungen des Bludrucks
%     \begin{itemize}
%         \item Hyp\Es{o}tonie
%         \item Hyp\Es{er}tonie
%         \begin{itemize}
%             \item Hypertensive Krise
%             \item Hypertensiver Notfall
%         \end{itemize}
%     \end{itemize}
% \end{itemize}


\subsection{Herz-Kreislauf-Stillstand}

Der Herz-Kreislauf-Stillstand wird
unter
\prettyref{topic:kreislaufstillstand}
und \prettyref{topic:reanimation-eh} besprochen.



\subsection{Herzinsuffizienz}
\label{topic:herzinsuffizienz}
\index{Herzinsuffizienz}
\index{Herzinsuffizienz!kompensierte}
\index{Herzinsuffizienz!dekompensierte}
\index{Rechtsherzinsuffizienz}
\index{Linksherzinsuffizienz}
\index{Globalinsuffizienz}
\index{Insuffizienz!Herz-}


\Definition{Unzureichende Auswurfleistung des Herzens.}

Eine Leistungseinschränkung des Herzmuskels führt zu einer verminderten
Auswurfleistung, es sind viele Ursachen möglich. Manche Ursachen treten
\E{akut} auf, manche sind eher \E{chronischer} Natur.  



\Layout{0}{Einteilung}
{
\subparagraph{Herzteil} 
Je nachdem welcher Herzteil betroffen ist unterteilt man in:

\begin{compactdesc}
  \item[\T{Rechtsherzinsuffizienz}] Unzureichende Leistung des \E{rechten}
  Herzens,
  \item[\T{Linksherzinsuffizienz}]
  Unzureichende Leistung des \E{linken} Herzens, und
  \item[\T{Globalinsuffizienz}] Unzureichende Leistung beider Herzteile.
\end{compactdesc}

\subparagraph{Kompensation}
Abhängig davon, wie gut der Körper mit der
Einschränkung zurecht kommt unterteilt man:

\begin{compactdesc}
  \item[\T{Kompensierte Herzinsuffizienz}] Der Körper hat noch
  genug Reserven, kann gegensteuern und kommt mit der
  eingeschränkten Leistung aus.
  \item[\T{\E{De}kompensierte Herzinsuffizienz}] Der Körper kommt aufgrund
  des Fortschreitens der Erkrankung oder aufgrund eines erhöhten
  Bedarfs nicht mehr mit der Leistung aus. Es treten deutliche
  Symptome auf.
\end{compactdesc}
}{
\begin{itemize}
  \item Herzteil:
  \begin{itemize}
    \item Rechtsherzinsuffizienz,
    \item Linksherzinsuffizienz,
    \item Globalinsuffizienz.
  \end{itemize}
  \item Kompensation:
  \begin{itemize}
    \item Kompensierte Herzinsuffizienz,
    \item \E{De}kompensierte Herzinsuffizienz.
  \end{itemize}
\end{itemize}
}{}{}{}




\Layout{0}{Beurteilung des Schweregrades}
{%
Es ist wichtig den Schweregrad einer Herzinsuffizienz zu beurteilen. Sie
muss nicht dramatisch sein, für viele Patienten ist sie ein
chronischer Zustand mit dem sie gut leben können, da sich
der Körper mit der Zeit an die abnehmende Leistungsfähigkeit
gewöhnt.

Problematisch wird es jedoch, wenn der Patient plötzlich
eine Symptomatik entwickelt und die Kompensationsmechanismen (köpereigene, Medikamente, {\dots})
versagen. Man spricht dann von der \Es{dekompensierten}
Herzinsuffizienz. Man muss die Symptome generell \E{anhand ihres
Verlaufes beurteilen}: Bekannte, seit längerer Zeit bestehende
Symptome sind weniger wichtig als plötzlich auftretende
Symptome. Plötzlich auftretende Symptome sind oft \Es{Alarmzeichen}. 
}{
\begin{itemize}
  \item Chronisch: oft Gewöhnungseffekte, Kompensation.
  \item Kompensation versagt: Plötzlich neu auftretende
  Symptome, Beschwerden, Alarmzeichen.
\end{itemize}
}{}{}{}


\Layout{0}{Alarmzeichen}
{%
Weitere klassische Alarmzeichen sind Symptome wie
\Es{Atemnot}, \E{plötzlich auftretende
Ödeme} oder \E{akute Schmerzen}. 

Bei einem plötzlichen
Leistungsabfall des \E{linken} Herzens kommt es zu einem
\E{Rückstau} des Blutes \E{in den Lungenkreislauf}. Aufgrund der Stauung
kommt es zum Austritt von Flüssigkeit in das
Zwischengewebe und es entsteht ein
\Es{(kardiales) Lungenödem}. Ein typisches Zeichen dafür
ist das \Es{brodelnde Atemgeräusch}\index{Atemgeräusch!brodelndes - bei
Herzinsuffizienz}. Fällt die Leistung des \E{rechten}
Herzens, so kann man Beinödeme und/oder gestaute Halsvenen
beobachten. 
}
{
\begin{itemize}
  \item \Es{Atemnot}.
  \item  \E{Akute Schmerzen}.
  \item Stauung:
  \begin{itemize}
    \item \Es{Lungenödem} mit \E{"`brodelndem
    Atemgeräusch"'} (\E{li.} Herz).
    \item \Es{Beinödeme}, gestaute Halsvenen (\E{re.}
    Herz).
  \end{itemize}
\end{itemize}
}{}{}{}


Das kardiale Lungenödem wird im unter
\prettyref{topic:lungenoedem} 
besprochen.

\MassnahmenNG{Spezielle Maßnahmen}{Herzinsuffizienz}{}{
\begin{itemize}
  \item Lagerung: Oberkörper hoch.
  \item Vitale Gefährdung beurteilen \par \TxBeiVit
  \TxMassVitBes .
  \begin{itemize}
    \item \ce{O2} 10\LMin, bei Bedarf bis zu 15\LMin .
  \end{itemize}
  \item \ce{O2}: je nach Symptomen 6-8\LMin, bei
  Bedarf bis zu 15\LMin .
  \item Vitalwerte erheben \& Monitoring.
  \item Ursachenforschung.
\end{itemize}
}

\subsubsection{Linksherzinsuffizienz}
\index{Herzinsuffizienz!Links-}\index{Insuffizienz!Linksherz-}\index{Linksherzinsuffizienz}

Die Linksherzinsuffizienz entsteht aufgrund mangelnder Auswurfleistung
des linken Herzens. Durch den \E{Rückstau von Blut in den
Lungenkreislauf} kann es zu einem \Es{kardialen Lungenödem}
kommen.

\Layout{0}{Kardiales Lungenödem und Linksherzinsuffizienz}
{\index{Lungenödem!kardiales - bei Herzinsuffizienz} Patienten mit einem 
\E{kardialem Lungenödem} sind \E{lebensbedrohlich} krank. Das Herz verbraucht oft die allerletzten Reserven, schon die
Aufregung durch den Transport im Stiegenhaus zum Auto kann das Herz
endgültig in den Herz-Kreislaufstillstand bringen. 
}{
\begin{itemize}
  \item \Ca{Kardiales Lungenödem = lebensgefährlich
  krank!}.
  \item Herz verbraucht die allerletzten Reserven!
\end{itemize}
}{}{}{}

Es kommt immer wieder vor, dass solche Patienten im
Stiegenhaus oder im Auto reanimationspflichtig werden. 

\Cave{Kein Transport von Patienten mit Lungenödem ohne Notarzt. Auch wenn
das Spital "`um{\textquotesingle}s Eck"'
ist.}


\Layout{0}{\TxSymptome}{%
Der Patient klagt über Atemnot (\Es{Dyspnoe}). Oft
berichten die Patienten, dass sie nur mehr im Sitzen Luft
bekommen, sie können nicht flach liegen
(\Z{Orthopnoe}\ZFootnote{\E{Orthopnoe}: Der Patient muss aufrecht sitzen um
die Atemnot zu lindern.}).

Die Atemfrequenz und Herzfrequenz sind oft erhöht
(\Es{Tachypnoe},
\Es{Tachykardie}\index{Hyperventilation!bei
Herzinsuffizienz}), oft besteht ein \Es{Hustenreiz}.
Mitunter wirken die Patienten unruhig und ängstlich.

Grundsätzlich ist der Blutdruck vermindert
(\Es{Hypotonie}). Allerdings kann ein erhöhter Blutdruck
bei einer Blutdruckkrise (\Es{Hypertonie}, \E{hypertensive
Krise} (\prettyref{topic:hypertensive-krise}), \E{hypertensiver
Notfall} (\prettyref{topic:hypertensiver-notfall})) der eigentliche
\E{Auslöser} einer Herzinsuffizenz sein. Der Blutdruck
kann dann natürlich erhöht sein. 

Ein besonders kritisches Symptom ist das \Es{Lungenödem}.
Wenn es sehr stark ausgeprägt ist, kann man es mit freiem
Ohr hören (\E{"`brodelndes Atemgeräusch"'}). Diese
Patienten sind als besonders kritisch einzuschätzen.
}{
\begin{itemize}
  \item \Es{Dyspnoe}.
  \begin{itemize}
    \item \Z{Orthopnoe:} Patient bekommt nur mehr im Sitzen Luft, kann
    nicht flach schlafen, etc.
  \end{itemize}
  \item \Es{Tachypnoe}\index{Hyperventilation!bei Herzinsuffizienz}.
  \item \Es{Tachykardie}.
  \item \Es{Hustenreiz}.
  \item Unruhe
  \item \Es{Hypotonie} oder \Es{Hypertonie} (\Abk{RR}\,↑↑ kann Auslöser
  sein).
  \item \Es{Lungenstauung/Lungenödem} (Blut staut sich
  zurück in den Lungenkreislauf).
\end{itemize}
}{}{}{}






        




\subsubsection{Rechtsherzinsuffizienz}
\label{topic:rechtsherzinsuffizienz}

Die Rechtsherzinsuffizienz entsteht aufgrund mangelnder Auswurfleistung
des rechten Herzens.
\index{Rechtsherzinsuffizienz}\index{Herzinsuffizienz!Rechts-}
\index{Insuffizienz!Rechtsherz-}

\Layout{0}{\TxSymptome}
{%
Durch den Rückstau des Blutes in das venöse Gefäßsystem
kommt es oft zu Ödemen (Beine, \ldots) und
gestauten Halsvenen. 

Die Patienten müssen häufig nachts urinieren. Über
längere Zeit kann sich die Leber vergrößern und ein sog.
"`Wasserbauch"' (\Z{Aszites}) bilden\Z{, dabei wird durch den Rückstau
Wasser in den freien Bauchraum abgepresst.} 
}{
\begin{itemize}
  \item \Es{Bein-}/\Es{Hautödeme}.
  \item \Es{gestaute Halsvenen}.
  \item Nächtliches Wasserlassen.
  \item Lebervergrößerung, evtl. mit Aszites
  ("`Wasserbauch"').
\end{itemize}
}{}{}{}


% \begin{wrapfigure}{r}
% %	\begin{figure}
% 		%\Captionof{figure}{Aszites und Bein"odeme}	
% 		\includegraphics[width=5cm]{./images/aszites-beinoedeme.png}
% %	\end{figure}
% \end{wrapfigure}





% \clearpage
\subsection{Koronare Herzkrankheit und Akutes Koronarsyndrom}
\index{Koronare
Herzkrankheit}\index{Herzkrankheit!Koronare}\index{KHK}
\index{Akutes Koronarsyndrom}\index{Koronarsyndrom!Akutes}


\begin{description}
  \item[Die \T{Koronare Herzkrankheit}
  (\T{KHK})] ist eine Erkrankung der
  \Es{Herzkranzgefäße} (Koronargefäße)
  bei der es zu einer Verengung der Gefäße kommt.
  \item[Das \T{Akute Koronarsyndrom}]\index{Akutes Koronarsyndrom} ist ein
  Sammelbegriff für \E{Herzinfarkt}, \E{stabile} und \E{instabile Angina pectoris}.
  Bei diesen Krankheitsbildern kommt es zu einer akuten
  \Es{Unterversorgung des Herzmuskels} mit Blut und
  Sauerstoff. Die Ursache sind in erster Linie verengte oder verstopfte
  Herzkranzgefäße.
\end{description}

\subsubsection[Angina pectoris]{Angina pectoris}
\index{Angina pectoris}

\Definition{Plötzliches Engegefühl in der Brust meist mit
Thoraxschmerz infolge einer \underline{vorübergehenden}
Unter\-ver\-sorgung des Herzmuskels mit sauerstoffreichen Blut.}

\Layout{0}{\TxBeschreibung}
{%
Es sind \E{zwei Auslösemechanismen} möglich: 
\begin{inparaitem}
  \item Einerseits kann eine Minderdurchblutung des Herzmuskels (Myokard)
  aufgrund einer oder mehrerer \E{Einengung(en) von Herzkranzgefäßen}
  (Koronararterien) die Ursache sein, oder 
  \item es kommt zu einem \E{Erhöhtem Sauerstoffbedarf} des Herzens bei
  Anstrengung.
\end{inparaitem}

In beiden Fällen kommt es zu einem
\Es{Missverhältnis zwischen Sauerstoffangebot und  
Sauerstoffbedarf} des Herzens, die Folge sind Schmerzen, Atemnot und ein
Listungsverlust!

Der akute Angina-Pectoris-Anfall gehört zu der Familie des Akuten
Koronarsyndroms. Dazu gehört auch der Herzinfarkt, dem der gleiche
Mechanismus wie der der Angina Pectoris zu Grunde liegt, wobei es beim Infarkt
im Gegensatz zur Angina pectoris zu einem \E{totalen}
Herzkranzgefäß\E{verschluss} mit Absterben von Herzmuskelgewebe kommt.
\A{KOCH: nein; GABS: Doch. }

Bei den meisten Patienten ist die Angina Pectoris schon bekannt, bzw. es
besteht eine "`\index{Koronare
Herz-Krankheit}\T{Koronare Herz-Krankheit}
(\index{KHK}\T{KHK})"'. Die KHK ist
eine Erkrankung der Herzkranzgefäße, diese werden durch
Ablagerung von Fetten und Kalk geschädigt und eingeengt, dadurch kann
weniger Blut durchfließen. So kommt es zu einer Mangelversorgung
des Herzmuskels und zu den Symptomen der Angina pectoris.  
}{
Missverhältnis zwischen Sauerstoffangebot und  
Sauerstoffbedarf.
\begin{itemize}
  \item Einengung(en) von Herzkranzgefäßen. 
  \item Erhöhter Sauerstoffbedarf bei Anstrengung. 
\end{itemize}

Koronare Herz-Krankheit (KHK).
}{}{}{}





\Layout{20}{\TxSymptome}
{%
Das Leitsymptom ist der \Es{akute, nicht atemabhängige,
Brustschmerz}. Der Schmerz ist meist hinter dem Brustbein lokalisiert und
\Es{strahlt häufig aus}, häufig in den Kiefer, (linken) Arm und in den Rücken,
evtl. auch in den Oberbauch. Der Schmerz wird als brennend bzw. stechend beschrieben.
Dazu kommt i.\,d.\,R. ein \Es{Druckgefühl} am Brustkorb, \emph{"`als würde jemand
auf dem Brustkorb stehen."'} \Es{Todesangst und Vernichtungsgefühl} ist
häufig. 

Ein weiteres wichtiges Leitsymptom ist die \Es{Atemnot}
(\E{Dyspnoe}).

Ein akuter Angina pectoris-Anfall ist primär \Es{nicht sicher von einem
Herzinfarkt unterscheidbar}. Die Symptome treten eher bei Belastung auf und
vergehen oft bei Entspannung, das Vorübergehen der Beschwerden darf aber
natürlich nicht abgewartet werden! 

}{
\begin{itemize}
  \item Brustschmerz, brenndend/stechend, Druckgefühl.
  \item Oft Ausstrahlung häufig in den Kiefer,
  linken Arm, Rücken, evtl. Oberbauch.
  \item Todesangst, Vernichtungsgefühl.
  \item Dyspnoe.
  \item Primär \Es{nicht sicher von MCI unterscheidbar}.
  \item Symptome eher bei Belastung, vergehen oft bei Entspannung.
\end{itemize}

}
{./images/hirtler-aass/Herzschmerz.pdf}
{Schmerzausdehnung beim ischämischen
Herzschmerz.}
{Hirtler.}




% \Married
% {
%  \includegraphics[width=\textwidth]{./images/noauth/AASdev20090228-img72.png}
% \Captionof{figure}{\AuthNotes{Bild fehlt!} Patient mit
% einem akuten Koronarsyndrom, vermutlich Angina Pectoris. Bei kaltem Wetter, beim
% Stiegen steigen mit schwerem Gepäck: Ein plötzliches Engegefühl
% im Brustkorb und ein stechender Schmerz im linken Thorax.}
% }{
% \includegraphics[width=\textwidth]{./images/noauth/AASdev20090228-img73.png}
% \Captionof{figure}{\AuthNotes{Bild fehlt!}
% Schmerzausdehnung beim ischämischen Herzschmerz.}
% }

\MassnahmenNG{Spezielle Maßnahmen}{Angina pectoris}{}{
\begin{itemize}
  \item \TxMassVitBes .
  \begin{itemize}
    \item Lagerung: Oberkörper hoch.
    \item \ce{O2} 10\LMin, bei Bedarf bis zu 15\LMin .
  \end{itemize}
  \item \Es{striktes Bewegungsverbot}!
  \item Patient darf eigenes Nitro nur  bei RR
  {\textgreater}100 syst. nehmen!
\end{itemize}
\AuthNotes{\#13: ?? eigenes Nitro nehmen?

Vorschlag: Nur wenn sein Nitro leer war und er es deswegen nicht
genommen hat bei RRsyst{\textgreater}110

Wenn er es schon genommen hat und es nicht geholfen hat
dann nicht <-- KOCH richtig wichtig
GABS: was soll ``richtig wichtig bedeuten? ob es der Pat. richtig genommen hat?

KOCH: Pat. darf eigenes Nitro nur  bei RR {\textgreater}100 syst.
\underline{selbst ein}nehmen!}
}

\subsubsection{Herzinfarkt - Myokardinfarkt (MCI)}
\index{Herzinfarkt}\index{Infarkt!Herz-}\index{Myokardinfarkt}\index{MCI}
\index{Herzkranzgefäße!Verschluss}\index{Koronargefäße!Verschluss}

% \DefinitionExp{Herzinfarkt}

\Definition{Verschluss eines oder mehrerer Herzkranzgefäße mit
  anschließendem Absterben des betroffenen Herzmuskelgewebes.}

\Layout{20}{\TxSymptome}
{%
Die Symptome entsprechen im Wesentlichen denen der Angina pectoris, doch
vergehen diese nicht mehr und es bleiben -- selbst bei optimaler Behandlung --
Schäden am Herzmuskel bestehen. Bei Diabetikern kann es aufgrund der
Nervenschädigungen auch zu einem schmerzlosen Infarkt kommen, man
spricht dabei auch vom \T{"`stummen
Infarkt"'}\index{Stummer Infarkt}.

Sonst steht normalerweise der \E{stechende/brennende}, nicht atemabhängige
Brustschmerz zusammen mit \Es{Atemnot} und \Es{Todesangst} im Vordergrund. Der
\Es{Brustschmerz} kann in Richtung Magen, rechter oder linker Hand, Kiefer
oder auch Rücken \Es{ausstrahlen}. Zusätzlich zu dem
Schmerz verspürt der Patient auch oft ein \Es{Engegefühl} im
Brustkorb, \emph{"`als würde jemand darauf stehen"'}. Die
Atemnot kann massiv ausfallen. 

Weiters können aufgrund des eintretenden \Es{kardiogenen
Schocks} die entsprechenden Schockzeichen beobachtet
werden: Blässe, Kaltschweißigkeit, Hypotonie, etc.
(\ref{topic:schock-kardiogener}).
}{
\begin{itemize}
  \item Symptome wie bei Angina Pectoris, vergehen aber nicht
  mehr
  \begin{itemize}
    \item Stechender/brennender \Es{Brustschmerz} mit Ausstrahlung.
    \item "`Stummer Infarkt"' bei Diabetikern.
    \item Druckgefühl, "`als würde jemand am
    Brustkorb stehen"'.
    \item Todesangst.
    \item \Es{Dyspnoe.}
    \item Kardiogener Schock.
    %	\item Keine ausreichende Besserung auf Nitro
  \end{itemize}
\end{itemize}
}
{./images/hirtler-aass/Herzschmerz.pdf}
{Schmerzausdehnung beim ischämischen
Herzschmerz.}
{Hirtler.}
  
\Layout{0}{\TxKomplikationen}
{%
\begin{itemize}
  \item Herzrhythmusstörungen (alle Arten), bis hin
  zum Atem-/Kreislaufstillstand und Reanimation.
  Besonders häufig kommt es als Frühkomplikation zum
  plötzlichen Kammerflimmern.
  \item Kardiogener Schock \index{Kardiogener
  Schock!beim Herzinfarkt}\index{Schock!Kardiogener!beim
  Herzinfarkt}
\end{itemize}

}{
\begin{itemize}
  \item \Es{Lebensgefahr!}
  \item Jede Herzrhythmusstörung als Komplikation möglich.
  \item Kammerflimmern als Frühkomplikation häufig!
\end{itemize}

}{}{}{}

\begin{table}[H]
\Dias{Herzinfarkt.}
{%
\includegraphics[width=\linewidth]{./images/wikipedia-cc-by/Heart_coronary_artery_lesion-lq.jpg}
\Captionof{figure}{Die Koronargefäße versorgen den
Herzmuskel. \SourcePicture{Patrick J. Lynch, CC-BY-2.5}} }{%
\includegraphics[width=\linewidth]{./images/wikipedia-cc-by/Heart_ant_wall_infarction-00800.jpg}
\Captionof{figure}{Abgestorbenes Muskelgewebe. \SourcePicture{Patrick J. Lynch,
CC-BY-2.5}} }{%
\includegraphics[width=\linewidth]{./images/wikipedia-cc-by/Hk_coronary_bionerd.jpg}
\Captionof{figure}{Herzkatheter: Darstellung der
Herzkranzgefäße. \SourcePicture{WmCo "`Bionerd"', CC-BY-3.0}} }
\end{table}

% \begin{center}
%  \TempPicMissing{Infarkt [Warning: Image ignored]
%  {AASdev20090228-img74.png}}
% \end{center}

\MassnahmenNG{Spezielle Maßnahmen}{Herzinfarkt}{}{   
\begin{itemize}
  \item \TxMassVitBes .
  \begin{itemize}
    \item Lagerung: Oberkörper hoch.
    \item \ce{O2} 10\LMin, bei Bedarf bis zu 15\LMin .
  \end{itemize}
  \item Beengende Kleidung öffnen.
  \item Striktes Bewegungsverbot!
\end{itemize}
    
    %%%%%%%%%%%%%%%%%%%%%%%%%%%%%%%%%%%%%%%%%%%%%%%%%%%%%%%%%%%
% TODO : ERRATUM: Nitroeinnahme MCI Fehlerhafter Text
\ERROR{GABS}{yyyy-mm-dd}
{ERRATUM: Fehlerhafter Text: Nitroeinnahme MCI}
{Nitroeinnahme MCI: Darauf wird nicht eingegangen. Dieser Teil ist fehlerhaft
oder sehr unvollständig. Die Aussagen sind unklar oder verwirrend. Eine
Überarbeitung ist dringend notwendig!}
% %%%%%%%%%%%%%%%%%%%%%%%%%%%%%%%%%%%%%%%%%%%%%%%%%%%%%%%%%%
}


\Layout{0}{Weitere Therapie}
{%
Herzinfarktpatienten werden "`besonders"'
behandelt und normalerweise nicht auf eine beliebige Interne Station
gebracht\cite{Kalla2006}. Der Notarzt muss i.\,d.\,R. zwischen dem Einleiten
einer Lyse-Therapie und der Einweisung in ein Herzkatheterlabor entscheiden.

\index{PTCA}
\index{Herzkatheterlabor}
Bei der \Es{Lyse-Therapie} versucht man mittels eines speziellen
Medikaments eine Verstopfung aufzulösen. Mit einem \Es{Herzkatheters}
versucht man die verengten oder verstopften Gefäße aufzudehnen. Dazu benötigt
man ein entsprechend ausgestattetes
\Es{Herzkatheterlabor}. Der Herzkatheter wird oft
mit "`\E{PCI}"' \Z{(Percutane Coronare Intervention)} oder "`\E{PTCA}"'
\Z{(Perkutane Transluminale Coronare Angioplastie)} abgekürzt.
}
{
\begin{itemize}
  \item Lyse.
  \item Herzkatheterlabor (PCI, PTCA).
  \opt{REGVIE}{
\begin{itemize}
  \item Dienstbereites Katheterlabor: eigener Dienstplan.
\end{itemize}
}  	
\end{itemize}

}{}{}{}

\opt{REGVIE}{
In Wien gibt es einen \E{Dienstplan} für die Katheterlabore, d.h. 
jedes Katheterlabor hat an bestimmten Tagen der Woche rund um die 
Uhr Bereitschaftsdienst. Eine Voranmeldung und vorangehende Absprache 
zwischen Notarzt und diensthabenden Oberarzt der Abteilung
ist erforderlich. Die PTCA ist in Wien u.\,a. im AKH, Donauspital und KH
Hietzing verfügbar. }

\Cave{{\bfseries Kein Transport ohne Notarzt!}

Auch nicht wenn
das Spital "`um{\textquotesingle}s Eck"'
ist.

Der Patient braucht vermutlich {\bfseries kein beliebiges
Spital} sondern ein {\bfseries dienst\-bereites
Herzkatheterlabor}. Der Notarzt muss entscheiden, ob das
notwendig und sinnvoll ist.

Auch wenn der Patient ein Einweisungsformular (Spitalszettel) von einem
Arzt hat (Ärztefunkdienst, Hausarzt) oder ein Hausarzt
meint "`es sei eh nicht dringend"' unbedingt
einen Notarzt beiziehen.
}



\subsection{Herzrhythmusstörungen}

\label{topic:herzrhythmusstoerungen}
\index{Herzrhythmusstörungen}
\index{Rhythmusstörungen!Herz-}

Herzrhytmusstörungen sind häufig, allerdings auch oft sehr unterschiedlich in
ihrer Bedeutung: Manche sind chronisch und eher wenig gefährlich, andere können
plötzlich auftreten und sogar akut lebensbedrohlich sein. Rhythmusstörungen
können sowohl selber das Problem darstellen, als auch \E{als Komplikationen} im
Rahmen anderer Krankheitsbilder (z.\,B. Herzinfarkt, etc.) auftreten.

\Layout{0}{Arten von Herzrhythmusstörungen}
{
Herzrhythmusstörungen können hervorgerufen werden durch eine \Es{Störung der
Frequenz} (zu langsam, zu schnell) oder einer \Es{Störung der Regelmäßigkeit}.
Wenn man die elektrischen Ströme mittels eines EKG darstellt, kann man
zusätzlich auch \Es{Störungen der Form der elektrischen Herzströme} erkennen
und beschreiben. Ein halbautomatischer Defibrillator (AED) erledigt dies
automatisch, er zeichnet die elektrischen Herzströme auf und beurteilt diese. 

Im Rahmen von Herzrhythmusstörungen kann es zu einer \Es{Entkoppelung der
elektrischen und der mechanischen Herzaktion} kommen, z.\,B. wenn nicht jedem
elektrischen Impuls ein Herzschlag folgt. 
}{
\begin{itemize}
  \item Störung der
  \begin{itemize}
    \item Frequenz.
    \item Regelmäßigkeit.
    \item Form der elektrischen Herzströme.
  \end{itemize}
\end{itemize}

Entkoppelung von elektrischer und der mechanischer Herzaktion möglich!
}{}{}{}

\Layout{0}{Hier zählt die Geschwindigkeit: Tachy- und Bradykardie}
{%
\index{Bradykardie}
\index{Tachykardie}
%
\begin{labeling}[~--]{Tachykardie}
  \item[\T{Bradykardie}]  zu langsamer Herzschlag
  ({HF \textless} 60\,\nicefrac{}{min} {\tiny (beim Erwachsenen)})
  
  Leistungssportler können aufgrund ihres trainierten Körpers auch sehr
  niedrige Ruhepulsfrequenzen haben.
  \item[\T{Tachykardie}]  zu schneller
  \marginpar{Tachometer im Auto:
  Geschwindigkeitsanzeiger, meistens zu schnell ;)}
  Herzschlag
  (HF {\textgreater} 100\,\nicefrac{}{min} {\tiny (beim
  Erwachsenen)})
\end{labeling}

Denke bei der Beurteilung auch immer an die Umstände und den
Erregungszustand des Patienten!
}{
\begin{itemize}
  \item Bradykardie \dotfill\ zu langsam. 
  \item Tachykardie \dotfill\ zu schnell. 
\end{itemize}
}{}{}{}

\prettyref{tab:herzrhythmusstoerungen} gibt einen Überblick über die häufigsten
Herzrhythmusstörungen.

\begin{table}[H]
\begin{gWideParEnv}
\Captionof{table}{Übersicht: Wichtige und häufige
Herzrhythmusstörungen.\label{tab:herzrhythmusstoerungen}}

% Index-Einträge für Tabelle ---
\index{Bradykardie} 
\index{Tachykardie}
\index{Kammerflimmern}
\index{Flimmern!Kammer-}
\index{Vorhofflimmern}
\index{Flimmern!Vorhof-}
\index{Asystolie}
\index{Pulslose Ventrikuläre Tachykardie} 
\index{Tachylardie!Pulslose Ventrikuläre}
% --- Ende


\begin{minipage}{\linewidth}
\begin{tabulary}{\textwidth}{LLL}
\TabHeading{Rhythmus}
& 
\TabHeading{Bedeutung}
& 
\TabHeading{Bemerkung}
\\\midrule\addlinespace
\TabSub{Bradykardie}

{\scriptsize (HF\,{\textless}\,60/min)}\MarginBugfix{(beim
Erwachsenen)} 
& 
Situations\-abhängig 
& 
\emph{"`zu langsam"'}
\\\addlinespace
\TabSub{Tachykardie} 

{\scriptsize (HF\,{\textgreater}\,100\MarginBugfix{(beim
Erwachsenen)}/min) }
& 
Situations\-abhängig 
& 
\emph{"`zu schnell"'}
\\\addlinespace\midrule \addlinespace
\TabSub{Kammerflimmern}
& Reanimation  & \emph{"`schockbarer"'} Rhythmus 
\\\addlinespace
\TabSub{Pulslose Ventrikuläre Tachy\-kardie} 
& Reanimation 
& \emph{"`schockbarer"'} Rhythmus 
\\\addlinespace
\TabSub{Asystolie} 

{\scriptsize (HF\,=\,0)} 
& 
Reanimation 
& 
Nulllinie, \emph{"`nicht-schockbarer"'} Rhythmus
\\\addlinespace
\TabSub{Vorhofflimmern}
& regellose Erregung im Vorhof

oft symptomlos 
& 
\E{Chronische Erkrankung.} Sehr viele Leute haben
Vorhofflimmern und leben ganz gut damit. 

Risikofaktor für:

\E{Tachykardie-Anfälle}

\E{Thrombenbildung}, dadurch erhöhtes Lungenembolie-
und Schlaganfallrisiko. Deshalb Vorbeugung mit
gerinnungs\-hemmenden ("`blut\-ver\-dünnenden"') Medi\-ka\-menten:
z.\,B. \E{"`Marcoumar"'} ${\rightarrow}$ Patient ist blutungsgefährdet
\\\bottomrule
\end{tabulary}
\end{minipage}
\end{gWideParEnv}
\end{table}







\subsubsection{Tachykarde Attacke}

\Definition{Plötzliche Tachykardie ohne erkennbare Ursache, wobei der Patient
(meistens) Symptome wahrnimmt und eine HF\,>\, 160\,\,\nicefrac{}{min}
vorliegt.}


\Layout{0}{\TxBeschreibung}
{%
Tachykarde Attacken sind ein häufiges Notfallbild. Da oft chronische
Erkrankungen die Ursache sind, haben die Patienten diese
Attacken immer wieder und kennen die Symptome bereits (und oft auch die Behandlung). Die
Ursachen können sehr unterschiedlich sein und sind auch vom Fachmann oft nicht
leicht zu ermitteln. 

Eine große Rolle spielt die \Es{Anamnese}: Oft sind entsprechende
Grunderkrankungen schon seit langem bekannt, oft gibt es bereits entsprechende Arztbriefe oder
Befunde. Ein rasches Hinweisen auf diese Befunde ist für den nachbehandelnden
Arzt wichtig um eine rasche, zielführende Behandlung einzuleiten.
}{
\begin{itemize}
  \item Oft Grunderkrankung bekannt.
  \item Anamnese!
  \item Befunde, Arztbriefe!
\end{itemize}
}{}{}{}

\Layout{0}{\TxSymptome}
{Es liegt klarerweise eine \Es{Tachykardie} vor. Der Puls
ist i.\,d.\,R. hoch, meist über 160\,\nicefrac{}{min}. Der Patient klagt oft über ein
\Es{"`Herzklopfen"'} und ist unruhig. Manchmal wird der Anfall von
Angstgefühlen begleitet.

Oft kommen \Es{Begleitsymtome} dazu, wie zum Beispiel Atemnot oder ein leichter
Brustschmerz. Dies sind erste Zeichen, dass das Herz an
seine Grenzen der Belastbarkeit stößt.
}{
\begin{itemize}
  \item Tachykardie.
  \item Spürbares Herzklopfen.
  \item evtl. Atemnot.
  \item evtl. Leichter Brustschmerz.
  \item evtl. Angstgefühl.
\end{itemize}
}{}{}{}

\Layout{0}{{\Lightning}Gefahr }
{%
Da das \Es{Herz plötzlich viel mehr Arbeit leisten
muss}, kann es sein, dass es an seine \Es{Grenzen}
stößt und \Es{versagt} bzw. die Blutversorgung
des Herzens für diese Arbeit nicht mehr reicht und ein
Angina-pectoris-Anfall oder sogar ein Herzinfarkt\A{ KOCH: passt mit der
Definition MCI nicht überein, AP wäre besser oder? GABS: AP
dazugenommen, Infarkt aber gut möglich. } entsteht. 
}{
\begin{itemize}
  \item Herz überfordert:
  \begin{itemize}
    \item Angina pectoris.
    \item Herzinfarkt.
  \end{itemize}
\end{itemize}
}{}{}{}





\MassnahmenNG{Spezielle Maßnahmen}{Tachykarde Attacke}{}{
\begin{itemize}
  \item \TxBeiVit \TxMassVitBes .
  \begin{itemize}
    \item Lagerung: Oberkörper hoch.
    \item \ce{O2} 10\LMin, bei Bedarf bis zu
    15\LMin .
    \item Notarzt: Eine medikamentöse Therapie ist meistens
    erforderlich oder ratsam.
  \end{itemize}
  \item Beengende Kleidung öffnen.
  \item Striktes Bewegungsverbot!
  \item Psychische Betreuung! Oase der Ruhe schaffen.
  \item Wenn der Patient versorgt ist (d.h. alle obigen Punkte erfüllt
  wurden) soll -- wenn möglich -- ein \E{EKG}
  (Extremitätenableitung, \prettyref{topic:ekg}) abgeleitet werden. Der
  Ausdruck dient v.a. der \E{Dokumentation} und der Information für das
  nachbehandelnde Personal (Notarzt, Spital, Facharzt, \ldots).
\end{itemize}
}



\subsubsection{Besondere Rhythmen}


\Layout{0}{Kammerflimmern}
{\Ca{Reanimationspflichtig}. 
\index{Kammerflimmern}
%
NICHT verwechseln mit
Vorhofflimmern! Die elektrische Erregung im Herzen ist ungerichtet, als würden alle Autos 
eines Parkplatzes gleichzeitig in alle Richtungen losfahren 
${\rightarrow}$ der Herzmuskel kontrahiert nicht
geordnet, er "`zittert"' nur (wie ein Pudding), es wird
\Es{keine Auswurfleistung} erreicht!
}
{
\begin{itemize}
  \item \Ca{Reanimation!}
  \item Wirre Erregung, keine sinnvolle Muskelarbeit ${\rightarrow}$ kein
  Auswurf.
\end{itemize}

}{}{}{}

    
\Layout{0}{Pulslose Ventrikuläre Tachykardie}
{\Ca{Reanimationspflichtig!} 
\index{Pulslose Ventrikuläre Tachykardie}
%
Vorstufe zum Kammerflimmern. Das Herz kontrahiert
so schnell, dass es sich nicht mehr füllen kann
${\rightarrow}$ \Es{keine Auswurfleistung}.
}
{
\begin{itemize}
  \item  \Ca{Reanimation!}
  \item Keine Zeit zum Füllen ${\rightarrow}$ kein Auswurf.
\end{itemize}

}




\Layout{0}{Asystolie}
{\Ca{Reanimationspflichtig}. 
\index{Asystolie}



Im Herz entsteht keine elektrische
Erregung, der Herzmuskel kontrahiert daher nicht und es erfolgt
{keine Auswurfleistung} 
}{
\begin{itemize}
  \item \Ca{Reanimation!}
  \item Keine elektrische Erregung, keine Muskelarbeit, kein Auswurf.
\end{itemize}
}{}{}{}
    
\Layout{0}{Vorhofflimmern}
{
\index{Vorhofflimmern}

Regellose Erregung im \E{Vorhof}, die oft symptomlos ist. 
Sehr viele Leute haben Vorhofflimmern und leben ganz gut damit.

Komplikationen:
\begin{itemize}
  \item Tachykardie-Anfälle
  \item Thrombenbildung
  \begin{itemize}
    \item deshalb Therapie mit \Es{gerinnungshemmenden}
    ("`blutverdünnenden"') Medikamenten:
    Marcoumar ${\rightarrow}$ Patient ist Blutungsgefährdet
  \end{itemize}
\end{itemize}
}
{
% TODO VIT: Slide fehlt!
\begin{itemize}
  \item Flimmern im \Es{Vorhof}, Kammern arbeiten normal.
  \item Häufige und oft chronische Erkrankung.
  \item Thrombenbildung, Tachykardieanfälle.
  \item Oft Therapie mit \E{gerinnngshemmenden Medikamenten}.
\end{itemize}

}{}{}{}


\subsection{Hypotonie}

\Definition{Unangemessen zu niedriger Blutdruck.}
 
\noindent \Es{Richtwert:} Beim Erwachsenen RR {\textless} 100\,mmHg
syst. Jeder Mensch hat allerdings im gewissen Rahmen einen
"`eigenen Normalwert"', bei Kindern gelten grundsätzlich
andere Werte.

% einvernähmlich gestrichen KOCH 2008-05-18
% 
% \paragraph{Ursachen}
% 
% \begin{itemize}
%     \item Cardial
%     \item Volumenmangel, Exsikkose
%     \item Regulationsstörung (vasovagal)
%     \item Sportler, junge Frauen
%     \item Schock
% \end{itemize}
% 
% \paragraph{\TxSymptome}
% 
% \begin{itemize}
%     \item Müdigkeit bis Kollaps
%     \item RR-Abfall mit Bewußtseinseintrübung
%     \item "`schwarzwerden vor Augen"', Ohnmacht
%     \item Sturz = "`Therapie"'
% \end{itemize}


\subsubsection{Kollaps, Synkope}

\Definition{passagere Kreislaufregulationsstörung  mit RR-Abfall,
    Blutverteilungsstörung  ohne Volumsverlust}

\Layout{0}{\TxSymptome}
{ 
Der Patient berichtet berichtet über Schwindel, evtl. ist er zusammengesackt,
war kurz bewusstlos. Der Patient wacht oft auf dem Boden liegend auf und weiß
nicht, was geschehen ist. Er erzählt  es wäre ihm schwarz vor den Augen
geworden. Der Blutdruck ist meist niedrig, der Patient ist eventuell bradykard oder
tachykard. 
}{
\begin{itemize}
  \item Schwindel. 
  \item evtl. kurze Bewusstlosigkeit, Ohnmacht, Schwarzwerden vor Augen. 
  \item RR ${\downarrow}$, evtl. Bradykardie oder Tachykardie. 
\end{itemize}
}{}{}{}




\Layout{2}{Ursachen}
{
Die Ursachen sind präklinisch oft
nicht sicher ermittelbar, mitunter bleibt die wahre Ursache
für immer ein Rätsel. Daher ist eine Hospitalisierung zur
bestmöglichen Abklärung angeraten. 

Ein Kollpas kann seine Ursachen in verschiedenen Organssystemen haben:

\begin{itemize}
  \item Herz (z.\,B. Herzrhythmusstörungen),
  \item Gefäßen (z.\,B. Hitzebedingte Erweiterung
  der Gefäße mit "`Versacken"' des Blutes in den Beinen),
  \item Blut (z.\,B. Blutarmut (\T{Anämie}),
  Regelblutung),
  \item Hirn,
  \item Hirngefäßen (z.\,B. TIA,
  \prettyref{topic:tia} ) \ldots
\end{itemize}

Es gibt jedoch auch typische "`Kollaps-Situationen"': Ein
schwüler Tag in einer überfüllten U-Bahn-Garnitur, nicht
gefrühstückt, auf dem Weg zu einem Vorstellungstermin und
seit gestern in der Menstruation. Oder im stickigen
Veranstaltungszelt bei 40\textdegree , wobei der der ganze
Flüssigkeitskonsum aus 5 (entwässernden) Dosen
RedBull\texttrademark
bestand. So mancher war auch ob eines Begräbnisses so
ergriffen, dass er sich fast in die Grube neben den Sarg
dazugelegt hätte \ldots 

}{

}{}{}{}


\MassnahmenNG{Spezielle Maßnahmen}{Kollaps/Synkope}{}{

\begin{itemize}
  \item Erheben ob eine akute vitale Bedrohung besteht
  \item Ausschluss von Notfallsdiagnosen soweit
  möglich, die gesamten zur Verfügung stehenden
  Möglichkeiten sind auszuschöpfen.
  \item Im Besonderen:
  \begin{itemize}
    \item Traumacheck (Sturz?),
    \item Neurocheck inkl. Blutzuckermessung,
    \item Temperatur,
    \item Ausführliche Anamnese \ldots
  \end{itemize}
  \item Lagerung: Nach Ausschluss von Herzbeschwerden,
  Atemnot und relevanten Verletzungen: Beine hoch.
  \item Hospitalisierung zur Abklärung
  anstreben\newline
  Abteilung: Innere Medizin.
\end{itemize}

}

\subsection{Hypertonie}

\index{Hypertonie}\T{Hypertonie}
(\index{Bluthochdruck}\T{Bluthochdruck}): Oft
symptomlose \textit{chronische} Erkrankung welche auf Dauer schwere
Schäden am Körper anrichten kann und im Zuge von plötzlichen
Blutdruckkrisen auch akut Probleme machen kann. 

\begin{itemize}
  \item Chronische Hypertonie: Obergrenze derzeit bei \RRmmHg{140}{90} (beim
  Erwachsenen).
\end{itemize}

\Cave{Bei \Es{Schwangeren} sind die Grenzwerte genau zu beachten!
(\Es{EPH-Gestose}, \prettyref{topic:eph-gestose})!}

\Layout{0}{Ursachen}
{
\begin{itemize}
  \item Widerstandserhöhung (durch Engstellung der Gefäße).
  \item Vermehrte Herzarbeit.
  \item Sympathikusaktivierung, Adrenalin+Cortison (Stresshormone) durch
  Dauerstress,  Bewegungsmangel, genetische Veranlagung.
  \item u.\,a. \ldots
\end{itemize}

}{

}{}{}{}


\Layout{60}{Langzeitfolgen}{%
Die Hypertonie hat in der Notfallmedizin auch eine wichtige
indirekte Bedeutung: Sie ist Ursache für viele andere
Erkrankungen, die ihrerseits akut lebensbedrohlich werden
können. 

Zwei Wirkungen der Hypertonie sind dabei besonders wichtig:
Die \Es{Gefäßschädigung} und die \Es{Belastung des Herzens durch
den hohen Gegendruck}. 

Bei der Gefäßschädigung kommt es zur Ablagerung durch Kalk
und sonstige Plaques und über längere Zeit zur \E{Verengung
oder Verstopfung der Gefäße}. Je nach Organ kann das sehr
massive Auswirkungen haben. Im Herz führt es zur \E{Angina
Pectoris} oder zum \E{Herzinfarkt}, im Hirn zu \E{Schlaganfällen}, in
den Nieren zu einer \E{Niereninsuffizienz}, usw. Selbst große
Gefäße wie die Hauptschlagader (Aorta) können derartig
geschädigt werden dass sie reißen können (Aortenaneurysma) --- eine
für den Patienten sehr lebensgefährliche Komplikation. 

Die Schädigung des Herzmuskels ist die Folge von dem
erhöhten Druck gegen den das Herz pumpen muss. Der Muskel
vergrößert sich um mit der Belastung fertig zu werden, kann
aber dann nicht mehr ausreichend mit sauerstoffreichen
Blut versorgt werden. 

Wir haben also gesehen, dass die Hypertonie ein wichtiger
\E{Risikofaktor} für viele schwere Notfälle ist. Daher ist es
auch verständlich dass die Behandlung der Hypertonie einen
großen Stellenwert in der medizinischen Versorgung der
Bevölkerung hat. 
}{
\begin{itemize}
  \item \Es{Gefäßschädigung}
  \begin{itemize}
    \item Herzkranzgefäße ${\rightarrow}$ Infarkt,
    \item Hirngefäße ${\rightarrow}$ Schlaganfall,
    \item Nierengefäße ${\rightarrow}$ Niereninsuffizienz ${\rightarrow}$
    Dialyse ${\rightarrow}$ typischer KTW-Patient,
    \item Hauptschlagader ${\rightarrow}$ Aneurysma.
  \end{itemize}
  \item \Es{Herzmuskelschädigung} \par  (Herz muss gegen
  größeren Widerstand pumpen).
  \begin{itemize}
    \item Krankhafte Vergrößerung des Herzmuskels.
  \end{itemize}
\end{itemize}
}
{./images/teaser/imgp0481.jpg}
{}
{GABS.}




\subsubsection{Hypertensive Krise und Hypertensiver Notfall}
\label{topic:hypertensive-krise}
\label{topic:hypertensiver-notfall}

\Layout{2}{\TxBeschreibung}
{
Eine plötzliche, sehr starke Erhöhung des Blutdrucks wird als hypertensive
Krise oder hypertensiver Notfall bezeichnet. Der Unterschied zwischen den
beiden Diagnosen ist das Fehlen bzw. Vorhandensein von typischen Symptomen, s.
\prettyref{tab:hxpertension-krise-notfall}.

}{}{}{}{}

\Layout{0}{\TxSymptome}
{
S. \prettyref{tab:hxpertension-krise-notfall}. Kopfschmerzen und Nasenbluten
(\T{Epistaxis}\index{Epistaxis}) sind häufige Leitsymptome für einen
hypertensiven Notfall. 

}{
S. \prettyref{tab:hxpertension-krise-notfall}.
}{}{}{}

\begin{table}[H]
\Captionof{table}{Unterscheidung Hypertensive Krise und Hypertensiver
Notfall.\label{tab:hxpertension-krise-notfall}}

\begin{tabularx}{\textwidth}{XX}
\TabHeading{Hypertensive Krise}
&
\TabHeading{Hypertensiver Notfall}
\\\addlinespace\midrule\addlinespace
\Abk{RR}\,{\textgreater}\,230/130\,mmHg\footnote{Richtwert}

\Es{Patient fühlt sich gut}
 &
\Abk{RR}\,{\textgreater}\,230/130\,mmHg\footnote{Richtwert}

\Es{+ Symptome}
\\\addlinespace
Zufallsbefund?

Oder steht der Blutdruck doch in Zusammenhang mit der
Berufung?

\textit{Patient hat überhaupt keine Beschwerden?} Warum hat
er Sie dann überhaupt gerufen?

Wenn es wirklich nur ein Zufallsbefund ist ("`Gesundheitstage"',
"`Vorsorgeuntersuchung"'):


&
\subparagraph{\Ca{Symptome:}}

\begin{itemize}
    \item Kopfschmerz.
    \item Sehstörungen.
    \item Brustschmerz.
    \item Evtl. Nasenbluten.
    \item Übelkeit, Erbrechen, Schwindel.
\end{itemize}

\Cave{Gefahr des Herzversagens -- Das Herz muss
plötzlich gegen großen Druck pumpen}

\\\addlinespace
Maßnahmen: Eine notfallmäßige Blutdrucksenkung ist (eher) nicht erforderlich.

&
Vitale Bedrohung

Eine frühzeitige Blutdrucksenkung (\caps{na}, \caps{nfs})
ist notwendig. 

Evtl. Anwendung der Notkompetenz gem. SanG. 
\\\addlinespace\bottomrule
\end{tabularx}
\end{table}
\marginpar{\cite{WHOISH2003,JNC7Report,BLIM3,IntroClinEmergMed4,Vlcek2008}}

\Co{\Z{Die Unterscheidung zwischen "`Hypertensiver
Krise"' und "`Hypertensiven Notfall"' ist in der
Literatur nicht einheitlich. Oft wird auch nur
zwischen "`Hypertensiver Krise ohne Symptome"' bzw. "`mit
Symptomen"' unterschieden oder generell eine andere
Einteilung gewählt. Von diesen Spitzfindigkeiten abgesehen
gilt das Vorhandensein von Symptomen durchgehend als
entscheidendes Kriterium.
\cite{ClassenInnere5,KlinLeitAlle3,Fauci2008}}}
% Renz-Polster2006 p442
% IntroClinEmergMed4 Kap. 28, p 393



\MassnahmenNG{Spezielle Maßnahmen}{Hypertensive Krise}{}
{
\begin{itemize}
  \item Lagerung: Oberkörper hoch.
  \item Evtl. \ce{O2}-Gabe.
  \item Beengende Kleidung öffnen.
  \item Monitoring.
  \item  Hospitalisierung oder Notarzt zur
  Behandlung/Belassung.
\end{itemize}
Eine notfallmäßige Blutdrucksenkung ist nicht erforderlich.
}

\MassnahmenNG{Spezielle Maßnahmen}{Hypertensiver Notfall}{}
{
\begin{itemize}
  \item \TxMassVitBes
  \begin{itemize}
    \item Lagerung: Oberkörper hoch
  \end{itemize}
  \item Beengende Kleidung öffnen
\end{itemize}

Eine frühzeitige Blutdrucksenkung (\caps{na}, \caps{nfs})
ist notwendig. 

\ASBFLO{\Abk{NFS} dürfen lt.
Algorithmus Nitrolingual-Spray verabreichen. Ist die Gabe erfolgreich kann u.U. von der
Anforderung eines Notarztes abgesehen werden.} 
}




\subsection{Gefäßaussackungen: Aneurysmen}

\subsubsection{Bauchaortenaneurysma}
\nocite{Meron2004}

%%%%%%%%%%%%%%%%%%%%%%%%%%%%%%%%%%%%%%%%%%%%%%%%%%%%%%%%%%%
% TODO : Überarbeitung
\TODO{GABS}{????-??-??}{Ticket 34: Erklärung Aneurysma.}{Dieser Teil ist nicht gut formuliert und/oder unvollständig, 
   er benötigt eine Überarbeitung!}
%%%%%%%%%%%%%%%%%%%%%%%%%%%%%%%%%%%%%%%%%%%%%%%%%%%%%%%%%%%
\Layout{0}{\TxBeschreibung{} und Ursachen}
{%%% Gerhard Gabriel <gerhard.gabriel@grissu.org>
Bei der Entwicklung eines Bauchaortenaneurysmas ist zumeist
eine krankhafte Veränderung der Aortenwand
verantwortlich. Als Folge eines Unfallgeschehens,
oder auch nach Hebeversuchen schwerer Lasten, kann die
Gefäßwand im Bereich eines Aneurysmas einreißen. 
}
{\SymbolText}{}{}{}

\Layout{0}{\TxSymptome}{
Je nach
Größe der Schädigung und entsprechendem Schweregrad der Blutung
im Läsionsbereich, kommt binnen kürzester Zeit unter
heftigem, akutem Bauchschmerz zum Tod. Kleinere Läsionen
sind ebenfalls von heftigem, akut einsetzendem Bauchschmerz
begleitet. Es bildet sich ein Blutungsschock (hypovolämischer
Schock) aus, der unbehandelt zum Tode führt.
}{
\begin{itemize}
  \item Heftigste Schmerzen
  ("`Vernichtungsschmerz"').
  \item Hoher Blutverlust ${\rightarrow}$ schnell im
  Schock!
  \item Evtl. bereits bekanntes Aneurysma
  ${\rightarrow}$\,Anamnese.
  \item Evtl. tastbare pulsierende Schwellung.
\end{itemize}
}{}{}{}



Oft ist bei den Patienten das Aneurysma bereits bekannt und
war noch nicht groß genug um behandelt zu werden
(${\rightarrow}$\,Anamnese!). Ansonsten ist es schwer das
Aneurysma präklinisch zu diagnostizieren. Im Vordergrund
steht der Volumenverlust und der entstehende hypovolämische
Schock. 


% \clearpage
\subsection{Gefäßverschlüsse in den Extremitäten}

\Layout{0}{Unterscheidung}
{
Grundsätzlich unterscheidet man je nach betroffenem Gefäß
zwischen \Es{arteriellen} und \Es{venösen}
Gefäßverschlüssen. 

Beim \E{arteriellen Verschluss} besteht
das Hauptproblem in der Unterversorgung der Extremität mit
Blut. Beim \E{venösen Verschluss} kommt es zu einer Ab- und
Rückflussbehinderung. Zusätzlich besteht die Gefahr, dass sich
ein Gerinnsel losreißt und in einem anderen Teil des
Körpers Schaden anrichtet\footnote{Zum Beispiel in der
Lunge im Rahmen einer Lungenembolie}.
}{
\begin{itemize}
  \item \textbf{\textcolor[rgb]{0.7882353,0.0,0.08627451}{Arteriell}}
  \item \textbf{\textcolor{blue}{Venös}}:
  \index{tVT}\T{tVT} (tiefe Venen Thrombose)
\end{itemize}
}{}{}{}



\subsubsection{Arterieller Gefäßverschluß}

\Layout{0}{Beschreibung}
{
Beim \E{arteriellen Gefäßverschluss} wird das
zuführende Gefäß verschlossen. Es kommt zu einer schmerzhaften 
Unterversorgung der Extremität mit Blut. Der komplette
Verschluss passiert plötzlich. 
}
{
\begin{itemize}
    \item Verschluß eines blut\E{zuführenden} Gefäßes.
\end{itemize}
}{}{}{}

\Layout{2}{Ursachen -- Warum kommt es zu einem Gefäßverschluss?}
{%
Meistens liegt eine 
\Margin{Chronische Gefäßschädigung}chronische Schädigung der Gefäße infolge anderer Erkankungen oder aufgrund eines ungesunden
Lebensstils vor. Ursachen für solch eine Gefäßschädigung
sind z.\,B. die Zuckerkrankheit (\T{Diabetes mellitus},
s.\ref{topic:diabetes}), Hypertonie, erhöhte Blutfette oder
das Rauchen. Die Gefäße verkalken und verändern sich auch sonst
nicht zu ihrem Vorteil.
}
{
\begin{itemize}
  \item Chronische Gefäßschädigung, z.\,B. Bei Diabetes mellitus, Hypertonie \ldots
\end{itemize}
}{}{}{}



\Layout{0}{pAVK}{
Wenn die Gefäße nur verengt sind spricht man von
der \Es{peripheren
arteriellen Verschluss\-krankheit}, abgekürzt \T{paVK}. Sie ist eine
häufige chronische Erkrankung, die man oft in Arztbriefen oder auf
Transportscheinen findet.
}{
\begin{itemize}
    \item Periphere arterielle Ver\-schluss\-krank\-heit.
\end{itemize}

}{}{}{}

 (\E{Achtung}: Obwohl es
"`\emph{Verschluss}krankheit"' heißt, sind die Gefäße nicht komplett
verschlossen, sondern nur verengt!)

\Layout{}{}{}{}{}{}{}

\Layout{0}{\TxSymptome}
{
Die Symptome ergeben sich aus der Unterversorgung der
Extremität mit sauerstoffreichem Blut: sie ist blass,
kalt, tut höllisch weh und die Rekap-Zeit ist nicht
mehr sinnvoll messbar. Zusätzlich wird die Extremität
gefühllos und gelähmt. Der Patient kann ein Kribbeln
empfinden (\E{"`Ameisenlaufen"'}. Tabelle \ref{tab:5p}
listet die klassischen Symptome auf \Z{(die \emph{"`5
englischen P's"'})}.
}{
\Captionof{table}{Die 5 Zeichen des arteriellen Gefäßverschluss.}\label{tab:5p}

\begin{tabulary}{\linewidth}{CLL}
\midrule\addlinespace
\centering 1. & \noindent\Z{Pain}        & \noindent\Es{schmerzend}
\\
\centering 2. & \noindent\Z{Pale}        & \noindent\Es{blaß}
\\
\centering 3. & \noindent\Z{Pulsless}    & \noindent\Es{pulslos}
\\
\centering 4. & \noindent\Z{Paresthesia} & \noindent\Es{gefühllos}
\\
\centering 5. & \noindent\Z{Paralysis}   & \noindent\Es{gelähmt}
\\\addlinespace\bottomrule
\end{tabulary}
}{}{}{}


\MassnahmenNG{Spezielle Maßnahmen}{Arterieller Gefäßverschluss}{}{
\begin{itemize}
  \item Lagerung: Extremität hängen lassen, weich und warm lagern.
  \item Schonend transportieren.
  \item Nüchtern.
  \item Hospitalisation: Abteilung für Chirurgie
  (Gefäßchirurgie).
\end{itemize}
}

\subsubsection{Tiefe Venenthrombose}

\Definition{Verstopfung einer Vene durch ein Blutgerinnsel
(Thrombus\footnote{Thrombus: Blutgerinnsel. Im Gegensatz dazu \T{Embolus}: Losgelöster Thrombus, welcher woanders ein Gefäß
verstopft und eine \T{Embolie} verursacht.}).}


\Layout{0}{Ursachen und Risikofaktoren}
{
Ursächlich für die Bildung eines Blutgerinnsels ist meistens eine verminderte
Blutflussgeschwindigkeit, welche wiederum verschiedene Ursachen haben kann: 

Die \Es{venöse Insuffizienz} ist ein weit verbreitetes chronisches Leiden.
 Es kommt  zu einem Rückfluss und Aussackungen von
Venen, welche im fortgeschrittenen Stadium als hässliche \E{Krampfadern} deutlich
sichtbar sind. Schuld sind oft insuffiziente Venenklappen, stehende bzw. sitzende
Tätigkeiten und fehlende Bewegung.

\index{Economy class Syndrom}
Eine länger dauernde \Es{Immobilisation} oder \Es{Bettlägrigkeit} sind auch
typische Ursachen für einen verminderten venösen Blutfluss. Darunter fallen auch
lange Reisen ohne Bewegung: So werden Thrombosen aufgrund von langen Flugreisen
auch scherzhaft als \T{Economy class Syndrom} bezeichnet.

Erschwerend wirkt sich auch eine \Es{unzureichende Flüssigkeitsaufnahme} aus, da
dann der flüssige Anteil des Blutes vermindert sein kann und das Blut folglich
dickflüssiger\ZFootnote{visköser} wird.

Ein besonderer Risikofaktor ist der Zigarettenkonsum. Die Einnahme von
hormonellen Verhütungsmitteln (\Es{Pille}) erhöht zusammen mit Zigarettenkonsum
das Thromboserisiko drastisch.

Im Rahmen anderer Erkrankungen, z.\,B. Tumorerkankungen, kann das Gerinnungsystem
gestört sein, wodurch es zu Thrombosen kommen kann.
}{
\begin{itemize}
  \item Venöse Insuffizienz (Krampfadern), Rückstau. 
  \item Immobilisation, Bettlägrigkeit. 
  \item Lange Reisen ohne Bewegung: (\T{Economy class Syndrom}). 
  \item Mangelnde Flüssigkeitsaufnahme. 
  \item Rauchen, Rauchen + Pille.
  \item Andere Erkrankungen. 
\end{itemize}
}{}{}{}

\Layout{0}{\TxSymptome}
{
Leitsymptom ist eine \Es{schmerzhafte Extremität}, zusammen mit einer (meist
einseitigen) \Es{Schwellung}, \Es{Überwärmung} und einer
dunkel-bläulich-violetten (lividen) \Es{Verfärbung}. 

In der Regel ist nur eine Extremität betroffen, d.h. die Symptome sind nur auf
\E{einer Seite} zu finden. 
}{
\begin{itemize}
  \item Schmerzen.
  \item Schwellung.
  \item Dunkel-bläuliche (livide) Verfärbung. 
  \item Überwärmung. 
  \item Symptome i.\,d.\,R. \E{einseitig}, \E{Seitenvergleich} beachten!
\end{itemize}
}{}{}{}

\Layout{0}{\TxGefahren}
{
Es kann zu einer \E{Loslösung von Teilen des Thrombus} kommen, dieser
losgelöste Teil wird \T{Embolus}\index{Embolus} genannt. Mit dem Blutstrohm gelangt der
Embolus über die Hohlvenen in das rechte Herz. Bis hierhin gibt es noch keine
Probleme, da die Blutgefäße an die Dicke zunehmen. Vom rechten Herzen gelangt
der Embolus jedoch in den Lungenkreislauf (Lungenarterien). Hier verzweigen
sich die Geäße jedoch wieder und der Durchmesser wird kleiner. Dies führt
dazu, dass der Embolus irgendwann zu groß ist und stecken bleibt. Dadurch wird
nun ein  Lungengefäß verstopft und ein Teil der Lunge nimmt nicht mehr am
Gasaustausch teil: Es entsteht eine \Es{Lungenembolie}
(\prettyref{topic:lungenembolie}).
}{
\begin{itemize}
  \item Lösung des Thrombus →
  \Es{Lungenembolie} (\prettyref{topic:lungenembolie})
\end{itemize}
}{}{}{}




\begin{gWideParEnv}
\Dias{Bilderserie: Thrombosen und die "`echte Welt"'.}
%1
{\includegraphics[width=\linewidth]{./images/economyclass.jpg}
\Captionof{figure}{Die Economyclass. Sorgsam geschlichtet verbringen Menschen
hier Stunden damit Thrombosen zu basteln.
\SourcePicture{Gabriel} \AuthNotesPicLegalOk{}{}}}
%2
{\includegraphics[width=\linewidth]{./images/lovenox1.jpg}
\Captionof{figure}{Thromboseprophylaxe. \Z{Niedermolekulares Heparin
(hier "`Lovenox\texttrademark") verhindert Thrombosen die z.\,B. durch lange
Immobilisation (Reisen, Bettlägrigkeit, Gips, \ldots)
entstehen können. Die Substanz wird unter die Haut ("`subkutan"') gespritzt.
\SourcePicture{Gabriel} \AuthNotesPicLegalOk{}{}}}}
%3
{~}
\end{gWideParEnv}






\MassnahmenNG{Spezielle Maßnahmen}{Beinvenenthrombose}{}{
\begin{itemize}
  \item Lagerung:  liegend, betroffenes Bein
  hochlagern.
  \item Bewegungsverbot (Emboliegefahr!).
\end{itemize}
}

\section{Respiratorische Notfälle}


% \begin{center}
%  \marginpar{\includegraphics[width=\linewidth]{./images/noauth/AASdev20090228-img75.png}
%  \AuthNotesPicLegalNotOk{img75}{Splash Pulmo Notfälle }}
% 
% \end{center}

\minisec{Querverweise}

\begin{itemize}
  \item Lungenentzündung (Pneumonie): \prettyref{topic:pneumonie}
  \item Tuberkulose: \prettyref{topic:tuberkulose}
\end{itemize}



% \clearpage
\subsection{Mechanische Atemwegsverlegung}
\label{topic:mechanische-atemwegsverlegung}

\Lang{Syn.} Fremdkörperaspiration.

% TODO Fremdkörperaspiration: Erwachsen -- Kind überarbeiten und zusammenführen
% (# 39)

\Layout{20}{\TxBeschreibung}
{%
\index{Aspiration}%
Wenn ein Fremdkörper den Atemweg teilweise oder vollständig  \E{blockiert}
oder \E{behindert}, spricht man von einer mechanische Verlegung. Abhängig von
der Größe des Fremdkörpers können sowohl die oberen als auch die unteren
Atemwege betroffen sein. Wenn der
Fremdkörper in die \E{unteren Atemwege} gelangt, spricht man von einer
(Fremdkörper-) \T{Aspiration}. 

Nicht nur feste Körper, sondern auch Flüssigkeiten können Auslöser
einer Verlegung oder Aspiration sein. Besonders Blut, Magensaft, Erbrochenes
oder Speisebrei stellen diesbezüglich eine Gefahr dar.

Besonders gefährdet hinsichtlich einer Atemwegsverlegung sind Kinder, ältere
Personen, bewusstseinseingschränkte Personen, Patienten mit Schluckstörungen oder
Allergiker (z.\,B. Anschwellen der Schleimhäute nach einem Bienenstich).

}
{%
\begin{itemize}
  \item Verlegung durch Fremdkörper.
  \item Obere und untere Atemwege möglich.
  \item Aspiration: untere Atemwege.
  \item Besonders gefährdet:
  \begin{itemize}
    \item Kinder oder alte Personen.
    \item Bewusstseinseingeschränkte Personen.
    \item Patienten mit Schluckstörungen.
    \item Allergiker.
  \end{itemize}
\end{itemize}
}
{./images/hirtler-aass/Fremdkoerper.pdf}
{Atemwegsverlegung (Schema). Der linke Bissen steckt in der Luftröhre fest, der
andere befindet sich in der Speiseröhre.}
{Hirtler.}



\Layout{60}{\TxSymptome}
{%
Die Symptome sind abhängig  von der Schwere der Verlegung. 
Man unterscheidet zwischen einer \Es{schweren} und einer \Es{milden
Atemwegsverlegung} \cite{Erc2005Sec2,Erc2005Sec6,Erc2005DeSec2,Erc2005DeSec6}.

\subparagraph{Schwere Verlegung}
Bei einer schweren Verlegung der Atemwege ist der Patient  \Es{unfähig zu reden
oder effektiv zu husten}. Eventuell kann er nur durch Nicken oder Gesten
kommunizieren. In besonders schweren Fällen kann es zu einem \Es{Atemstopp} und
nach kurzer Zeit zu einer \Es{Bewusstlosigkeit}, gefolgt von einem
\Es{Kreislaufstillstand} kommen.


\subparagraph{Milde Verlegung} Bei der milden Verlegung kann der Patient noch
sprechen und atmen. Er klagt meist über Atemnot und ein Fremdkörpergefühl, oft
versucht er zu Husten oder zu würgen. Oft ist auch ein pfeifendes Atemgeräusch
während der Ein- und/oder Austamung zu hören (in-/exspiratorischer \E{Stridor}).



\Fazit{Die (Un-) Fähigkeit zu sprechen und zu husten ist das wesentliche
Kriterium um die Schwere einer Atemwegsverlegung einzuschätzen.} 
}{
\begin{itemize}
  \item Schwere Verlegung:
  \begin{itemize}
    \item Unfähigkeit zu Sprechen oder effektiv zu husten.
    \item Evtl. Bewusstlosigkeit → Kreislaufstillstand
  \end{itemize}
  \item Milde Verlegung:
  \begin{itemize}
    \item \Es{Kann Sprechen \& Husten}.
    \item Atemnot. 
    \item Fremdkörpergefühl.
    \item Pfeifendes Atemgeräusch (Stridor).
  \end{itemize}
\end{itemize}
} 
{./images/wikipedia-pd/Abdominal_thrusts3.jpg} 
{Heimlich-Manöver.} 
{U.S. Naval Hospital, Okinawa, Japanese intern Naoko Uehara teaches
proper technique for treating a choking victim with the assistance of Hospital
Corpsman Second Class Ross Francis. Amanda M. Woodhead, USMC / WmCo, PD.}



\Layout{60}{Heimlich-Manöver}
{%
\Lang{Syn.} Heimlich-Handgriff. Beim Heimlich-Manöver versucht
man durch Kompression Bauchraumes den Fremdkörper durch den entstandenen
Überdruck aus den Atemwegen zu entfernen.  

Der Helfer umfasst von hinten den Oberbauch des Patienten und bildet mit
der einen Hand eine Faust und unschliesst diese mit der anderen Hand. Nun legt
er die Faust unterhalb der Rippen und des Brustbeins an und zieht sie dann
ruckartig gerade nach hinten zu sich. 

\Ca{Verletzungsgefahr!} Durch Anwendung des Heimlich-Handgriffes kann es zu
schweren inneren Verletzungen kommen, die Anwendung ist daher auf bedrohliche Situationen
beschränkt.
}{
\begin{itemize}
  \item Kompression des Bauchraumes → Überdruck.
  \item Ansatzpunkt unterhalb des Rippenbogens, ruckartige Bewegung.
  \item \Ca{Verletzungsgefahr!}
\end{itemize}
}
{./images/wikipedia-pd/US_medic_teaches_the_Heimlich_manuever_to_laughing_Afghans_00500.jpg}
{Heimlich-Manöver.}
{"`Air Force Staff Sgt. Abraham Jara provides instruction in the proper
application of the Heimlich maneuver to teachers at the Kharana Girls Middle
School, Dara district, Panjshir province, Afghanistan."' Photo by 
Capt. John T. Stamm, USAF / WmCo, PD.}




 

 
\Layout{2}{Schläge zwischen die Schulterblätter}
{%
Durch kräftige Schläge auf den Rücken zwischen die Schulterblätter kann  ein
festsitzender Fremdkörper evtl. gelockert und ausgehustet werden.}
{}{}{}{}
 

 

 

 

 

 

 

 

 

 

 

 





\MassnahmenNG{Spezielle Maßnahmen}{Mechanische Atemwegsverlegung}{}
{%
\begin{itemize}
  \item Atemwege soweit möglich freimachen.
  \item \TxBeiVit \TxMassVitBes
  \begin{itemize}
    \item Lagerung: Situationsgerecht, z.\,B. Oberkörper erhöht. 
    
    Bei Schlag
    zwischen die Schulterplätter soll der Patient abwärts schauen.
    \item \ce{O2} 10\LMin, bei Bedarf bis zu 15\LMin .
  \end{itemize}
  \item Maßnahmen abhängig von der Beurteilung der Schwere der
  Atemwegsverlegung\ZFootnote{gem.
  \cite{Erc2005DeSec2}}:
\end{itemize}

\noindent\begin{tabularx}{\linewidth}{XXX}
\toprule\addlinespace
\multicolumn{3}{c}{\TabSub{Beurteilung der Schwere der Atemwegsverlegung}}
\\\addlinespace\midrule\addlinespace
\TabSub{Schwer}

Kann nicht sprechen

Hustet ineffektiv
&
&
\TabSub{Mild}

Kann sprechen

Hustet effektiv
\\\addlinespace\cmidrule(lr){1-2}\cmidrule(lr){3-3}\addlinespace
\TabSub{Ohne Bewusstsein}
&
\TabSub{Bei Bewusstsein}
&
\TabSub{Bei Bewusstsein}
\\\addlinespace\cmidrule(lr){1-1}\cmidrule(lr){2-2}\cmidrule(lr){3-3}\addlinespace
5 Initialbeatmungen

Reanimation
&
5 Schläge zwischen Schulterblätter

5 Heimlich-Manöver (\textgreater{}\,1\,a) oder Thoraxkompressionen
(\textless{}\,1\,a)

Wiederholung 
&
Husten lassen

Überwachung
\\\addlinespace\bottomrule\addlinespace
\end{tabularx}





Durch Anwendung des Heimlich-Handgriffes kann es zu
    schweren inneren Verletzungen kommen. Ein Patient ist nach
    Anwendung unbedingt zu hospitalisieren!

}

\Layout{2}{\dsliterary}{%
\cite{Wuttig2009}
}{}{}{}{}



\subsection{Asthma bronchiale}
\label{topic:asthma-bronchiale}

\Definition{Chronische Erkrankung mit mehr oder weniger
  häufigen Episoden von Anfällen mit massiver Atemnot oder
  Atembehinderung durch Verengung der Bronchien.}

% \DefinitionExp{Asthma_bronchiale}

\Layout{0}{\TxBeschreibung{}}
{
Im akuten Asthmaanfall kommt es zu einer Verengung der
Bronchien und vermehrter Schleimbildung, die \Es{Ausatmung}
wird massiv erschwert und der Patient hat das Gefühl nicht genug Luft zu bekommen.
}{%
Wiederholte Anfälle von Atemnot durch

\begin{itemize}
  \item Verengte Bronchien.
  \item Vermehrte Schleimbildung.
  \item Behinderte Ausatmung.
\end{itemize}
}
{}
{}
{}





Asthma ist bei den meisten Patienten oft schon bekannt und
kann bzw. muss erfragt werden. Viele Patienten haben auch
einen oder mehrere Sprays oder Inhalationsgeräte zur Dauer-
oder Akutbehandlung. 

\AuthNotes{GABS: Gabe von
Beta-Mimmetikum-Spray muss noch diskutiert werden.

KOCH: nein nur NKA

GABS: wenn, dann NFS (Medi-Liste 1).

Das Problem ist aber: Darf der Patient seinen (vom Arzt
vielleicht gerade für diese Situation
\underbar{verschriebenen}) Spray nehmen? wann darf er und
wann nicht? Macht ja keinen Sinn wenn man dem Pat seinen
Spray generell vorenthält.}

Besonders betroffen vom Asthma bronchiale sind Patienten
mit bekannten Allergien. Folglich kann alles, was bei dem
Patienten eine Allergie verursacht auch einen Asthmaanfall
auslösen (Pollen, Gräser, Katzen, Zigarettenrauch, \ldots). 

Ist bei Eintreffen des Rettungsdienstes noch keine
Besserung der Atemnot eingetreten ist diese Atemnot als
massiv und lebensbedrohend einzuschätzen. Es ist nicht
damit zu rechnen, dass der Zustand von alleine wieder
vergeht und bedarf einer medikamentösen Behandlung durch einen
Notarzt. 

\Co{Unterschätze nie die Schwere eines Asthmaanfalles! Gefährlich sind
besonders die "`ruhigen"', wenig klagenden Patienten. Asthma-Patienten haben
oft ein anderes Luftnot-Empfinden. Diese Patienten sind
HochrisikoPatienten und machen auch bei der \ce{O2}-Gabe
Probleme! \ce{O2} nur bis max. 95\% Sauerstofsättigung
!\cite{Olschewski2009}\A{\footnote{Vortrag Univ-Prof.Dr.
Horst Olschewski, 5. AGN-Kongress 2009; Notfall Rettungsmed (2009) 12:322}}}


\A{Foto Inhaltation V}


\Layout{0}{\TxSymptome}
{%
Das Leitsymptom ist die Atemnot. Oft ist ein pfeifendes,
brummendes oder giemendes Atemgeräusch besonders während der
Ausatmung (\E{Exspiration}, \E{exspiratorischer Stridor})
wahrnehmbar. Die Exspiration ist ausserdem verlängert. Trockerner Husten ist häufig.

Der Patient nimmt meist automatisch und unbewusst eine
Position ein die seine Atemhilfsmuskulatur unterstützt: Er
sitzt aufrecht und stützt sich oft nach vorne oder hinten
ab. 

Die Atemfrequenz ist erhöht, oft auch die Herzfrequenz
(Tachypnoe und Tachykardie). Die Sauerstoffsättigung kann
erniedrigt sein. 

\subparagraph{\Ca{Schwerer Anfall}}

Im schweren Anfall ist oft \Es{kein}, oder nur mehr ein
\Es{leises Atemgeräusch} hörbar (da kaum mehr
Luft bewegt wird, \Z{"`silent chest"'}). Es kommt zu einer
Kreislaufdepression, Herzfrequenz und Blutdruck stürzen ab.
Dies kann zu Bewusstseinsstörungen führen. Die
Sauerstoffsättigung ist stark vermindert und der Patient ist zyanotisch.


}{
\begin{itemize}
  \item Akute Atemnot mit trockenem Husten.
  \item Exspiratorisches Atemgeräusch (Stridor).
  \item Einsatz der Atemhilfsmuskulatur.
  \item Verlängerte Ausatmungsphase.
  \item Schwerer Anfall:
  \begin{itemize}
    \item Kein oder leises Atemgeräusch.
    \item HF, RR $\downarrow$.
    \item Sp\ce{O2} $\downarrow\downarrow$.
    \item Bewusstseinsstörungen.
  \end{itemize}
\end{itemize}
}
{}
{}
{}

% \Layout{2}{\TxKomplikationen}
% {
% \begin{itemize}
%   \item Zyanose
%   \item Status asthmaticus
%   \item Panik
% \end{itemize}
% }{
% 
% }
% {}
% {}
% {}


\paragraph{Status Asthmaticus}
\index{Status!asthmaticus}\label{topic:status-asthmaticus}Der
"`Status"' ist eine gefürchtete Komplikation des
Asthmaanfalles bei der die Atemnot nicht wieder vergeht. Es
besteht die Gefahr, dass der Patient nach einiger Zeit keine
Kraft mehr zum Atmen hat (Erschöpfung). Der Status
asthmaticus kann tödlich enden.

% \begin{figure}[h]
%     \includegraphics{./images/copd-bronchus.png}
%     
%     \Captionof{figure}{Bronchus bei COPD.
%     \AuthNotesPicLegalNotOk{copd-bronchus}{}}
% \end{figure}



\MassnahmenNG{Spezielle Maßnahmen}{Akuter Asthmaanfall}{}{
Einschätzen der vitalen Bedrohung.

\Es{\TxBeiVit} \TxMassVitBes

\begin{itemize}
  \item Lagerung: s.\,u.
  \item \ce{O2}: s.\,u.
\end{itemize}

\Es{weiters:}

\begin{itemize}
  \item Lagerung: Oberkörper hoch, evtl. abstützen
  lassen (Atemhilfsmuskulatur unterstützen)
  \item \ce{O2} bis eine Sauerstoffsättigung von 95\,\% erreicht ist,
  \Es{nicht mehr!} Wenn keine Pulsoxymetrie vorhanden ist 8\LMin.
  
  \Ca{\mdseries Achtung: In seltenen Fällen
  kann es vorkommen, dass bei hochdosierter
  Sauerstoffgabe der Patient zu atmen aufhört! Siehe "`\Abk{COPD}"'
  (\prettyref{topic:copd}). Bei sorgfältiger Überwachung
  des Patienten stellt die \ce{O2}-Gabe in der Praxis
  jedoch kein Problem
  dar\cite{ScholzNotfallmedizin2}.}
  
  \item Vitalwerte erheben \& Monitoring
  \item Beruhigen, voratmen, durch fast
  geschlossene Lippen ausatmen lassen (Lippenbremse)
  \item Sprays, die Patient evtl. bei sich hat dürfen
  nicht mehr genommen werden.
\end{itemize}
\AuthNotes{\#14: Einnahme von verordneten Sprays analog
zu Nitro?\par $O2$-Gabe?}

 }

\Layout{22}{}
{
Patientin mit einem akuten Asthmaanfall.

Sie sitzt auf einer Treppe und stützt sich nach hinten
mit den Armen ab. Die Erstmaßnahmen bei vital bedrohten Patienten wurden ergriffen:
\begin{compactitem}
  \item Situationsrechte Lagerung.
  \item Sauerstoffgabe.
  \item Reanimationsbereitschaft\footnote{Defibrillator eingeschaltet,
  sonstiges Zubehör nicht am Foto sichtbar}.
  \item Notarztnachforderung\footnote{Paramedic
  vor Ort, am Foto nicht sichtbar.}.
  \item Vitalwerte erheben und Monitoring:
  Blutdruck wird gerade gemessen, die Patientin
  ist am Monitor angeschlossen.
\end{compactitem}
}{

}
{./images/asthma-paramedic.jpg}
{Patientin mit einem akuten Asthmaanfall.\label{fig:asthma-paramedic}}
{Gabriel.}





\subsection{Chronische Bronchitis und 
Chronisch-obstruktive Lungenerkrankung
(COPD)}
\label{topic:copd}
\nocite{Renz-Polster2006} 

% Renz-Polster2006 p442

\emph{\Lang{Syn.} \T{Chronic obstructive pulmonary disease}}

\Definition{\T{Chronische Bronchitis}:
Chronisch entzündliche Schleimhautschädigung  der Atemwege.}
\Definition{\T{COPD}: Chronisch-entzündliche
Schleimhautschädigung, welche durch inhalative Noxen
verursacht wird und im Verlauf eine zunehmende obstruktive
Atemwegseinschränkung aufweist.}

Am Anfang steht die chronische Bronchitis, welche durch
Husten mit schleimigem Auswurf gekennzeichnet ist
("`Raucherhusten"'). Es kommt dabei zu einer gesteigerten
Entzündungsantwort auf eingeatmete Stoffe (Zigarettenrauch,
Umweltschadstoffe, \ldots). Wenn man in der
Lungenfunktionsuntersuchung eine Atemwegseinschränkung
nachweisen kann, spricht man von der chronisch-obstruktiven
Lungenerkrankung (\Abk{COPD}). 


\Layout{0}{\TxSymptome}{
Die \Abk{COPD} ist durch eine voranschreitende
Verschlechterung der Atemleistung gekennzeichnet, erstes
Anzeichen ist ein Husten mit Auswurf, der klassische
"`Raucherhusten"'.

Je nach
Schweregrad kommt es zu Zeichen der Atemwegsverlegung
(Obstruktion) und Ateminsuffizienz:
Die \E{Ausatmung} (Exspiration) ist erschwert, der Patient klagt
über \E{Belastungsdyspnoe}, ein Engegefühl, \E{Husten} und hat Tachypnoe und
evtl. Zyanose. Man kann oft ein \E{brummendes oder giemendes Atemgeräusch}
hören.

Durch die erschwerte Ausatmung kommt es zu einer chronischen Überblähung der
Lungenbläschen und zu einer "`Faß-förmigen"' Verformung des Brustkorbes. Im
Endstadium zeigen sich Zeichen einer Rechtsherzinsuffizienz (\prettyref{topic:rechtsherzinsuffizienz}) aufgrund einer Störung im Lungenkreislauf.

Bei fortgeschrittener Ateminsuffizienz bekommt der Patient oft eine
\E{Heimsauerstofftherapie} verordnet. 
}{
\begin{itemize}
  \item Husten mit Auswurf, anfänglich
  "`Raucherhusten"'.
  \item Voranschreitende, sich verschlechternde \E{Ateminsuffizienz}:
  \begin{itemize}
    \item Belastungsdyspnoe (bei schwererer \Abk{COPD}
    auch in Ruhe).
    \item Ausatmung (Exspiration) erschwert.
    \item Evtl. chronische Zyanose.
    \item Brummendes, giemendes Ausatemgeräusch.
    \item Benötigt evtl. Heimsauerstoff.
  \end{itemize}
\end{itemize}
}{}{}{}



\Layout{0}{Exazerbationen}{
Der \Abk{COPD}-Patient erleidet öfters plötzliche
Verschlechterungen, sog. \E{Exazerbationen}. Diese sind
meist durch Infekte ausgelöst und durch vermehrte Atemnot,
Husten und Auswurf gekennzeichnet.
}{
\begin{itemize}
  \item Plötzliche Verschlechterungen.
  \item Meist Infekt-bedingt.
  \item Vermehrte Atemnot,
  Husten und Auswurf.
\end{itemize}
}{}{}{}


% 
% \subparagraph*{Sammelbegriff für}

% \begin{itemize}
%     \item Lungenemphysem (Überblähung der Lungenbläschen)
%     \item chron. obstruktive Bronchitis
% \end{itemize}




\begin{figure}[H]
    
     \includegraphics[width=\linewidth]{./images/hirtler-aass/Alveolen-Bronchus-COPD-edited2.png} 
     \Captionof{figure}{\AuthNotesPicLegalOk{}{}Veränderung der
     Atemwege bei der COPD. \TabSub{Links:} Schema eines gesunden
     Bronchus und einer gesunden Alveole. \TabSub{Oben rechts:} Bei der COPD sind
     die kleinen Luftwege verschleimt und verengt. \TabSub{Unten rechts:} Die
     Lungenbläschen (Alveolen) sind überbläht, weil die Luft nur erschwert wieder entweichen kann.
 \SourcePicture{Hirtler.}}
\end{figure}





\begin{figure}[h]
\begin{center}
% \includegraphics{./images/copd-stufen.png}
\includegraphics[width=\linewidth]{./images/hirtler-aass/COPD-Stufen_edited.pdf}

\Captionof{figure}{Eine COPD
entsteht nicht plötzlich: Ein COPD{\textquotesingle}ler hat eine
"`Karriere"' hinter sich. \SourcePicture{Hirtler}}
\end{center}
\end{figure}


\Layout{0}{Probleme mit Sauerstoff bei
COPD-Patienten}
{
\index{Sauerstoff!COPD-Patient}
Normalerweise wird der Atemantrieb über die
\ce{CO2}-Konzentration im Blut gesteuert. Die
\ce{CO2}-Konzentration ist bei \Abk{COPD}-Patienten jedoch chronisch
hoch, der Körper gewöhnt sich an diesen Zustand. Er
regelt dann den Atemantrieb direkt über den \ce{O2}-Gehalt im Blut.

Wird der \ce{O2}-Gehalt plötzlich künstlich erhöht, fehlt
der \index{COPD!Atemantrieb}Atemantrieb und der Patient kann
\Es{Bewusstseinsstörungen} ("`\ce{CO2}-Narkose"') und einen
\Es{Atemstillstand} erleiden.
}{
\Ca{Achtung!}

\SymbolText
}{}{}{}




\MassnahmenNG{Spezielle Maßnahmen}{COPD-Exazerbation}{}{
\begin{itemize}
  \item
  Wenn bei Eintreffen noch keine Besserung
  eingetreten ist, ist von einer vitalen Bedrohung
  auszugehen, \TxMassVit .
  \item \Es{\ce{O2}-Gabe:}
  
  \Ca{\mdseries Sauerstoff vorsichtig dosieren,
  anfänglich nur \VonBis{2}{3}{\LMin} und Atmung überwachen.
  
  Wenn der Patient bereits Heimsauerstoff benutzt, kann
  man bei Atemnot versuchen \VonBis{1}{2}\LMin{} höher zu dosieren. Die Atmung muss
  dabei überwacht werden!
  
  Heimsauerstoff soll weiter gegeben werden!}
  
  \item Lagerung: Oberkörper hoch,
  \item voratmen, Lippenbremse.
  \item Eigener Spray darf nicht mehr genommen werden.
  
\end{itemize}
\AuthNotes{Spray bei COPD -- Regelung wie bei Nitro??}
}

\Layout{2}{\Z{Literatur}}
{
\Z{
\begin{inparadesc}
  \item[Kasuistiken:] \cite{Knacke200601}
\end{inparadesc}
}
}{}{}{}{}


\subsection{Lungenembolie}\label{topic:lungenembolie}

\Definition{Einschwemmen eines Gerinnsels in die  Strombahn des Lungenkreislaufs}


\Layout{20}{Herkunft des Embolus}
{%
Der Embolus, der die Lungenarterien verstopft, entsteht
normalerweise nicht in den Lungengefäßen selbst, sondern er
wird angeschwemmt. Sehr oft löst sich ein Thrombus bei
einer tiefen Beinvenenthrombose und schwimmt im Blut bis in
die Lunge. Daher gelten die
Risikofaktoren der tiefen Beinvenenthrombose
(Immobilisation, lange Reisen, \ldots) auch bei der
Lungenembolie. 
}{

}
{./images/hirtler-aass/Thrombus-Embolie.pdf}
{Herkunft des Thrombus.}
{Hirtler.}



\marginlabel{\begin{itemize}\item Tiefe
(Bein-)Venenthrombose\item Immobilisation (Gips, lange Reise)\item Herz: Bei Rhythmusstörungen\item Tumorpatienten
\end{itemize}}Auch bei Herzrhytmusstörungen wie einem
Vorhofflimmern kann ein Thrombus im Vorhof selber entstehen, 
sich losreißen und Lungengefäße verstopfen. Bei
Gerinnungstörungen oder auch Krebserkrankungen kann das
Gerinnungssystem verrückt spielen und Thromben bilden die
zu einem Embolus werden.

\Layout{0}{\TxSymptome}
{%
\begin{itemize}
  \item Atemnot, Tachypnoe
  \item Akuter Thoraxschmerz, eher atemabhängig (Atemstopp)
  \item Angst
  \item Tachykardie
\end{itemize}

Eine Lungenembolie kann man leicht mit einem Herzinfarkt oder Angina
pectoris verwechseln. Dies hat aber bei der Erstversorgung kaum
Konsequenz, die Maßnahmen sind anfänglich die Gleichen.


}{
\begin{itemize}
  \item Atemnot, Tachypnoe.
  \item Akuter Thoraxschmerz, eher atemabhängig.
  \item Angst.
  \item Tachkardie.
\end{itemize}
}{}{}{}


\Layout{2}{Typische Befunde}
{%
\begin{itemize}
  \item Zeichen einer tiefen Beinvenenthombose (mögliche
  Quelle des Embolus): Bein: einseitig geschwollen,
  überwärmt, rot/rot-bläulich verfärbt, schmerzhaft
  \item evtl. gestaute Halsvenen
\end{itemize}
}{

}{}{}{}


\Layout{0}{Typische Anamnese}
{%
\begin{itemize}
  \item Immobilisation oder Operation innerhalb
  der letzten 4 Wochen (\Es{Lange Reise}, Flug, Auto,
  Bahn, Bettlägrigkeit, Gips, {\dots}
  \item \Z{Schwangerschaft}
  \item \Z{Raucher(in)}
  \item \Z{Verhütungspille}
\end{itemize}
}{
\begin{itemize}
  \item Immobilisation.
  \item \Z{Schwangerschaft.}
  \item \Z{Raucher(in).}
  \item \Z{Verhütungspille.}
\end{itemize}
}{}{}{}


\Layout{2}{Wann ist die Diagnose \emph{"`Lungenembolie"'}
 wahrscheinlich?}
{%
nach \cite{Wells2000}

\begin{itemize}
  \item {\bfseries Klinischer Verdacht} basierend auf
  Erfahrungswerten.
  \item {\bfseries Andere Diagnosen sind weniger
  wahrscheinlich.}
  \item Herzfrequenz \textgreater 100\,\nicefrac{}{min}
  (Tachykardie).
  \item Immobilisation oder Operation \Z{innerhalb
  der letzten 4 Wochen} (\Es{Lange Reise}, Flug, Auto,
  Bahn, Bettlägrigkeit, Gips, {\dots}
  \item Bereits einmal durchgemachte tiefe
  Beinvenenthrombose oder Lungenembolie.
  \item Bluthusten.
  \item Krebserkrankung \Z{innerhalb der letzten 6 Monate}.
\end{itemize}

Die ersten beiden, fett gedruckten, Kriterien sind die
wichtigsten Kriterien. 
}{

}{}{}{}


\MassnahmenNG{Spezielle Maßnahmen}{Lungenembolie}{}{
\begin{itemize}
  \item \TxMassVitBes .
  \begin{itemize}
    \item Lagerung: Oberkörper hoch.
    \item \ce{O2} 10\LMin, bei Bedarf bis zu 15\LMin .
  \end{itemize}
  \begin{itemize}
    \item Striktes Bewegungsverbot.
    \item Beengende Kleidung öffnen.
  \end{itemize}
\end{itemize}
}


\subsection{Lungenödem}\label{topic:lungenoedem}

\Definition{Flüssigkeitsansammlung durch Stau in den
Alveolen und im Zwischengewebsraum der Lunge.}
\nopagebreak

\Layout{2}{Ursachen}
{%
\begin{itemize}
  \item Blutrrückstau: akute Linksherzinsuffizienz,
  Linksherzversagen.
  \item z.\,B. bei MCI, zusätzliche Belastung bei vorgeschädigtem Herz
  (Infekt).
  \item Durch Entzündung: Inhalation toxischer Substanzen.
\end{itemize}
\Cave{Das gefährliche am Lungenödem sind vor allem dessen Ursachen!}
}{

}{}{}{}



\Layout{2}{\TxSymptome}
{%
\begin{itemize}
  \item Atemnot.
  \item Zyanose.
  \item \index{Brodelnde Atemgeräusche}Brodelnde Atemgeräusche, oft
  ohne Hilfsmittel hörbar, im Extremfall schon von der Wohnungstür
  aus.
  
  \T{Brodelnde Atemgeräusche}: Als ob jemand mit einem
  Strohhalm in ein Glas Wasser bläst.
\end{itemize}
}{

}{}{}{}



\MassnahmenNG{Spezielle Maßnahmen}{Lungenödem}{}{
\begin{itemize}
  \item \TxMassVitBes .
  \begin{itemize}
    \item Lagerung: Oberkörper hoch, \E{wenn möglich Beine
    herabhängen lassen}.
    \item \ce{O2} 10\LMin, bei Bedarf bis zu 15\LMin .
    \item  Notarzt!
    \item Reanimationsbereitschaft!
  \end{itemize}
  \item Striktes Bewegungsverbot.
  \item Beengende Kleidung öffnen.
  \item \Z{Evtl. Gabe von Nitroglycerin-Spray durch einen
  Notfallsanitäter}.
\end{itemize}
}

Patienten mit einem Lungenödem (v.a. bei kardialer Ursache) sind
\Es{lebensbedrohlich krank}. Das Herz verbraucht oft
schon die allerletzten Reserven, schon die Aufregung durch den
Transport im Stiegenhaus zum Auto kann das Herz endgültig in den
Herz-Kreislaufstillstand treiben. 

Es kommt häufig vor, dass diese Patienten \Es{im
Stiegenhaus} oder \Es{im Auto}
\Es{reanimationspflichtig} werden. 

\Cave{Kein Transport ohne Notarzt! 

Auch wenn das Spital 
"`um{\textquotesingle}s Eck"' 
ist. Der Patient schafft es möglicherweise nicht 'mal  mehr 
ins Auto!}



\subsection{Hyperventilationstetanie}
\label{topic:hyperventilationstetanie}
\index{Hyperventilation!-stetanie}
\index{Hyperventilation!psychisch bedingte --}



\Layout{0}{\TxBeschreibung}{%
Durch die Atmung wird der Säure-Basen-Haushalt
des Körpers wesentlich beeinflusst (\prettyref{topic:saeure-basen-haushalt}). Atmet
ein sonst gesunder Mensch unnatürlich tief und schnell, so wird sein Blut-pH
sehr rasch basisch. Dadurch geraten die Elektrolyte im Blut in ein
Ungleichgewicht, es kommt zu Symptomen. 
}{
\begin{itemize}
  \item Tiefe und schnelle Atmung bei sonst gesundem Menschen
  \item Blut-pH wird basisch $\rightarrow$
\end{itemize}
}{}{}{}

\Layout{0}{\TxSymptome}
{%
Die Patienten beklagen \Es{Atemnot}, obwohl die \ce{O2}-Zufuhr nicht
beeinträchtigt und das Blut maximal gesättigt ist (Sp\ce{O2} von 100\,\%). Sie
haben eine \Es{tiefe und schnelle Atmung} (\E{Tachypnoe}).

Durch das Elektrolyt-Ungleichgewicht kommt es zuerst zu Schwindelgefühlen, etwas
später zu einem Kribbeln in den Fingern. In weiterer Folge zeigen sich tonische
Krämpfe, klassisch ist die sog. \T{Pfötchenstellung}\index{Pfötchenstellung}:
Die Arme sind angezogen und Handhaltung erinnert an Pfoten eines Tieres. Wird
die Hyperventilation nicht rechtzeitig beendet, kann es zur
Eintrübung bis hin zur Bewusstlosigkeit kommen.

}{
\begin{itemize}
  \item Subjektiv: \emph{"`Atemnot"'}, Sp\ce{O2} von 100\,\% .
  \item Tiefe, schnelle Atmung.
  \item Schwindel, Bewusstseinsstörungen.
  \item Kribbeln in den Fingern.
  \item Tonische Krämpfe ("`Pfötchenstellung"').
\end{itemize}
}{}{}{}

\Layout{0}{Ursachen}
{%
Die Hyperventilationstetanie ist meistens psychisch bedingt und meist Folge
einer \Es{psychischen Belastungsreaktion} bzw. Folge einer Panikattacke. Auch
im Rahmen von (Pop-)Konzerten und ähnlichen Veranstaltungen sind oft Patienten mit
Hyperventilationstetanie anzutreffen.

}{
\begin{itemize}
  \item Meist psychisch bedingt.
  \begin{itemize}
    \item Psych. Belastungsreaktion, Panikattacke.
    \item Veranstaltungen.
    %   \item Lungenembolie
    %   \item diabetische Stoffwechselentgleisung (Ketoazidose, kompensatorisch)
  \end{itemize}
\end{itemize}
}{}{}{}

\Layout{0}{Differentialdiagnosen}
{%
Die psychisch bedingte Hyperventilation muß unbedingt von der Hyperventilation
aufgrund anderer Ursachen abgegrenzt werden. Bei \Es{Ateminsuffizienz} oder
\Es{durch Erkrankungen verursachte Atemnot} kommt es natürlich auch oft zu
einer (kompensatorischen) Hyperventilation (Herzinfarkt, Lungenembolie, Pneumothorax,
Asthma-Anfall, \ldots). 

Bei einer stoffwechselbedingten \Es{Übersäuerung} (Azidose) kommt es zu einer
Hyperventilation um \ce{CO2} abzuatmen (dadurch wird Säure im Körper abgebaut). Die \E{Kussmaul'sche Atmung} beim \Es{diabetischen Koma} wäre ein typisches
Beispiel dafür (\prettyref{topic:koma-diabetisches}).

}{
\begin{itemize}
  \item Atemnot mit körperlicher Ursache. 
  \item Ateminsuffizienz. 
  \item Stoffwechselbedingte Azidose (z.\,B. Diabetisches Koma).
\end{itemize}
}{}{}{}


% \paragraph{\TxSymptome}
% 
% \begin{itemize}
% \item Unruhe, Aufregung, evtl. Angstgefühle, Panik
% \item Tachypnoe mit Alkalose $\rightarrow$ durch Änderung
% des pH-Werts Änderung des Kalziumspiegels im Blut, dadurch
% \begin{itemize}
%     \item Kribbeln, 
% \item Krämpfe (Gesichtsmuskulatur),  Pfötchenstellung
% \end{itemize}
% \end{itemize}

\MassnahmenNG{Spezielle Maßnahmen}{Hyperventilationstetanie}{}{
\begin{itemize}
  \item Psychische Betreuung: Patient beruhigen, mit ruhiger, langsamer Stimme
  mit dem Patienten sprechen.
  \item Kommandoatmung.
  \item \ce{CO2}-Rückatmung: In ein Plastiksackerl über kurze Zeit Ein- und
  Ausatmen lassen. Ebenso eignet sich eine \Es{\ce{O2}-Maske mit Reservoir}. 
  Handschuh)
  \item Grundsätzlich sollte kein \ce{O2}  verabreicht werden. Aus
  psychologischen Gründen kann es jedoch sinnvoll sein \Es{\ce{O2} in geringer
  Dosierung (1\LMin{})} über eine \ce{O2}-Maske zu
  verabreichen.
  \item Wenn nicht erfolgreich und Krämpfe bestehen bleiben bzw.
  Bewusstseinsstörungen auftreten: \TxMassVit .
\end{itemize}
}





\clearpage
\section{Neurologische Erkrankungen und Notfälle}

\Definition{Erkrankungen, Notfälle und sonstige Störungen
die das Nervensystem betreffen. Neurologische Störungen
gehen oft mit einer Störung des Bewusstseins oder Schmerzen einher.}

Eine Übersicht und willkürliche Einteilung:

\begin{description}
  \item [Schlaganfall]
  \item [Vergiftungen] können fast alles verursachen, inkl. neurologischer
  Symptome.
  \item [Schädel-Hirn-Trauma] Abkz. 'SHT'; Ein Schlag auf
  den Kopf erhöht \emph{nicht} das Denkvermögen \ldots
  \item [Stoffwechselstörungen] welche Einfluss auf das Nervensystem nehmen, z.\,B. Zuckerentgleisungen
  \item [Infektionen] können das Nervensystem beeinflussen, z.\,B. Verwirrtheit bei Fieber
  \item [Alter] beeinflusst auch das Nervensystem, im Rahmen einer Demenz kommt es z.\,B. zum Abbau der Denkleistung
  \item [Krampfanfälle] können vom zentralen Nervensystem ausgehen
  \item [Erkrankungen des Stützapparates] Aufgrund der räumlichen Nähe zum Stütz- und Bewegungsapparates kann das Nervensystem in Mitleidenschaft gezogen werden, z.\,B. bei Bandscheibenvorfällen
\end{description}

\minisec{Querverweise}

\begin{itemize}
  \item Infektiöse (bakterielle) Meningitis:
  \prettyref{topic:meningitis}
\end{itemize}


\subsection{Schlaganfall, Apoplektischer Insult, Hirninfarkt}
\index{Schlaganfall}
\index{Insult}
\label{topic:tia}

\emph{\Lang{Syn.} \index{Insult}\index{Apoplektischer
Insult}Apoplektischer Insult, \index{Apoplexie}Apoplexie}

\Definition{siehe Arten.}

\Layout{0}{Arten}
{%
Grundsätzlich gibt es 2 Formen, die aber vor Ort nicht
unterschieden werden können:
\begin{labeling}[~--]{xxxxxxxxxxxxxxxxxx}
  \item[\T{Trockener Insult}] Verschluss eines Hirngefäßes mit anschließendem
  Absterben der versorgten Hirnregion ("`trockener
  Insult"'), oder
  \item[\T{Blutiger Insult}] Hirnblutung, z.\,B. durch Platzen eines
  Hirngefäßes, kann spontan oder aufgrund eines Traumas auftreten.
\end{labeling}

}
{
\begin{itemize}
  \item "`Trockener"' Insult
  \item "`Blutiger"' Insult
\end{itemize}

}{}{}{}

\Layout{0}{\T{TIA}}
{\index{TIA}
\index{Transitorisch-Ischämische Attacke}
TIA ist die Abkürzung für Transitorisch-Ischämische Attacke. Sie ist eine Sonderform des
Schlaganfalls bei der die Symptome innerhalb von 24 Stunden vergehen und kein
Hirngewebe abstirbt.
    % TODO INHALT: TIA: Ischämie erklären (evtl. in EIN)
}
{
% TODO MED: Slide fehlt!
\begin{itemize}
  \item Transitorisch-Ischämische Attacke.
  \item Symptome <\,24\,h.
\end{itemize}

}{}{}{}

% \Dias
% {
%     \includegraphics[width=\linewidth]{./images/noauth/AASdev20090228-img83.png}
%     \Captionof{figure}{\AuthNotesPicLegalNotOk{img83}{}Insult: Eine spastische
%     Halbseitenlähmung kann sich als Spätfolge Wochen nach einem durchgemachten
%     Schlaganfall e ntwickeln.
%     \SourcePicture{Quelle nicht eruierbar.}}
% }{
%     \includegraphics[width=\linewidth]{./images/noauth/AASdev20090228-img84.png}
%     \Captionof{figure}{\AuthNotesPicLegalNotOk{img84}{}Insult:
%     Halbseitenlähmung im Gesicht.
%     \SourcePicture{Quelle nicht eruierbar.}}
% }{}


\Layout{0}{\TxSymptome}
{
Die Symptome sind davon abhängig welche Hirnregion betroffen ist und 
können sehr unterschiedlich sein und umfassen:

%\begin{multicols}{2}
\begin{labeling}[~--]{xxxxxxxxxxxxxxxxxx}
  \item[Motorik] Störungen der Kraft und Beweglichkeit sind sehr häufig:
  
  \T{Halbseitenzeichen}: Aufgrund des Ausfalls eines Zentrums im
  Gehirn kann es zu einer \E{Kraftminderung} einer Seite kommen, dies kann bis
  zu einer \E{schlaffen, kompletten Lähmung} führen.
  
  Ein \E{herabhängender Mundwinkel} ist ein Zeichen einer Lähmung
  der Gesichtsmuskulatur.
  
  Durch die Lähmungen kann es zu \Es{Schluckstörungen} kommen, es
  besteht dann eine erhöhte Aspirationsgefahr.
  
  Der \T{Herdblick} ist eine "`Blickdiagnose"': Es kann zu Blicklähmungen
  kommen, der Patient blickt in Richtung des "`Herd des Geschehens"'.
  
  \item[Sprache] Störungen der Sprache sind ebenfalls sehr
  häufig
  
  Oft nimmt man eine \Es{verwaschene Sprache} wahr, welche (fälschlich) an
  Trunkenheit erinnert.
  
  Weiters kann es zu einer "`wirren Sprache"' kommen: Worte fehlen dabei oder
  Silben werden vertauscht. Im Extremfall kann die Sprache gänzlich "`verloren"'
  gehen. Auf der anderen Seite kann die Sprache zwar erhalten, aber das
  Sprach\E{verständnis} geschädigt sein
  
  \item[Wahrnehmung]
  Es kann zu Wahrnehmungstörungen wie z.\,B. \Es{Sehstörungen}, plötzliche
  Erblindung, Doppelbilder, Gesichtsfeldausfälle kommen.
  
  \item[Bewusstsein] Störungen des Bewusstseins sind gefährlich und müssen bei
  Insult-Patienten immer erwartet werden. Der Bewusstseinszustand kann sich
  auch plötzlich verschlechtern.
  
  \item[Allgemein] Weiters können unspezifische Symptome wie z.\,B.
  Schwindel und Kopfschmerzen auftreten.
  Ein plötzlicher, \Es{donnerschlagartiger Kopfschmerz} ist ein klassisches
  Symptom für eine spontane Hirnblutung.
  
\end{labeling}
%\end{multicols}
}
{
Je nach betroffener Hirnregion sehr unterschiedlich.
\begin{itemize}
  \item Halbseitenzeichen (Schwäche, Lähmung)
  \item Herdblick
  \item Verwaschene Sprache
  \item Sehstörungen
  \item Bewusstseinsstörungen
  \item Kopfschmerzen
\end{itemize}
}{}{}{}


\Layout{0}{Differentialdiagnose}
{
Die wichtigste Differntialdiagnose ist eine Störung des
Blutzucker-Stoffwechsels. Sowohl die Hyper(${\uparrow}$)-, als auch die
Hypoglykämie (${\downarrow}$) können einen Schlaganfall vortäuschen. Bei
neurologischen Symptomen muß der Blutzucker immer gemessen werden!
Ebenso ist immer an Vergiftungen zu denken.

Die Abgrenzung zu anderen neurologischen Erkrankungen ist oft schwierig bzw.
kaum möglich.
}{
\begin{itemize}
  \item Blutzucker ${\downarrow}$, ${\uparrow}$.
  \item Sonstige neurologische Erkrankung.
  \item Vergiftungen
\end{itemize}

}{}{}{}


\MassnahmenNG{Spezielle Maßnahmen}{Insult}{}{

\begin{itemize}
  \item \TxMassVitBes
  \begin{itemize}
    \item Lagerung: \E{leicht erhöhter Oberkörper}
    \item \ce{O2} 10\LMin, bei Bedarf bis zu 15\LMin
    % ERRATUM 0.8.0 Nachberunfung des Notarztes beim Schlaganfall
    
    \item Notarzt: Grundsätzlich ist ein Schlaganfall-Patient aufgrund seiner Diagnose als vital bedroht einzustufen. Es gibt jedoch eine
    Besonderheit: Eine definitive Therapie in einer Spezialabteilung
    (Stroke Unit) kann andere, noch nicht abgestorbene Teile des Hirns
    vor weiterem Schaden bewahren, d.h. es ist sehr wichtig, dass diese
    Therapie möglichst schnell erfolgt.
    
    Daher soll von einer
    \Es{Nachberufung eines Notarztes \underline{abgesehen} werden,
    \E{wenn folgende Umstände gegeben sind}}:
    \begin{itemize}
      \item Die Verdachtsdiagnose gilt als \E{sehr sicher}.
      \item Der Patient ist sonst als \E{stabil} eingestuft
      (\prettyref{m:ersteinschaetzung}) und
      bewusstseinsklar\footnote{Bei der Einschätzung des Bewusstseins
      darf man sich nicht durch die durch den Schlaganfall
      verursachten Symptome (Sprachstörungen, Blickstarre \ldots)
      beirren lassen.}.
      \item Es gibt sonst \E{keine anderen Umstände} die die Berufung
      eines Notarztes erforderlich machen.
      \item Es muss versucht werden, eine Aufnahme auf eine
      \E{Spezialabteilung} für Schlaganfälle zu
      organisieren\footnote{Der Erfolg ist abhängig von der
      Verfügbarkeit eines entsprechenden freien Bettes.} (Stroke Unit,
      Kontaktaufnahme mit der Leitstelle).
    \end{itemize}
    \A{HART: Muss es unbedingt eine Spezialabteilung sein?
    "`bewusstseinsklar"' könnte missverstanden werden.}
    \A{\par Diskussion: }
    \ACh{KOCH}{2010-02-13}{ok.}
    \ACh{HART}{2010-02-13}{Muss es unbedingt eine Spezialabteilung sein?
    "`bewusstseinsklar"' könnte missverstanden werden.}
    \ACh{GABS}{2010-02-13}{Fußnote wegen Bewusstseinseinschätzung
    eingefügt. Spezialabteilung: muss nur \underline{versucht} werden
    eine Stroke Unit zu organisieren, je nach Verfügbarkeit.}
    \ACh{GABS/STEU}{2010-04-06}{Ticket 38:
    Besprechung gestern mit STEU: kann so übernommen werden, sieht
    allerdings Problem wegen Inkonsistenz und für den Nutzer
    möglicherweise verwirrender wenn-dann-Verzweigung. Bei Hirnblutung
    ist das Vorgehen ein Nachteil. (Einwand GABS: Ischämisch häufiger,
    Abwägung Benefit Narkose, etc. vs. schnelles CT)}
  \end{itemize}
  \item \Es{Neurocheck} inkl. Messung des \Es{Blutzuckers}!
  \item Hospitalisation: Wenn verfügbar: \Es{Stroke Unit}, sonst
  Abklärung mit der Leitstelle wegen Alternativen (Abt. für
  Neurologie oder Innere Medizin, mit Möglichkeit zur
  Intensivüberwachung).
\end{itemize}
}

\subsection{Vom Gehirn ausgehende Krämpfe --- Zerebrale Krampfanfälle}
\label{topic:zerebrale-krampfanfaelle}
\label{topic:krampfanfall}

\Definition{Vom Gehirn ausgehende motorische Krampfanfälle aufgrund von akuten
oder chronischen Erkrankungen des Gehirns. Dabei entladen sich unkontrolliert
Nervenzell-Gruppen im Gehirn.}


\Layout{0}{Einteilung}
{%
Die zerebralen Krampfanfälle kann man auf verschiedene Arten einteilen

\begin{enumerate}
  \item Nach Art:
  \begin{labeling}[~--]{Generalisierter Anfall }
    \item[Tonisch] "`Verkrampfung"', erhöhter
    Muskeltonus
    \item[Klonisch] Zuckungen
    \item[Tonisch-klonisch] Kombination aus beidem, das
    klassische Erscheinungsbild eines Krampfanfalles
  \end{labeling}
  \item Nach Lokalisation:
  
  \begin{labeling}[~--]{Generalisierter Anfall }
    \item [Fokaler Anfall] Der Krampfanfall ist auf
    eine Körperregion beschränkt, siehe
    \ref{topic:krampfanfall-fokaler}
    
    \item [Generalisierter Anfall] Er betrifft den
    ganzen Körper, siehe
    \ref{topic:krampfanfall-generalisierter}
    
  \end{labeling}
\end{enumerate}
}{%
Die zerebralen Krampfanfälle kann man auf verschiedene Arten einteilen

\begin{itemize}
  \item Nach Art:
  \begin{description}
    \item Tonisch
    \item Klonisch
    \item Tonisch-klonisch
  \end{description}
  \item Nach Lokalisation:
  \begin{description}
    \item Fokaler Anfall
    \item Generalisierter Anfall
  \end{description}
\end{itemize}
}{}{}{}

\Layout{2}{Ursachen}
{%
Für zerebrale Anfälle gibt es eine Reihe von unterschiedlichen Ursachen die
direkt oder indirekt das Gehirn betreffen. Einige betreffen die Hirn- und
Nervenstruktur, andere sorgen dafür, dass das Hirn zu wenig
Sauerstoff und andere Nährstoffe erhält. Die häufigsten Ursachen und
Grunderkrankungen sind im Folgenden aufgeführt: 

\begin{labeling}{xxxxxxxxxxxxxxx}
  \item[Epilepsie] Die Epilepsie ist eine chronische Erkrankung welche durch wiederkehrende
  Episoden von zerebralen Krampfanfällen gekennzeichnet ist.
  \item[Vergiftungen] (Medikamente/Drogen/...) können wieder fast alles machen,
  also auch Krampfanfälle. \cite{Supady2009}
  \item[Hirntumor] Durch den Tumor wird das Hirn beschädigt, auch dadurch kann es
  zu Anfällen kommen.
  \item[Narbengewebe] nach Insult oder Hirnverletzung.
  \item[Hirnblutung]
  \item[Sauerstoffmangel] z.\,B. Kollaps, Ischämie aufgrund eines Gefäßverschlusses
  \item[BZ] Der Dauerbrenner. Eine Blutzuckerentgleisung kann fast alle
  neurologischen Symptome vortäuschen!
  \item[angeboren] ohne organisch fassbare  Ursache
  \item[Fieberkrampf] Fieber bei Kleinkindern
  \item[Infektionen]
  \item[Entzug] Entzug von Alkohol oder anderen Drogen
\end{labeling}
}{
}{}{}{}






\subsubsection{Fokaler Anfall}
\index{Krampfanfall!fokaler}
\label{topic:krampfanfall-fokaler}
Der Krampfanfall ist auf eine Körperregion beschränkt, z.\,B.
auf eine Hand. Der Patient kann bei Bewusstsein sein. Ein
fokaler Anfall kann sich ausweiten und in einen
generalisierten Anfall übergehen. 

\Z{Der fokale Anfall wird
auch \emph{"`petit mal"'}\ZFootnote{Petit mal: \emph{franz.:} "`kleines
Unglück"'} genannt.}
      


\subsubsection{Generalisierter Anfall}
\index{Krampfanfall!generalisierter}
\label{topic:krampfanfall-generalisierter}

\Layout{0}{Beginn}
{%
Er betrifft den ganzen Körper. Dem Anfall kann ein
\E{Dämmerzustand} 
vorausgehen, \Z{eine sog. "`Absence"'}. Der Patient wirkt
dabei plötzlich abwesend und \emph{"`irgendwie komisch"'}.
Dann verliert er das Bewusstsein und die Muskeln werden
anfangs schlaff. Wenn der Patient gestanden ist, fällt er
ungebremst auf den Boden (Cave: \Es{Verletzungen!}).
Manchmal kommt es vor dass der Patient einen
Initialschrei\index{Initialschrei} von sich gibt, der
gruselig klingen kann. 
}{
\begin{itemize}
  \item Dämmerzustand -- Absence
  \item Initialschrei
  \item evtl. Sturz
  \item evtl. Folgeverletzungen
\end{itemize}

}{}{}{}

\Layout{0}{Krampf}
{%
Dann kommt es zu
\Es{tonisch-klonischen Krämpfen},
die den ganzen Körper betreffen. Während des Anfalls ist der
Patient nicht bei Bewusstsein, es besteht außerdem i.\,d.\,R.
ein Atemstillstand. Er kann jedoch starke Kräfte entwickeln, es
besteht massive
Verletzungsgefahr.
}{
\begin{itemize}
  \item Tonisch-klonische Krämpfe
  \item Verletzungsgefahr
\end{itemize}
}{}{}{}


\Layout{0}{Danach}
{%
Der Krampf dauert meist wenige Minuten und hört von selbst
auf. Der Patient ist danach erstmal bewusstlos, im weiteren
Verlauf mit mehr oder weniger großem Aufwand erweckbar. Er
befindet sich dann in der sogenannten
\T{Nachschlafphase}\index{Nachschlafphase}.

Der Patient \E{klart im Verlauf mehr und mehr auf}. Es ist
wichtig zu versuchen mit dem Patienten in Kontakt zu treten
um seinen Wachheitsgrad und dessen
\Es{Verlauf} beurteilen zu können. Für den Anfall selbst besteht
Amnesie! Daher muss man den Patienten aufklären was passiert
ist, wo er sich befindet, etc. Ein \Es{Neuro-} und
\Es{Traumacheck} ist nach \underline{\E{jedem}} Anfall
unbedingt notwendig! 
}{
\begin{itemize}
  \item Nachschlafphase
  \item Aufklaren
  \item Amnesie
  \item \Es{Neurocheck}
  \item \Es{Traumacheck} nach \underline{\E{jedem}} Krampfanfall
\end{itemize}
}{}{}{}


\Layout{0}{Sonstige Besonderheiten}
{%
Häufig kommt es während des Anfalles zu einem Zungenbiss,
außerdem zum Einnässen oder Einkoten.
}{
\begin{itemize}
  \item Zungenbiss
  \item Einnässen
  \item Einkoten
\end{itemize}
}{}{}{}


\Z{Der generalisierte Anfall wird auch \emph{"`grand
mal"'}\ZFootnote{Grand mal: \emph{franz.:} "`großes Unglück"'} genannt.}

\subsubsection{Besondere Krampfanfälle}

\Layout{60}{Fieberkrampf}
{%
\label{topic:fieberkrampf}
\index{Fieberkrampf}
\index{Krampfanfall!Fieberkrampf}
%
\Es{Kinder} unter 5~Jahren können bei schnellem, hohen Fieberanstieg
({\textgreater}\,39{\textcelsius}) sogenannte \E{Fieberkrämpfe}
entwickeln. Sie entsprechen einem vom Gehirn ausgehenden Krampfanfall und sind grundsätzlich
genau so zu behandeln. 

Mitunter ist die Neigung des Kindes zu Fieberkrämpfen schon bekannt und es
wurden bereits spezielle Zäpfchen für den Bedarfsfall verschrieben (und bei
Eintreffen eventuell schon von den Eltern gegeben). 
}{
\begin{itemize}
  \item Kinder {\textless} 4 Jahre.
  \item Hohes Fieber.
  \item Wie generalisierter Anfall.
\end{itemize}
}
{./images/gabriel-sebastian-aass/fieberthermometer.jpg}
{}
{GABS.}



\Layout{0}{Eklamptischer Krampfanfall während der Schwangerschaft}
{%
In Folge einer Präeklampsie (EPH-Gestose)
kann es zu einem sog. \E{eklamptischen Krampfanfall}
kommen. Dieser wird unter \prettyref{topic:eklampsie}
besprochen.
}{
\begin{itemize}
    \item Folge einer Präemklampsie.
    \item \prettyref{topic:eklampsie}.
\end{itemize}
}{}{}{}


\subsubsection{Status Epilepticus}
\index{Status!epilepticus}
\Layout{0}{Lebensgefahr}
{
Der Status Epilepticus ist eine
lebensgefährliche Steigerung
des generalisierten Krampfanfalls. Er besteht dann, wenn der
Anfall nicht nach kurzer Zeit vergeht oder innerhalb kurzer
Abstände mehrere Anfälle auftreten (z.\,B. 2 Anfälle
innerhalb von 24 Stunden). Der Krampf muss dann mit Medikamenten
durchbrochen werden.
}
{
\begin{itemize}
  \item Lange Dauer.
  \item Wiederholtes Auftreten (>\,2\,/\,24\,h).
  \item Akute \Ca{Lebensgefahr!}
\end{itemize}
}{}{}{}


\Fazit{Im Rettungsdienst gilt praktisch jeder Krampfanfall,
den das Personal selber miterlebt, als Status-Verdächtig.}

\subsubsection{Maßnahmen}

\MassnahmenNG{Spezielle Maßnahmen}{Krampfender Patient}{}{
\begin{itemize}
  \item \TxMassVitBes
  \begin{itemize}
    \item Lagerung: sichern
    \begin{itemize}
      \item Patient vor Verletzungen schützen
      \item Möbel etc. entfernen oder polstern
      \item NICHT festhalten -- aber sichern, Eigenschutz
      beachten!\footnote{\textbf{Sichern}: Patient darf nicht von der
      Trage oder aus dem Bett fallen, der Kopf darf nicht am Boden oder Tischbein anschlagen, etc.}
    \end{itemize}
  \end{itemize}
  \item nach Anfall BAK-Schema, Atemwege freihalten
  \item stabile Seitenlage  (Erholungsschlaf wie  Bewusstlosigkeit)
  \item Besondere Diagnostik:
  \begin{itemize}
    \item Neurocheck inkl. Messung des \Es{Blutzuckers}!
    \item Traumacheck
  \end{itemize}
  \item Reizabgeschirmter Transport
  \item Hospitalisierung: Abt.f. Innere Medizin. (evtl. auch
  Intensivüberwachung)
\end{itemize}
}

\MassnahmenNG{Spezielle Maßnahmen}{Nicht-(mehr) krampfender Patient}{}{
\begin{itemize}
  \item Umstände klären und entscheiden ob \Es{Notarzt}
  notwendig ist:
  \begin{itemize}
    \item 2 oder mehr Krampfanfälle innerhalb von 24 Stunden
    \item Anfälle "`erheblich anders"' als sonst üblich
    \item Erster Krampfanfall überhaupt
  \end{itemize}
  \item Besondere Diagnostik:
  \begin{itemize}
    \item Neurocheck inkl. Messung des \Es{Blutzuckers}!
    \item Traumacheck
  \end{itemize}
  \item Reizabgeschirmter Transport
  \item Hospitalisierung: Abt.f. Innere Medizin.
\end{itemize}
}



\subsection{Lumbago, Lumbo-ischialgie und Bandscheibenvorfall} 
\Z{(Discusprolaps)} 
\label{topic:lumbago}

Unter Ischialgie (\T{Lumbalgie}, Lumboischialgie,
\emph{"`\T{Lumbago}"'}) versteht man Schmerzen im
Versorgungsbereich des Ischias-Nerv, welche u. a. durch eine Entzündung oder durch Reizung bzw. Kompression des Nervs oder seiner
Wurzeln verursacht werden können. 

\paragraph{Ursachen} können eine Reizung bzw. Druck
im Bereich der Lendenwirbel bzw. der Kreuzbeinwirbel,
ein Bandscheibenvorfall, sowie andere Erkrankungen der Wirbelsäule, aber auch ein Unfall sein. 

\paragraph{\TxSymptome} sind Schmerzen in der
Lendengegend, die in das betroffene Bein bis zum Fußaußenrand ausstrahlen, 
Schonhaltung, lokale Druck- und Klopfempfindlichkeit der Wirbelsäule,
Sensibilitätsstörungen und eventuell Lähmung der Zehenmuskulatur. Je nach
Schwere und zeitlicher Entwicklung (langsames/plötzliches Auftreten) bezeichnet
man das Zustandsbild als Lumboischialgie oder Lumbago (Hexenschuss). 
Gegebenenfalls ist im Krankenhaus ein Bandscheibenvorfall abzuklären.




\clearpage
\section{Metabolische Störungen (Stoffwechselstörungen)}

Die relevanteste Stoffwechselstörung im Rettungsdienst ist der
\index{Diabetes mellitus}\T{Diabetes mellitus}. Dieser
stellt eine Störung des Kohlehydrat-(zucker-)stoffwechsels dar.
\Es{Zucker ist ein Kohlehydrat}. 

\subsection{Diabetes Mellitus}
\label{topic:diabetes}
\label{topic:diabetes-mellitus}

\Definition{Diabetes mellitus ist eine chronische
Stoffwechselerkrankungen bei der die Blutzuckerregulation gestört ist.} 



\Layout{2}{Exkurs: Der Blutzuckerwert}
{%
%
%\vspace{1em}
%
\begin{center}
\begin{tabular}{cccl}
\toprule
\\\midrule
0 & -- & 20 & Lebensgefahr
\\\addlinespace
21 & -- & 60 & stark erniedrigt, Bewusstseinsstörungen
\\\addlinespace
61 & -- & 80 & erniedrigt, leichte Symptome
\\\addlinespace
81 & -- & 120 & Nüchtern-Blutzucker
\\\addlinespace
121 & -- & 200 & erhöht
\\\addlinespace
201 & -- & 399 & stark erhöht
\\\addlinespace
400 & -- & ?? & sehr stark erhöht
\\\addlinespace\bottomrule
\end{tabular}
\end{center}

Durch die Erkrankung oder die unsachgemäße Behandlung kann es zu 
folgenden Notfällen kommen:

\begin{description}
\item[\T{Hyp\underline{\E{o}}glykämie}] \ldots "`Unterzucker"'

\item[\T{Hyp\underline{\E{er}}glykämie}] \ldots "`Überzucker"'

\end{description}
\TODO{GABS}{2010-04-18}{Bild: }{Interpretation des
BZ-Wertes: Richtwerte.}
}
{}
% {./images/noauth/AASdev20090228-img88.png}
{./icons/under-construction.png}
{Interpretation des
BZ-Wertes: Richtwerte.}
{}

\Layout{2}{Insulin}
{%
\T{Insulin} \index{Insulin} ist ein Hormon und wird in den \T{Beta-Zellen}
der Bauchspeicheldrüse (Pankreas) produziert. Es ist für
den Zuckerstoffwechsel wichtig. Es sorgt dafür, dass Zucker
(\T{Glukose}) aus dem Blut in die Zellen aufgenommen wird.

\Fazit{Insulin senkt den Blutzuckerspiegel.} 


Wenn man künstlich Insulin spritzt, besteht die Gefahr eines
Unterzuckers (zu viel Insulin gespritzt oder zu wenig
gegessen). ${\rightarrow}$  Spritz/Essabstand
er\-fra\-gen.
}{

}{}{}{}




Zwei Störungen sind möglich:

\begin{enumerate}
  \item Es wird zu wenig Insulin gebildet.
  \item Der Körper hat einen gesteigerten Bedarf an Insulin weil er schon
  zu sehr daran gewöhnt ist. Der Körper reagiert auf die 'normale' Dosis
  nicht mehr so wie er sollte, er braucht größere Mengen Insulin um
  den gleichen Effekt wie ein Gesunder zu erreichen (Gewöhnungseffekt),
  irgendwann kann aber der Körper nicht mehr Insulin produzieren.
\end{enumerate}

Folglich unterscheidet man 2 Typen:



\Layout{0}{Diabetes mellitus: Typen}
{
\begin{labeling}[~--]{xxxxxxx }
  \item[Typ I]  \Z{(IDDM,
  insuline dependent DM)} Die Insulinproduktion ist gestört und
  der Körper produziert zu wenig Insulin. Der Patient
  muss als Ersatz künstlich Insulin spritzen.  Die Patienten sind eher jung
  (Jugendliche, junge Erwachsene) und sportlich.
  
  \Z{Ursache für die Störung ist eine Störung des
  Immunsystems, bei der es zu einer Zerstörung der Insulin-produzierenden
  ${\beta}$-Zellen des Pankreas kommt. Durch den Produktionsausfall kommt es
  zu einem absoluten Mangel.}
  
  \item[Typ II]
  \Z{(NIDDM, non insuline dependent DM)} Durch eine dauerhaft zu hohe
  Zuckerzufuhr kommt es im Körper zu Gewöhnungseffekten und einem
  Wirkungsverlust des vom Körper produzierten Insulins
  \Z{(Insulinresistenz)}. Als Folge muss der Körper immer mehr Insulin produzieren um den gleichen Effekt zu
  erreichen. Wenn das Produktionslimit erreicht ist kommt es zu
  einer Erhöhung des Blutzuckerspiegels und zum Typ-II-Diabetes.
  
  Der Patient mit Typ-II-Diabetes benötigt oft kein Insulin. Die
  Erkrankung kann zuerst mittels Lebensstiländerungen und
  \Es{Medikamenten} behandeln. Erst im Endstadium muss mitunter auch Insulin
  gespritzt werden.
  
  Der Typ-2-Diabetes betrifft eher Patienten im mittleren Alter und aufwärts,
  aufgrund der zunehmend fett- und kohlehydratreichen Ernährungsgewohnheiten
  werden die Patienten aber immer jünger.
\end{labeling}
}{
\begin{itemize}
  \item Typ I
  \begin{itemize}
    \item Insulinproduktion gestört.
    \item Insulin muss gespritzt werden.
    \item Eher Jugendliche und junge Erwachsene.
  \end{itemize}
  \item Typ II
  \begin{itemize}
    \item Lebensstil, Ernährung.
    \item Insulin-Gewöhnungseffekt.
    \item Eher ab mittlerem Alter.
    \item Lebensstiländerung, Medikamente, Insulin.
  \end{itemize}
\end{itemize}
}{}{}{}

\Layout{0}{Spätfolgen}
{
Ein hoher Blutzuckerspiegel über Jahre schädigt den gesamten
Körper. Es kommt zu:
\begin{labeling}[~--]{xxxxxxxxxxxxxxx }
  \item[Gefäßschäden] durch Verkalkung
  \begin{itemize}
    \item Gefäßverschlüsse (kleine und große Gefäße)
    \item \Es{Beine}: "`Diabetischer
    Fuß"', in Folge der Mangeldurchblutung in der Peripherie heilen auch kleine
    Wunden nicht mehr, in weiterer Folge muss immer weiter amputiert
    werden.
    \item \Es{Herz}: Die Herzkranzgefäße
    (Koronar\-arterien) sind ebenso betroffen. Diabetiker sind Risikopatienten für
    Herz-Kreislauf-Erkrankungen wie zB Herz\-infarkt.
    \item \Es{Hirn}: Wie beim Herzen kommt es zu
    Gefäßschäden, daher ist das Risiko für
    Schlaganfälle erhöht.
    \item \Es{Niere}: Die sehr kleinen Gefäße der
    Niere werden geschädigt. Im weiteren Verlauf stellt die Niere ihre
    Arbeit ein, es kommt zu einer "`\T{chronischen
    Niereninsuffizienz}"', die im Endstadium mittels
    Blutwäsche (Dialyse) behandelt werden muss. Dialyse-Patienten machen
    einen großen Teil der Krankentransporte aus.
  \end{itemize}
  \item [Nervenschäden] \Z{(periphere Polyneuropathie)}:
  Gefühlslosigkeit vor allem in den Extremitäten, aber auch am
  gesamten Körper: Ein Herzinfarkt kann zum Beispiel schmerzlos
  ablaufen!
  \item [Augenschäden] bis hin zur Erblindung
  
\end{labeling}

Krankheiten, die sonst schmerzhaft sind, können beim Diabetiker aufgrund
der Nervenschäden ohne Schmerzen ablaufen. 

\Cave{Achtung: "`Stummer"', schmerzarmer/-loser Herzinfarkt beim Diabetiker!}
}{
\begin{itemize}
  \item Gefäßschäden (Herz, Hirn, Nieren, Extremitäten, \ldots)
  \item Nervenschäden
  \item \Ca{Cave "`Stummer Infarkt"'!}
  \item Augenschäden
\end{itemize}
}{}{}{}








\subsubsection{Hypoglykämie}
\label{topic:hypoglykaemie}
\index{Hypoglykämie}

\Definition{Maßgebliche Erniedrigung des Blutzuckerspiegels.}

\Layout{0}{Ursachen}
{
Die häufigsten Ursachen für eine plötzliche Hypoglykämie sind Diätfehler
und/oder eine falsche Anwendung von Insulin oder bestimmten anderen
Diabetes-Medikamenten. Aufgrund besonderer Anstrengungen kann es zu einem
gesteigerten Zuckerbedarf des Körper kommen, wodurch der Blutzuckerspiegel
sinkt. Dies kommt häufig im Rahmen von Begleiterkrankungen, wie z.\,B. grippalen
Infekten, vor. Auch die Alkohol-Intoxikation kann zu einer
Hypoglykämie führen\opt{NIVDOC}{ (Das beim Ethanolabbau
entstehende \Abk{NADH} hemmt die
Glukoneogenese.)}\cite{LoefflerBiochemie7}. 
}{
\begin{itemize}
  \item zu viel Insulin (spritzen ohne zu essen).
  \item Fasten (spritzen ohne zu Essen).
  \item Besondere Anstrengungen, Begleiterkrankungen,
  insb. Infekte.
  \item Alkohol-Intoxikation.
\end{itemize}
}{}{}{}






\Layout{0}{\TxSymptome}
{
Der Patient mit Hypoglykämie ist meistens unruhig, desorientiert und agitiert,
oft auch unkooperativ und aggressiv. Oft fällt ein Zittern und eine schweißige
Haut auf. Auch Kopfschmerzen sind häufig. Eine Tachykardie und Tachypnoe ist
häufig feststellbar. Manchmal klagt der Patient über Übelkeit und Heißhunger.

Je tiefer der Blutzuckerspiegel sinkt, desto deutlicher werden die
\Es{Bewusstseinsstörungen}. Von der anfänglichen
Agitiertheit/Desorientiertheit kann es rasch zu einer Somnolenz, zum Sopor und
zum Koma kommen!

Eine Hypoglykämie kann fast alle neurologischen  Störungen
aus\-lösen/imitieren! Häufig sind Kopfschmerzen und Bewusstseinsstörungen, evtl.
auch Krampfanfälle. Bei jeder neurologischen Symptomatik muss man daher an eine Hypoglykämie denken!

\Cave{Bei der Hypoglykämie können Symptome auftreten, die
den Patient wie alkoholisiert wirken lassen! \E{Hüte dich vor vorschnellen
Diagnosen!}}

\Fazit{${\rightarrow}$ BZ-Messung bei jeder neurologischen Störung. }

}{
\begin{itemize}
  \item Unruhe, Zittern, Schweißausbruch
  \item Aggressives Verhalten, unkooperativ
  \item Neurologische Symptomatik.
  \begin{itemize}
    \item Bewusstseinsstörungen, Verwirrtheit, Kopfschmerzen, Krampfanfälle,
    Bewusstlosigkeit.
  \end{itemize}
  \item Heißhunger, Übelkeit.
  \item Tachykardie, Tachypnoe.
  \item Haut feucht/blass.
\end{itemize}
}{}{}{}












\MassnahmenNG{Spezielle Maßnahmen}{Hypoglykämie}{}{
\begin{itemize}
  \item Vitale Bedrohung einschätzen. \TxBeiVit \TxMassVit
  
  \item BZ messen! Zur Diagnosestellung bzw. \E{vor} und \E{nach}
  jeder Zuckergabe
  \item Zucker zuführen: Wenn der Patient
  bewusstseinsklar ist Traubenzucker oder Zuckerwasser geben. Danach "`feste
  Speisen"' nachessen lassen, da der Traubenzucker nur eine kurze Wirkdauer hat.
  
  Wenn der Patient nicht mehr bewusstseinsklar ist kann die Zuckergabe
  nur über eine venöse Gabe erfolgen (getrübte Patienten dürfen aufgrund
  der Aspirationsgefahr nichts mehr essen).
  
\end{itemize}
}

\subsubsection{Hyperglykämie}
\label{topic:hyperglykaemie}
\index{Hyperglykämie}

% Done: XXX INHALT-NEU: Hyperglykämie, Text und Überarbeitung,
% POI 11
\A{\#11: Muss noch überarbeitet werden}

\Definition{Maßgebliche Erhöhung des Blutzuckerspiegels.}

\Layout{0}{Hyperglykämie: Zustand oder Notfall?}
{
Eine Erhöhung des Blutzuckerspiegels ist nicht automatisch ein Notfall. Die
Bedrohlichkeit ist z. B. stark davon abhängig wie
rasch sich diese Erhöhung entwickelt hat. So hat ein
Patient, bei dem seit Jahren ein unerkannter Diabetes mellitus besteht auch mit sehr hohen Blutzuckerwerten (akut) keine
Probleme. Der Körper hat sich über die Zeit daran \E{gewöhnt}. 

Ein anderer
Patient, bei dem der Zustand plötzlich eingetreten ist (z.\,B.
Erstmanifestation eines Diabetes mellitus Typ I bei einem 19-jährigen, oder wenn
das Insulin nicht gespritzt wurde, \ldots), kann mit dem gleichen
Blutzuckerwert bereits im Koma liegen. 

\Fazit{Beurteile nicht nur den Messwert, sondern immer auch den Zustand des
Patienten.} }
{←}
{}{}{}



\Layout{0}{Auswirkungen}
{
Im Grunde bereitet der Zucker im Akutfall dem Körper zwei
Probleme. Einerseits wirkt der Zucker wie ein Elektrolyt auf den \Es{Flüssigkeitshaushalt} ein und
stört diesen. Es kommt zu einer ungünstigen Wasserumverteilung. Andererseits
mangelt es den Körperzellen an Zucker wenn das Insulin fehlt um es in die
Zellen hineinzubekommen, es entsteht ein \Es{Nährstoffmangel}.
}{
\begin{itemize}
  \item Störung des Wasser- und Elektrolythaushaltes.
  \begin{itemize}
    \item Zucker wirkt wie ein Elektrolyt.
  \end{itemize}
  \item Störung der Nährstoffversorgung.
\end{itemize}
}{}{}{}

\Layout{0}{Ursachen}
{
Die Hyperglykämie betrifft praktisch nur Diabetes mellitus-Patienten. Bei
Patienten, die einen noch unerkannten DM~Typ~I haben, ist ein hyperglykämischer
Notfall oft die Erstmanifestation (die Diagnose Diabetes ist noch nicht
bekannt!). Sonst sind Diätfehler oder eine mangelnde Insulinzufuhr häufige
Ursachen. Auch Stress und Begleiterkrankungen wie z.\,B. akute Infektionen
kommen in Betracht, da hier der Insulinbedarf steigt.
}{
\begin{itemize}
  \item Erstmanifestation eines DM~Typ~I
  \item Diätfehler
  \item Insulinzufuhr unzureichend/vergessen
  \item Stress
  \item Infektionen (Insulinbedarf ${\uparrow}$)
\end{itemize}
}{}{}{}









\Layout{0}{\TxSymptome}
{
Der Patient hat i.\,d.\,R. ein erhöhtes Durstgefühl und gleichzeitig eine
\Es{vermehrte Harnproduktion} (Zucker wird ausgeschieden). Er klagt über
Abgeschlagenheit, in weiterer Folge kann es zu \Es{Bewusstseinsstörungen} bis
hin zum \Es{Koma} (s.u.) kommen. Die Geschwindigkeit der Veschlechterung ist je
nach Patient sehr unterschiedlich.

\subparagraph{Koma} \label{topic:koma-diabetisches} \Z{Wie oben bereits
beschrieben wurde, kann es sowohl zu einem Problem des
Flüssigkeitshaushaltes und/oder der Nährstoffversorung der
Zellen kommen. Je nachdem welches Problem im Vordergrund steht, kommt es zu} \E{unterschiedlichen Symptomen}. 

Ist vorwiegend der Flüssigkeitshaushalt gestört kommt es zum sog.
\Z{Hyperosmolaren Koma}, dessen Leitsymptom erwartungsgemäß die
Bewusstseinsstörung ist.

Steht der Nährstoffmangel im Vordergrund passiert etwas Interessantes: Der
Körper produziert sog. \Z{Ketonkörer}. Da diese \E{sauer} sind kommt es zu
einer stoffwechselbedingten Übersäuerung (metabolische \Es{Azidose}
\prettyref{topic:azidose-metabolische}).

Unter \prettyref{topic:alkalose-respiratorische} wurde besprochen dass der pH-Wert
bei einer Hyperventilation in Richtung basisch wandert. Bei der diabetischen
Ketoazidose (sauer) versucht der Körper diesen Mechanismus auszunutzen: Er \Es{hyperventiliert} um sich "`vom sauren pH
    zum normalen pH zu atmen"'. Es ist daher bei diesen Patienten eine
    \Es{schnelle, tiefe Atmung} feststellbar\ZFootnote{Diese
    charakteristische Form der Atmung wird auch "`Kußmaul-Atmung"' genannt}.

Mitunter kann auch ein blumig/fruchtiger oder Azeton-artiger Atemgeruch
wahrgenommen werden. Dieser darf nicht mit einer "`Alkoholfahne"' verwechselt
werden! 
}{
\begin{itemize}
  \item Durstgefühl
  \item Vermehrte Harnproduktion
  \item Abgeschlagenheit, Bewusstseinstrübung bis
  \item Bewusstlosigkeit (Koma)
  \begin{itemize}
    \item Problem Flüssigkeitshaushalt
    \begin{itemize}
      \item \Z{Hyperosmolares Koma}
    \end{itemize}
    \item Problem Nährstoffmangel
    \begin{itemize}
      \item \Z{Ketoazidotisches Koma}
      \item Azidose → Hyperventilation zur Kompensation
    \end{itemize}
  \end{itemize}
  \item Azeton-Geruch (Nicht mit Alkoholfahne verwechseln!)
\end{itemize}
}{}{}{}



\MassnahmenNG{Spezielle Maßnahmen}{Hyperglykämisches Koma}{}{
\begin{itemize}
  \item \TxMassVitBes
  \begin{itemize}
    \item Lagerung: Stabile Seitenlage
    \item \ce{O2}: ab 10\LMin
  \end{itemize}
\end{itemize}
}





\clearpage
%%%%%%%%%%%%%%%%%%%%%%%%%%%%%%%%%%%%%%%%%%%%%%%%%%%%%%%%%%%
%%% Infektionskrankheiten                               %%%
%%%%%%%%%%%%%%%%%%%%%%%%%%%%%%%%%%%%%%%%%%%%%%%%%%%%%%%%%%%

\section{Infektions- und Entzündungskrankheiten}


%%%%%%%%%%%%%%%%%%%%%%%%%%%%%%%%%%%%%%%%%%%%%%%%%%%%%%%%%%%
%%%%%%%%%%%%%%%%%%%%%%%%%%%%%%%%%%%%%%%%%%%%%%%%%%%%%%%%%%%
%%%%%%%%%%%%%%%%%%%%%%%%%%%%%%%%%%%%%%%%%%%%%%%%%%%%%%%%%%%
\subsection{Grippe -- Influenza}
\index{Grippe}\index{Influenza}\index{Influenza}
\label{topic:grippe}

\Definition{Eine schwere, durch Viren verursachte
Infektionskrankheit.}

\Layout{2}{\TxBeschreibung}
{
Die Grippe ist eine \E{virale} Erkrankung, welche je nach
Patient und Virenstamm (s.u.) einen sehr schweren
Verlauf nehmen und auch tödlich enden kann. Gefürchtet sind
\Es{Komplikationen} wie Lungen- oder
Herzmuskelentzündungen oder weitere (Folge-)Infektionen.
Besonders gefährdet sind ältere und immunsupprimierte Personen.
}{}{}{}{}


\Fazit{Da es sich bei der Grippe um eine durch Viren
verursachte Krankheit handelt sind \Es{Antibiotika
wirkungslos}.}

Die Grippe tritt besonders in der kalten Jahreszeit auf und tritt
regelmäßig als \index{Epidemie}\Es{Epidemie} auf. Es
kann auch zum Ausbruch einer weltweiten \index{Pandemie}Pandemie kommen, \Z{z.\,B.
Spanische Grippe (1918, 25~Mio~Tote), Asiatische Grippe
(1957, 1~Mio~Tote) und  Hongkong-Grippe (1968,
800.000~Tote)\footnote{http://de.wikipedia.org/wiki/Influenza}}.
\marginlabel{\cite{wikiSpanischeGrippe}}

In Österreich erkranken pro Jahr ca. 380.000 Personen an der
Influenza, ca. 4.500 Personen müssen stationär behandelt werden und
ca. 120 Personen sterben an den
Folgen\cite{Strauss0000}.

%\footnote{http://www.bmg.gv.at/cms/site/standard.html?channel=CH0742\&doc=CMS1075899626901}.

\Fazit{Die Grippe tritt in Wellen auf und kann tödlich
enden.}

\Layout{60}{\TxSymptome}
{%
Typisch für die Grippe ist das \Es{rasche, hohe
Anfiebern (38{\textdegree}-40{\textdegree}C) mit
Schüttelfrost,  Kopf-, Hals-, Muskel-, Gelenksschmerzen, Husten und stark reduziertem
Allgemeinzustand} sowie lang dauernde Abgeschlagenheit nach
Ende der Erkrankung. Bei schon vorher geschwächten
Patienten kann der Verlauf weniger spektakulär aussehen (z.\,B. kaum Fieber), das Risiko einer Komplikationen ist allerdings natürlich höher. 
}
{
% TODO VIT: Slide fehlt!
}
{./images/gabriel-sebastian-aass/fieberthermometer.jpg}
{}
{GABS.}


\Layout{2}{Stämme}
{
Grippeviren kommen in verschiedenen
Stämmen (\Code{A}, \Code{B} und \Code{C}) vor, welche
wiederum durch unterschiedliche Arten von
\Z{Hämagglutinin} (\Code{H}) und \Z{Neuramidase}
(\Code{N}) unterschieden werden. Dadurch kommen so
klingende Namen wie "`\Code{A(H5N1)}"' oder
"`\Code{A(H1N1)}"' zu Stande. Je nach Unterart verhalten sich die Viren anders und man beobachtet Unterschiede hinsichtlich der
Ansteckungsgefahr, Schwere der Erkrankung usw. 
\Z{}
}
{

}
{}
{}
{}

\Layout{2}{Schutzimpfung}
{
Gegen die Grippe wird eine
\Es{Schutzimpfung} angeboten. Sie basiert auf der Voraussage
welcher Stamm in der jeweils bevorstehenden Grippewelle
vorherrschen wird und wirkt daher nur für ein Jahr und muss
jede Saison wiederholt werden. Die Impfung bietet keinen
hundertprozentigen Schutz, wird aber trotzdem für
Risikopatienten und Personen die einem erhöhten
Infektionsrisiko ausgesetzt sind (darunter fällt auch
Personal im Gesundheitsbereich) empfohlen.
}
{
% TODO VIT: Slide fehlt!
}
{}
{}
{}

\cite{Gallaher2009,H1N1InvTeam2009,Peiris2009,Wolf2009,Cutler2009}



\Layout{2}{Schweinegrippe, Vogelgrippe -- Warum die große
Aufregung?}
{
Grippeviren befallen nicht nur den Menschen, auch
(Wasser-)Vögel und andere Tiere erkranken an der Grippe.
Normalerweise erfolgt keine Übertragung vom Tier zum
Menschen. Überschreitet ein Erreger diese Grenze wird es
gefährlich: Die Tierwelt beinhaltet dann einen nahezu
unbegrenzten Vorrat an Keimen (Keimreservoir). Dieser kann 
praktisch nicht kontrolliert werden, schließlich lassen
sich Zugvögel nicht so leicht unter Quarantäne stellen
\ldots 

Glücklicherweise entwickelten sich in letzter Zeit nur
Viren, welche zwar vom Tier auf den Menschen und dann
weiter auf den Menschen übertragbar waren, jedoch waren sie
(bis jetzt) nicht übermäßig infektiös oder zeigten keinen
besonders schweren Verlauf. 

\subparagraph{Die Gefahr} Die
Tier\,→\,Mensch\,→\,Mensch-Übertragung stellt eine ständige
Gefahr dar. 

Ein weiteres Problem ist jedoch die mögliche Kreuzung
mit Stämmen die sehr infektiös und schwer verlaufend sind. Dies kann passieren wenn ein
Patient mit mehreren Stämmen gleichzeitig infiziert ist.
Daher ist es wichtig die Infektionen so rasch als möglich
einzudämmen. 
}
{
% TODO VIT: Slide fehlt!
}
{}
{}
{}



\Fazit{Die große Gefahr ist die Entwicklung eines hoch
infektiösen und schwer verlaufenden Stammes, welcher eine
Tier-Mensch-Mensch-Übertragung aufweist.}
\Fazit{Eine rasche Isolierung der Infzierten soll die
Wahrscheinlichkeit einer bösartigen Kreuzung minimieren!}

\subparagraph{Neue Grippe ("`Schweinegrippe"')} Die jüngste
Entwicklung war die Übertragung des Stammes \Code{A(H1N1)}
vom Schwein auf den Menschen in Mexiko im Frühjahr 2009. Der
neue Stamm breitete sich sehr rasch um die ganze Welt aus,
\E{die Erkrankung verlief jedoch recht mild}. Bemerkenswert
war jedoch, dass -- anders als bei einer "`normalen
Grippe"' üblich -- besonders junge Menschen infiziert wurden. 

Man darf gespannt sein was die Zukunft bringt \ldots

% \begin{tabulary}{\linewidth}{lLL}\toprule
% 
% \TabHeading{Übertragung}
% &
% \TabHeading{~}
% &
% \TabHeading{Einschätzung}
% \\\midrule
% Mensch → Mensch
% &
% &
% \\
% Tier → Mensch
% &
% &
% \\
% Tier → Mensch → Mensch
% &
% &
% \\
% Tier → Mensch → Mensch
% &
% hohe Ansteckungsrate
% 
% schwerer Verlauf
% &
% sehr schlimm
% \\\bottomrule
% \end{tabulary}

%%%%%%%%%%%%%%%%%%%%%%%%%%%%%%%%%%%%%%%%%%%%%%%%%%%%%%%%%%%
%%%%%%%%%%%%%%%%%%%%%%%%%%%%%%%%%%%%%%%%%%%%%%%%%%%%%%%%%%%
%%%%%%%%%%%%%%%%%%%%%%%%%%%%%%%%%%%%%%%%%%%%%%%%%%%%%%%%%%%
\subsection{Grippaler Infekt}\index{grippaler
Infekt}

\emph{\Lang{Syn.} \index{Verkühlung}Verkühlung,
 \index{Erkältung}Erkältung} 

\Layout{2}{\TxBeschreibung}
{
Die Verkühlung ist eine Infektion der oberen Luftwege mit Viren,
allerdings mit anderen als bei der "`echten
Grippe"' (Influenza). Die Symptome ähneln denen der
Grippe, allerdings sind sie bei weitem nicht so stark ausgeprägt. Im
Vordergrund stehen Halsschmerzen, Husten, rinnende oder verstopfte Nase
und eine eher leicht erhöhte Temperatur. 
}{}{}{}{}



%%%%%%%%%%%%%%%%%%%%%%%%%%%%%%%%%%%%%%%%%%%%%%%%%%%%%%%%%%%
%%%%%%%%%%%%%%%%%%%%%%%%%%%%%%%%%%%%%%%%%%%%%%%%%%%%%%%%%%%
%%%%%%%%%%%%%%%%%%%%%%%%%%%%%%%%%%%%%%%%%%%%%%%%%%%%%%%%%%%
\subsection{Lungenentzündung -- \T{Pneumonie}}
\label{topic:pneumonie}

Streng genommen ist die Lungenentzündung eine Erkrankung der
Atemwege. Meistens steht bei den Symptomen eher die
entzündliche Komponente im Vordergrund (Fieber, allgemeine
Schwäche, \ldots).

\Layout{2}{Ursachen}
{
\begin{itemize}
  \item Infektiös: Bakterien, Pilze
  \Z{(Pneumokokken,Chlamyden, Candida bei
  HIV/Immunsupression)}
  \item Chemische Schädigung: Spätfolge nach
  \index{Pneumonie!Aspiration}Aspiration
\end{itemize}

Höheres Lebensalter begünstigt Infektionen wie z.\,B.
Pneumonien, da ältere Patient ein zunehmend schwächeres
Abwehrsystem haben.
}{}{}{}{}


\Layout{2}{\TxSymptome}
{
Die Symptome sind oft nicht charakteristisch, eine Diagnose kann meist
erst im Krankenhaus gestellt werden. 

\begin{itemize}
  \item Dyspnoe, Husten
  \item Fieber (muss aber nicht sein)
  \item Reduzierter Allgemeinzustand (AZ)
\end{itemize}
}{

}{}{}{}





\MassnahmenNG{Spezielle Maßnahmen}{Pneumonie}{}{

\begin{itemize}
  \item Symptomatisch
  \item Wärmeerhalt
  \item Hospitalisierung, Abteilung: Innere Medizin
\end{itemize}
}

\subsection{Bakterielle Meningitis}
\index{Meningitis}
\label{topic:meningitis}

\Definition{Eitrige Entzündung der Hirnhäute aufgrund
einer Infektion mit Bakterien.}

\noindent Die bakterielle Meningitis ist eine gefürchtete Erkrankung. Einerseits
ist sie hoch ansteckend, andererseits kann sie rasch tödlich enden. Gefährdet sind
grundsätzlich alle Altersgruppen, besonders aber Kinder, Jugendliche und alte
Menschen.
 
 
\Layout{0}{\TxSymptome}
{%
Leitsymptome sind \Es{Fieber}, \Es{Nackensteifigkeit} (Meningismus) und
\Es{Lichtscheu}. Darüber hinaus kann es zu anderen unspezifischen
Symptomen wie Kopfschmerzen, Übelkeit, Erbrechen, etc. kommen. 

\E{Die Meningitis wird oft mit einem einfachen grippalen Infekt oder einer
Grippe verwechselt!}

\Z{Im fortgeschrittenen Stadium kann es zu massiven
Gerinnungsstörungen kommen, am Körper bilden sich dann dunkle
Hämatome. Einer der
häufigsten Erreger sind die Meningokokken, man spricht dann
von einer Meningokokken-Meningitis.}

}{
\begin{itemize}
  \item Fieber, Kopfschmerzen.
  \item Nackensteifigkeit.
  \item Lichtscheu.
  \item Erbrechen, Übelkeit.
  \item Hämatome am ganzen Körper (im Endstadium).
  \item Verwechslungsgefahr mit "`normalem Infekt"'!
\end{itemize}

}{}{}{}

\opt{STATUS-EXPERIMENTAL}{

\Layout{22}{}
{
% Endstadium einer bakteriellen Meningitis.
% Die Flecken auf der Haut sind Blutungen in Folge einer nicht mehr
% funktionierenden Blutgerinnung. Der Tod ist bei dieser
% Patientin innerhalb weniger Stunden zu erwarten.

\TODO{GABS}{2010-04-18}{Bild: Meningitis von KAR + Beschreibung}{}
}{}
% {./images/meningitis-kind-waterhouse.png}
{./icons/under-construction.png}
{\label{fig:meningitis-baby}}
{unbekannt}
}



\AuthNotes{FEHLT: Querverweis Meningimus\\FEHLT: Maßnahmen}


\MassnahmenNG{Spezielle Maßnahmen}{Meningitis}{}{
\begin{itemize}
  \item Schutzmaske für Personal und Patient (wenn zumutbar).
  \item Rückmeldung an Leitstelle vor Transport, Voranmeldung im Krankenhaus.
  \item Patient erst in das Spitalsgebäude bringen, wenn
  Ziel und Weg dorthin klar ist.
  \item Desinfektionsmaßnhamen gem. Hygieneplan.
  \item Personal bekommt im Spital evtl. eine Antibiotika-Prophylaxe.
\end{itemize}

}

\subsection{Wundstarrkrampf -- Tetanus}
\index{Tetanus}
\index{Wundstarrkrampf}


%\includegraphics[width=5.984cm,height=3.539cm]{images/AASdev20090228-img97.gif}
%\Captionof{figure}{ \AuthNotes{Tetanus Bild in gif:img97} }



Der Erreger des Wundstarrkrampfes\ZFootnote{Clostridium tetani} kommt
\E{überall in der Umgebung} vor. Er produziert einen Giftstoff (Toxin),
welches einen tonischen Krampf der Skelettmuskulatur
auslöst. \Z{Die Inkubationszeit beträgt ca. 4-21 Tage.
Erste Vorzeichen sind Müdigkeit und Inappetenz.}


\Layout{20}{\TxSymptome}
{%
Es kommt zu tonischen \Es{Krämpfen der
quergestreiften Muskulatur}, oft beginnend bei der
Kaumuskulatur. Der Patient hat starke Schmerzen und
entwickelt in Folge der Krämpfe einen \Es{Atemlähmung}. Die
\E{Sterblichkeit ist hoch}, die Letalität beträgt ca.
50\,\%.
}
{
% TODO VIT: Slide fehlt!
}
{./images/hygiene/6373_lores_tetanus.jpg}
{\AuthNotesPicLegalOk{PD}{}Patient mit Tetanus.}
{Public Health Image Library, ID\#: 6373, Public Domain.}

\Layout{0}{Schutzimpfung}
{
\index{Impfschutz!Tetanus}
\index{Tetanus!Impfschutz}
\label{topic:tetanus-impfschutz}
Es existiert
eine Schutzimpfung, diese muss dem Patienten \E{rechtzeitig} nach der
Verletzung gegeben werden. Der Impfschutz beträgt
normalerweise ca \E{10 Jahre}, bei einer
Verletzung wird sicherheitshalber allerdings schon bei einer
\Es{länger als 5 Jahre} zurückliegenden Impfung eine neuerliche Gabe verabreicht.

Auch bei kleinen Verletzungen muss nach dem Impfschutz
gefragt werden, dies muss auch unbedingt \Es{dokumentiert
werden!} Ist das Datum der letzten Impfung nicht
sicher erinnerlich soll der Patient unverzüglich zu Hause im
Impfpass nachschauen. Erforderlichenfalls soll der Patient
unverzüglich eine Unfallabteilung aufsuchen. 

Die entsprechende Aufklärung muss am Revers bestätigt werden!
}
{
% TODO VIT: Slide fehlt!
}
{}
{}
{}

\Cave{Eine mangelhafte Aufklärung oder Dokumentation
bezüglich Tetanus und dessen Impfschutz ist grob fahrlässig und ein unverzeihbarer
Kunstfehler.}


%%%%%%%%%%%%%%%%%%%%%%%%%%%%%%%%%%%%%%%%%%%%%%%%%%%%%%%%%%%
%%%%%%%%%%%%%%%%%%%%%%%%%%%%%%%%%%%%%%%%%%%%%%%%%%%%%%%%%%%
%%%%%%%%%%%%%%%%%%%%%%%%%%%%%%%%%%%%%%%%%%%%%%%%%%%%%%%%%%%


\subsection{Tuberkulose}
\label{topic:tuberkulose}

\emph{\Z{Mykobacterium Tuberculosis, "`Schwindsucht"'.}}
Abk.: "`\T{Tbc}"'

\Layout{0}{Beschreibung}
{
Die Tuberkulose kann praktisch alle Organe befallen, der
Verlauf wird vor allem davon bestimmt, welches
Organ befallen ist \Z{(Lunge, Skelett, Darm, Haut,
Genitalien, ...)}. Ausgangspunkt ist fast immer die \Es{Lunge}.

Die Übertragung erfolgt vor allem \Es{aerogen}, aber unter
Umständen auch oral, durch Hautkontakt oder plazentär. 

Zunehmend treten auf Antibiotika resistente Erreger auf \Z{(MDR-TB,
XDR-TB)\footnote{Quelle: AGES}}
% TODO LIT, HYG: MDR/XDR-TB Quelle bei AGES
.

\TODO{GABS}{2010-04-18}{Bild: Bild Lungenröntgen TB-Patient als Eye-Catcher.
--> Hauer}{}
    
\Z{Die Inkubationszeit beträgt 4-12 Wochen.}
}{}
% {./images/AASdev20090228-img98.jpg}
{./icons/under-construction.png}
{\AuthNotesPicLegalNotOk{img98}{Durch Foto von Hauer
ersetzen}Röntgenbild eines TB-Pat\-ient\-en.}{}


\Layout{0}{\TxSymptome}
{%
Es dominieren vor allem Allgemeinsymtptome wie erhöhte
Temperatur, Nachtschweiß, Lymphknotenschwellung und Husten.
Auffallend ist manchmal die Schwere des Hustens, Blutbeimengungen gelten
als Alarmsignal. 

In der Anamnese ist mitunter eine bereits druchgemachte
Tuberkulose erkennbar.
}{
\begin{itemize}
  \item Husten, Bluthusten.
  \item erhöhte Temperatur.
  \item \Es{Nachtschweiß, Lymphknotenschwellung}
  \item Anamnese!
\end{itemize}
}
{}
{}
{}


\Layout{2}{Offene Tuberkulose}
{%
\index{Offene Tb}%
Wenn der Infektionsherd in der Lunge \E{Luftkontakt} hat, spricht man von der \T{Offenen
Tuberkulose}. Der Patient ist dann infektiös,
eine Tröpfcheninfektion ist möglich. Die Verwendung von
Schutzmasken ist anzuraten. 

Es ist äußerlich \E{nicht sicher erkennbar} ob der Patient
an einer offenen oder geschlossenen Tuberkulose leidet.
Bluthusten wäre allerdings ein deutlich Hinweiszeichen.
}
{
% TODO VIT: Slide fehlt!
}
{}
{}
{}


\Z{
\paragraph{Ansteckungsgefahr}
Bei Einhaltung der Schutzmaßnahmen (s. Spezielle Maßnahmen) ist eine Ansteckung
von sonst gesunden und immunkompetenten Personal unwahrscheinlich. Für eine
Ansteckung  wäre ein ca. 8-stündiger Kontakt erforderlich
\footnote{National Institute for Clinical Excellence (NICE) Tuberculosis.
National clinical Guideline for Diagnosis, management, prevention and control
(2006)}
% TODO HYG, LIT: National Institute for Clinical Excellence (NICE)
% Tuberculosis. National clinical Guideline for Diagnosis, management,
% prevention and control (2006)
.
}

\Massnahmen{
\paragraph{Spezielle Massnahmen: Tuberkulose}~

\begin{itemize}
  \item Bei Verdacht auf eine offene Tuberkulose
  Schutzmasken verwenden: FFP~3 für das Personal, "`OP-Maske"' für den
  Patienten, wenn zumutbar.
  \item \Es{Meldung an Leitstelle \& Betriebsleitung!}
  \item Desinfektion laut Hygieneplan.
\end{itemize}
}


%%%%%%%%%%%%%%%%%%%%%%%%%%%%%%%%%%%%%%%%%%%%%%%%%%%%%%%%%%%
%%%%%%%%%%%%%%%%%%%%%%%%%%%%%%%%%%%%%%%%%%%%%%%%%%%%%%%%%%%
%%%%%%%%%%%%%%%%%%%%%%%%%%%%%%%%%%%%%%%%%%%%%%%%%%%%%%%%%%%
\subsection{HIV und AIDS}
\index{HIV}\index{Humanes Immundefizienz
Virus}\label{topic:hiv}\label{topic:aids}

\T{HIV}: \E{Humanes Immundefizienz Virus}. 
\T{AIDS}: \E{Aquired Immune Deficiency Syndrom},
\E{erworbenes Immunschwäche-Syndrom}

\Layout{22}{HIV und AIDS}
{
\subparagraph{HI-Virus} Das HIV befällt vor allem die Zellen des
Abwehrsystems. Es vermehrt sich in ihnen, setzt sie außer Funktion und zerstört sie schließlich. Das körpereigene Abwehrsystem kann HIV nicht aus dem Körper
    entfernen.
Es kommt zu einer Immunschwäche.

\subparagraph{AIDS} HIV verursacht das \E{AIDS} (Aquired Immune Deficiency
Syndrom, erworbenes Immunschwäche-Syndrom), welches die
eigentliche Erkrankung darstellt. Dabei kommt es durch das
\E{geschwächte Immunsystem} zu Folgeerkrankungen, an denen der Patient versterben kann:

\begin{itemize}
  \item \E{Infektionen} mit normalerweise harmlosen
  Erregern
  \begin{itemize}
    \item Pilze, "'fakultative{\textquotedblright} Erreger,
    "'Opportunisten{\textquotedblright}, seltene Arten der
    Lungenentzündung
  \end{itemize}
  \item \E{Tumore}
\end{itemize}

Abhängig von der durchgeführten Therapie bricht das AIDS
einige Jahre nach der eigentlichen HIV-Infektion aus.
}{
% TODO VIT: Slide fehlt!
% \begin{itemize}
%     \item \Z{Point}
% \end{itemize}
}
{./images/hygiene/Niaid-hiv-virion-mod.jpg}
{Schema eines HI-Virus.}
{US National Institute of Health, PD-USGov-HHS-NIH.}









\Layout{2}{Übertragung}
{%
Übertragen wird das HI-Virus v.\,a. durch Körperflüssigkeiten
wie Blut, Blutprodukte, Sperma und Scheidensekret.
\E{"`Ungeschützter Sex"'}, Transfusionen und
\E{Nadelstichverletzungen} gelten als klassische Risiken.
Es erfolgt \Es{keine Übertragung durch normale
Sozialkontakte} wie Händeschütteln, wohnen im
gemeinsamen Haushalt, gemeinsame Benutzung von Geschirr und
Toilettenanlagen, etc.
}{
% TODO VIT: Slide fehlt!
}
{}
{}
{}

\Layout{2}{Test}
{Erst nach 6-12 Wochen sind Antikörper nachweisbar. Das bedeutet dass der
\Es{HIV-Test nicht "`aktuell", sondern "`12 Wochen blind"} ist!
}
{
% TODO VIT: Slide fehlt!
}
{}
{}
{}


\Layout{2}{Therapie}
{%
\begin{itemize}
  \item Derzeit ist \Es{keine Heilung} möglich, jedoch
  kann eine \E{deutliche Lebensverlängerung} durch
  Medikamente erreicht werden.
\end{itemize}
}
{
% TODO VIT: Slide fehlt!
}
{}
{}
{}

\Layout{2}{Prophylaxe}
{%
% TODO VIT: Fließtext fehlt!
\begin{itemize}
  \item Allgemeine Selbstschutzmaßnahmen bei der
  Patientenversorgung: Schutzhandschuhe, Kontakt mit
  Körperflüssigkeiten vermeiden.
  \item Postexpositionsprohylaxe bei Nadelstichverletzungen
  \item Kondome
\end{itemize}
}
{
% TODO VIT: Slide fehlt!
}
{}
{}
{}




%%%%%%%%%%%%%%%%%%%%%%%%%%%%%%%%%%%%%%%%%%%%%%%%%%%%%%%%%%%
%%%%%%%%%%%%%%%%%%%%%%%%%%%%%%%%%%%%%%%%%%%%%%%%%%%%%%%%%%%
%%%%%%%%%%%%%%%%%%%%%%%%%%%%%%%%%%%%%%%%%%%%%%%%%%%%%%%%%%%
\subsection{Hepatitis}
\label{topic:hepatitis-virale}%
\label{topic:hepatitis-toxisch}%
\label{topic:hepatitis}%
\index{Hepatitis!infektiöse}%
\index{Hepatitis!toxische}%

\Definition{Leberentzündung durch verschiedene (virale)
Erreger (A -- E) oder durch sonstige
(chronische) Schädigung.}



\Layout{0}{\TxBeschreibung{} und Ursachen}
{%%% Gerhard Gabriel <gerhard.gabriel@grissu.org> 
Hepatitis ist eine
entzündliche Erkrankung der Leber. Diese Entzündung kann durch Erreger wie
Viren, oder Mikroorganismen hervorgerufen werden
(infektiöse, virale Hepatitis). Aber auch
regelmäßige Schädigungen der Leber durch Giftstoffe wie
Alkohol, Medikamente \Z{(z.\,B. Paracetamol)}, etc. können
zu einer (chronischen) Entzündung führen. Auch ein länger
bestehender Rückstau der Gallenflüssigkeit, z.\,B. bei einer
Verstopfung des Gallenganges, kann zu einer Entzündung
führen.

Je nach Ursache der Hepatitis kann es vorkommen, dass es zu
keiner Abheilung kommt und die Hepatitis dauerhaft bestehen
bleibt. Man spricht dann von der \T{chronischen
Hepatitis}.} 
{
\begin{itemize}
  \item Virusinfektionen (Hep. A, B, C, \ldots ).
  \item Gallen-Abflussstau.
  \item Vergiftungen Lebensmittel(Pilze)~/
  Medikamente (Paracetamol)~/ chron. Alkoholismus.
  \item Chronische Verlaufsform möglich.
\end{itemize}
}{}{}{}

\Layout{0}{\TxSymptome}
{%
Die Symptomatik einer Hepatitis ist
uncharakteristisch und eher allgemein: Unspezifisches mehr oder minder
ausgeprägtes Krankheitsgefühl, ähnlich einem grippalem Infekt, Appetitlosigkeit.
Auch erhöhte Körpertemperatur und Schmerzen im
rechten Oberbauch kommen bisweilen vor. Die Leber ist oft
vergrößert.

Im
fortgeschrittenen Stadium kann es zu einer Gelbfärbung
des Augenweiß und eine gelbliche Verfärbung der Haut
(Ikterus) kommen. 

Bei einem chronischen
Verlauf kommt es zu einem Umbau von
funktionierendem Lebergewebe zu
funktionslosem Bindegewebe. Man spricht
dann von der
\T{Leberzirrhose}\label{topic:leberzirrhose}\index{Leberzirrhose}.
Schließlich kann es zu einem Leberversagen und in kurzer
Folge zum Tod kommen, da die Leber ein überlebenswichtiges
Organ ist. 
}{
\begin{itemize}
  \item uncharakteristische Allgemeinsymptome (grippale Symptome,
  Appetitlosigkeit).
  \item Schmerzen im re. Oberbauch.
  \item Gelbfärbung Augen/Haut (Ikterus ).
  \item Vergrößerte, schmerzhafte Leber.
  \item Ikterus (Gelbsucht).
  \item Unterscheide "'akut{\textquotedblright} und
  "'chronisch{\textquotedblright}.
  \begin{itemize}
    \item Chronisch: wenn Hepatitis nicht abheilt.
  \end{itemize}
  \item Ende: Leberversagen, Tod.
\end{itemize}
}{}{}{}



\paragraph{Verschiedene Arten der viralen Hepatitis}~

Abhängig vom jeweiligen Virus unterscheidet man
verschiedene Arten der viralen Hepatitis:

\begin{table}[H]
\begin{gWideParEnv}
\Captionof{table}{Übersicht der wichtigsten viralen Hepatitis-Arten.}

\begin{tabular}{p{3cm}p{\linewidth-4cm}}
\TabHeading{Typ}
&
\TabHeading{Bemerkungen}
\\\midrule\addlinespace
\TabSub{Hepatitis-A} 

Hepatitis-A-Virus (HAV)
&
\begin{itemize}
  \item \Z{Inkubationszeit 10 - 40 Tage}
  \item typische Reisekrankheit (Mittelmeerraum, Südamerika, Orient),
  meist epidemisch,
  \item Dauer einige Wochen,
  \item kein Übergang in chronischen Verlauf
\end{itemize}

\\\addlinespace
\TabSub{Hepatitis-B} 

Hepatitis-B-Virus (HBV)
&
\begin{itemize}
  \item \Z{Inkubationszeit 4 - 12 Wochen,}
  \item \Z{meist sporadisches Auftreten}
  \item Übertragung durch Speichel, Blut(produkte),
  Sperma und andere Körperflüssigkeiten,
  \item chronischer Verlauf in ca. 5 \%
\end{itemize}
\\\addlinespace
\TabSub{Hepatitis C}

Hepatitis-C-Virus (HCV)
&
\begin{itemize}
  \item \Z{Inkubationszeit 6 - 12 Wochen,}
  \item häufig nach Bluttransfusion oder i.v. Drogenabusus,
  \item \Es{oft chronischer Verlauf} (in 70 -- 80~\% der
  Fälle)
\end{itemize}
\\\bottomrule
\end{tabular}
\end{gWideParEnv}
\end{table}



\Layout{0}{Impfung}
{%
\index{Impfschutz!Hepatitis}%
\index{Hepatitis!Impfschutz}%
Für die Hepatitis A und B ist eine Impfung
verfügbar (z.\,B. Twinrix{\texttrademark}). Für die Hepatitis
C ist derzeit noch keine Impfung regulär am Markt. Gegen
die toxischen Formen der Hepatitis gibt es natürlich auch keine Impfung.

Achtung: Non-Responder möglich!
% FEATURE Hepatitis Laborwerte Interpretation einfügen.
}{
\begin{itemize}
  \item Für Hep. A \& B.
  \item Nicht für Hep. C.
\end{itemize}
}
{}
{}
{}


\Fazit{Für den Rettungsdienst sind besonders die Hepatitis
B und C wichtig: Hohe Infektiosität, chronischer Verlauf!

Gegen Hepatitis C gibt es derzeit keine Impfung!

Für Risikopersonal (dazu zählt auch Sanitätspersonal)
übernimmt i.\,d.\,R. die Unfallversicherungsanstalt die
Kosten für die Schutzimpfung.}



%%%%%%%%%%%%%%%%%%%%%%%%%%%%%%%%%%%%%%%%%%%%%%%%%%%%%%%%%%%
%%%%%%%%%%%%%%%%%%%%%%%%%%%%%%%%%%%%%%%%%%%%%%%%%%%%%%%%%%%
%%%%%%%%%%%%%%%%%%%%%%%%%%%%%%%%%%%%%%%%%%%%%%%%%%%%%%%%%%%
\subsection{Sepsis}
\label{topic:sepsis}

Siehe auch: Septischer Schock und Toxischer Schock,
Kap.\,\ref{topic:schock-septischer}.

\begin{itemize}
  \item Wenn Keime und deren Produkt außer Kontrolle geraten,
  erfolgt Amoklauf des Körpers:
  \item Ungehemmte Freisetzung von Stoffen des Entzündungs-, Gerinnungs-
  und Immunsystems
  \item Auch ohne Infektion möglich:
  \begin{itemize}
    \item \Z{Systemic Inflammatory Response Syndrome (\Abk{SIRS})}
    \item nach schwerem Trauma, Verbrennungen, Gewebshypoxie
  \end{itemize}
\end{itemize}


\Layout{2}{\TxSymptome}
{%
Die Symptome sind sehr vom Fortschritt und
körperlichen Abbau abhängig. Grundsätzlich:


\begin{itemize}
  \item Sehr stark reduzierter Allgemeinzustand
  \item Fieber ({\textgreater}\,38{\textcelsius}) oder
  Hypothermie ({\textless}\,36{\textcelsius})
  \item Herzfrequenz {\textgreater}\,90\nicefrac{}{min}
  \item Atemfrequenz {\textgreater}\,20\nicefrac{}{min} oder
  Hyperventilation
  \item unkontrollierbare Gerinnung/Blutung
  \item Hypotonie ${\rightarrow}$ Septischer Schock
  \item Organversagen
\end{itemize}

}{}{}{}{}


%%%%%%%%%%%%%%%%%%%%%%%%%%%%%%%%%%%%%%%%%%%%%%%%%%%%%%%%%%%
%%%%%%%%%%%%%%%%%%%%%%%%%%%%%%%%%%%%%%%%%%%%%%%%%%%%%%%%%%%
%%%%%%%%%%%%%%%%%%%%%%%%%%%%%%%%%%%%%%%%%%%%%%%%%%%%%%%%%%%





% \clearpage
%%%%%%%%%%%%%%%%%%%%%%%%%%%%%%%%%%%%%%%%%%%%%%%%%%%%%%%%%%%
%%% Abdomen                                             %%%
%%%%%%%%%%%%%%%%%%%%%%%%%%%%%%%%%%%%%%%%%%%%%%%%%%%%%%%%%%%

\section{Abdominelle Erkrankungen}



\Layout{22}{Der Bauch als Herausforderung}
{%
Der Bauch ist eine echte Herausforderung. Er
beinhaltet viele unterschiedliche Organe und Strukturen
welche weh tun und
erkranken können. Viele dieser, oft auch mit starken Schmerzen verbundenen, Störungen
sind nicht übermäßig schlimm, einige sind jedoch sehr gefährlich oder können rasch
gefährlich werden. 
}
{}
{./images/operation.jpg}
{\AuthNotesPicLegalOk{}{}Oft die
    Abteilung der Wahl für einen abdominellen Notfall: Die
    Chirurgie. Aber auch die Abteilungen für
    Kinderheilkunde, Kinderchirurgie, Gynäkologie (!) oder
    die Interne können manchmal passende Ziele sein.
    }{Gabriel}
% \begin{wrapfigure}{r}{5cm}
%     \includegraphics[width=\linewidth]{}
%     \Captionof{figure}{
%     \SourcePicture{}}
% \end{wrapfigure}


\paragraph{Anamnese und Untersuchung} Es ist daher wichtig
typische Symptome zu erheben und Alarmzeichen sofort zu erkennen.

Wesentlich ist bei allen unklaren Bauchbeschwerden die Erhebung einer exakten Stuhlanamnese:

\begin{itemize}
  \item Auffälligkeiten? (Geruch, Farbe, Konsistenz, Frequenz)
  \item Wann erfolgte der letzte Stuhlgang?
  \item Winde?
  \item Genaue Beschreibung von Harn und eventuell Erbrochenem.
  \item Bei Frauen im gebärfähigem Alter Regelanamnese.
\end{itemize}

Der Patient muss bis zur endgültigen Klärung der Situation
\Es{nüchtern} belassen werden!

Wenn es der Patient toleriert, muss eine Betrachtung und
Abtastung des unbekleideten Abdomens beim liegenden
Patienten erfolgen wie unter
\ref{topic:untersuchung-abdomen-abtasten} beschrieben.

\subsection{Der schmerzende Bauch}

\Layout{0}{Der Bauch -- unendliche Weiten}
{%
Gleich vorweg: In vielen
Fällen wird der Grund des Bauchschmerzes nicht mit
hinreichender Sicherheit feststellbar sein und es bietet sich
keine sichere Verdachtsdiagnose an. Statt einer Diagnose wird
dann das \Es{Beschwerdebild beschrieben} (Ort des Schmerzes, Art des
Schmerzes, Dauer, Verlauf, Austrahlung, Ergebnis der
abdominellen Tastuntersuchung: Abdomen weich oder hart,
Druckschmerz, Loslassschmerz).
}{
\begin{itemize}
  \item In vielen Fällen Krankheitsbild unklar.
  \item Im Zweifel: Beschreibung statt Diagnose.
\end{itemize}
}{}{}{}

\Layout{2}{Beispiele}
{
\begin{quote}
{\sffamily\bfseries Diagnose:} {\ttfamily Unklarer
Bauchschmerz: Kolikartige Schmerzen im linken Unterbauch seit 1/2h }

{\sffamily\bfseries Anmerkungen:} {\ttfamily keine Ausstrahlung, Druckschmerz im li. Unterbauch, Abdomen sonst  weich.}
\vspace{0.5em}

{\sffamily\bfseries Diagnose:} {\ttfamily Unklarer
Bauchschmerz: Brennend-stechender Schmerz im mittleren
Oberbauch seit 2h mit Ausstrahlung in den Rücken}

{\sffamily\bfseries Anmerkungen:} {\ttfamily Abdomen weich,
keine Druck- oder Loslassschmerzen}
\vspace{0.5em}

{\sffamily\bfseries Diagnose:} {\ttfamily Unklarer
Bauchschmerz: „Dumpfer“ Schmerz im linken und rechten
Oberbauch seit gestern. }

{\sffamily\bfseries Anmerkungen:} {\ttfamily Abdomen
weich, keine Druck- oder Loslassschmerzen}
\end{quote}
}{

}{}{}{}



Dennoch muss man bei jeden Patienten versuchen, so gut als
möglich die in Frage kommenden Krankheitsbilder einzugrenzen
und gefährliche Erkrankungen nach Möglichkeit auszuschließen. Im
folgenden sind einige häufige Krankheitsbilder beschrieben,
die \E{erkennbar}  und relativ häufig sind.

\MassnahmenNG{Allgemeine Maßnahmen}{Abdominalerkrankung}
{\label{m:am-abdominalerkrankungen}}{
\begin{itemize}
  \item \Es{Vitale Bedrohung einschätzen.} \TxBeiVit
  \TxMassVit
  \item Lagerung: grundsätzlich bauchdeckenentspannend
  (Knierolle, Beine angezogen), oder nach Wunsch des Patienten
  \item nüchtern lassen (in Hinblick auf möglicherweise bevorstehende
  Operation, nichts essen oder trinken lassen)
  %     \item Monitoring. Der Zustand kann sich abrupt
  %     verschlechtern.
  \item Hospitalisierung: Situationsgerecht, z.\,B. Abt. f. Chirurgie oder
  Innere Medizin.
  
  Bei Kindern Abt. f. Kinderchirurgie oder Abt. f. Kinderheilkunde (je nach
  Verdachtsdiagnose und regionaler Regelung).
  
  Bei weiblichen Patienten Abt. für Gynäkologie (evtl. im "`Hintergrund"') erwägen.
\end{itemize}
}


\subsection{Häufige und gut erkennbare Erkrankungen}



\subsubsection{Das Akute Abdomen --- Alarmzeichen mit
bretthartem Bauch}\label{topic:akutes-abdomen}



\Definition{\T{Akutes Abdomen}: Akute lebensbedrohliche
Erkrankung des Bauchraumes mit plötzlichen, starken
\Es{Bauchschmerzen} und
\Es{bretthartem Bauch}. Es sind viele Ursachen  möglich!}

\Z{Die Definition des Akuten Abdomens ist oft nicht ganz
eindeutig: Einerseits kann man da\-runter jede plötzlich auftretende
Baucherkrankung verstehen\footnote{Longmore et.al.: Oxford Handbook of
Clinical Medicine, 7\textsuperscript{th} edt., Oxford University Press,
2007\cite{Longmore2007}}\footnote{\cite{Renz-Polster2006}
S.611}.} 

Allerdings ver\-stehen viele Leute darunter eher
eine bedrohliche Bauch\-erkrankung, bei der man akut (schnell) handeln muss (Ab\-wehrspannung, harter Bauch). In der
Notfallmedizin halten wir uns an letztere Definition. 


\Layout{2}{\TxBeschreibung}
{Das „\T{Akute Abdomen}“ ist ein interventionspflichtiger
Symptomenkomplex, der auf einem sich akut oder chronisch
entwickelnden (Stunden, Tage oder auch Wochen) pathologischen
Prozess im Bauchraum beruht.}
{

}{}{}{}

Es hat grundverschiedene Ursachen, die eine gemeinsame
Symptomatik haben.

\Layout{0}{\TxSymptome}
{%
Der Patient liegt meistens im Bett. Oft ist
eine Schonhaltung (angezogene Beine, gekrümmte
Körperhaltung) zu beobachten. Er klagt über Schmerzen im
Bauch die können als diffus (im gesamten Bauch) oder auf ein Bauchsegment
beschränkt, angegeben werden. 

Die harte Bauchdecke („\E{bretthart}“) ist neben den
Schmerzen das Leitsymptom des Akuten Abdomens. 

Es besteht oft Übelkeit oder auch Erbrechen. Der Kreislauf
kann kompensiert sein, oder Zeichen von Schwäche zeigen,
bis hin zu einer Schocksymptomatik (Blutverlust, Sepsis).
} 
{
\begin{itemize}
  \item \Es{Schmerz}: diffus (ganzer Bauch) oder
  lokalisierbar
  \item \Es{Harte  Bauchdecke} (bretthart)
  \item Schonhaltung  (angezogene Beine, Patient
  krümmt sich)
  \item Übelkeit, Erbrechen
  \item evtl. Stuhl-, Windverhalten
  \item evtl. Fieber
  \item evtl. Schock
\end{itemize}
}{}{}{}


Diese Auflistung der Symptome ist in der Praxis nicht
unbedingt vollständig bei einem Patienten nachzuweisen. Die
Intensität und die Anzahl der angeführten Symptome ist
einerseits vom Allgemeinzustand, vom Lebensalter, von
anderen vorbestehenden Erkrankungen des Patienten und der
Dauer des aktuellen Zustandes abhängig.

Wesentlich nach Eintreffen beim Patienten, ist die
Beantwortung der Fragen \emph{"`Wie schlecht, bzw. wie gut
geht es dem Patienten? Ist der Patient akut vital
bedroht?"'}.


\Layout{0}{Ursachen}
{%
Grundsätzlich kommt fast jede Erkrankung oder Verletzung in
Betracht die Strukturen und Organe im Bauchraum betrifft,
entsprechend vielfältig sind auch die Ursachen: 

\begin{multicols}{2}
\begin{itemize}
  \item Bauchfellentzündung (Peritonitis)
  \item Darmverschluss (Ileus)
  \item Blinddarmentzündung (Appendizitis)
  \item Entzündung der Bauchspeicheldrüse  (Pankreatitis)
  \item Entzündung des Magens / Magengeschwür
  \item Durchbruch / Zerreißung eines Hohlorganes
  \item Blutungen in den Bauchraum
  \item Bauchaortenaneurysma
  \item Mesenterialinfarkt
  \item andere abdominelle Erkrankungen
  \item Entzündung der Leber (Hepatitis )
  \item Gynäkologische Erkrankungen
  \item \ldots
\end{itemize}
\end{multicols}
}{
\begin{itemize}
  \item Bauchfellentzündung (Peritonitis)
  \item Blinddarmentzündung (Appendizitis)
  \item Entzündung der Bauchspeicheldrüse  (Pankreatitis)
  \item Durchbruch / Zerreißung eines Hohlorganes
  \item Blutungen in den Bauchraum
  \item Bauchaortenaneurysma
  \item andere abdominelle Erkrankungen
  \item Gynäkologische Erkrankungen
  \item u.v.a.m. \ldots
\end{itemize}
}{}{}{}





\MassnahmenNG{Spezielle Maßnahmen}{Akutes Abdomen}
{\label{m:akutes-abdomen}}
{
\begin{itemize}
  \item Allgemeine Maßnahmen bei Abdominalerkrankungen
  (\prettyref{m:am-abdominalerkrankungen})
  \item \Es{Vitale Bedrohung einschätzen.} \TxBeiVit
  \TxMassVitBes
  \begin{itemize}
    \item \ce{O2} 10\LMin
    \item Lagerung: je nach Situation: bauchdeckenentspannend,
    Beine hoch, \ldots
  \end{itemize}
  \item Lagerung: grundsätzlich bauchdeckenentspannend
  (Knierolle, Beine angezogen), oder nach Wunsch des Patienten
  \item nüchtern lassen (in Hinblick auf möglicherweise bevorstehende
  Operation, nichts essen oder trinken lassen)
  \item Monitoring. Der Zustand kann sich abrupt
  verschlechtern.
  \item Hospitalisierung: Abt. f. Chirurgie, bei Kindern ${\rightarrow}$
  Abt. f. Kinderheilkunde, bei Frauen evtl. Gynäkologie oder Chirurgie mit
  Abt. f. Gynäkologie im Haus \opt{REGVIE}{(z.\,B. Donauspital, Rudolfstifung,
  Wilhelminenspital, etc.)}
\end{itemize}
}


\subsubsection{Darmverschluss -- Ileus}\index{Ileus}

\Definition{Unterbrechung der normalen  Darmpassage}

\Z{\index{Ileus}\T{Ileus}  (vollständig)

\index{Subileus}\T{Subileus}  (unvollständig)}

\Layout{2}{\TxBeschreibung}
{%
Der Verschluss kann mechanisch bedingt sein (durch
Hindernisse wie Verwachsungen oder Tumore), oder durch eine
Lähmung der Darmmuskulatur (Medikamente, Gifte, Entzündung
\ldots) verursacht sein. Bestimmte Schmerzmittel
\Z{(Opiate)}, welche oft bei Krebspatienten eingesetzt
werden, sind bekannt dafür Darmlähmungen zu erzeugen. 
}{
\begin{itemize}
  \item Mechanische Blockade oder Lähmung.
\end{itemize}
}{}{}{}


\Layout{0}{\TxSymptome}
{
Typisch für den Darmverschluss ist die mangelhafte
Stuhlproduktion. Es kommt zu diffusen Bauchschmerzen und einem
Blähbauch, sowie zu Übelkeit und Erbrechen bis hin zum
Erbrechen von zurück gestautem Kot (selten).
}{
\begin{itemize}
  \item Diffuse Bauchschmerzen
  \item Stuhlverhalten
  \item Übelkeit, Erbrechen (im schlimmsten Fall Koterbrechen)
\end{itemize}
}{}{}{}



\MassnahmenNG{Spezielle Maßnahmen}{Darmverschluß}{}{
\begin{itemize}
  \item Allgemeine Maßnahmen bei Abdominalerkrankungen
  (\prettyref{m:am-abdominalerkrankungen})
  \item \Es{Vitale Bedrohung einschätzen.} \TxBeiVit
  \TxMassVitBes
  \begin{itemize}
    \item \ce{O2} 10\LMin
    \item Lagerung: je nach Situation: bauchdeckenentspannend,
    Beine hoch, \ldots
  \end{itemize}
  \item Patient nüchtern lassen.
  \item Hospitalisierung: Abt. f. Chirurgie.
\end{itemize}
}


\subsubsection{Blinddarmentzündung -- Appendizitis}



\Definition{Entzündung des \Es{Wurmfortsatzes}
(\Es{Appendix} \Z{vermiformis}) des Blinddarms}

\Layout{0}{\TxSymptome{} und Verlauf}
{%
\ACh{GABG}{2009-05-07}{NEU}Im Vordergrund der akuten Appendizitis
steht anfänglich meist ein als eher diffus empfundener, heftiger Schmerzzustand im Bereich der
Nabelgegend\ACh{GABS}{2009-05-07}{Magengegend --> Nabelgegend}, der
innerhalb weniger Stunden in den rechten Unterbauch
wandert. Übelkeit, Brechreiz und Verstopfung
\ACh{GABS}{2009-05-07}{Obstipation --> Verstopfung} können das Krankheitsbild begleiten. Meist ist die Körpertemperatur bis 38\textcelsius{} leicht erhöht
\ACh{GABS}{2009-05-07}{gelöscht da unerheblich für RS: , wobei eine rektale und
axilläre Messung eine Differenz um ca. 1\textcelsius aufweist}. 
\ACh{GABS}{2009-05-07}{alt: Unbehandelt kann der weitere Verlauf durch eine
Abszeßbildung und Perforation in
die Bauchhöhle gekennzeichnet sein. neu: Unbehandelt kann es zu einer Eiterung
und einem Durchbruch in die Bauchhöhle kommen.} Unbehandelt kann es zu einer
Eiterung und einem Durchbruch in die Bauchhöhle kommen. In diesem Stadium entwickelt sich ein
\ACh{GABS}{2009-05-07}{alt: Schockzustand, der eine absolute Lebensbedrohung
darstellt; neu: akutes Abdomen, siehe \ref{topic:akutes-abdomen}} akutes
Abdomen (\prettyref{topic:akutes-abdomen}). 

Der Verlauf einer akuten Appendizitis ist
jeweils durch individuelle Faktoren wie Lebensalter und Allgemeinzustand des
Patienten geprägt. 
}{
\begin{itemize}
  \item Schmerzen im re. Unterbauch
  \item \E{Druck}schmerzen im re. Unterbauch
  \item \E{Loslass}schmerz li. Bauch
  \item Abwehrspannung
  \item Übelkeit, Erbrechen
  \item evtl. Fieber
\end{itemize}
}{}{}{}

\MassnahmenNG{Spezielle Maßnahmen}{Appendizitis}{}{
\begin{itemize}
  \item Allgemeine Maßnahmen bei Abdominalerkrankungen
  (\prettyref{m:am-abdominalerkrankungen})
  \item Patient nüchtern lassen.
  \item Hospitalisierung: Abt. f. Chirurgie, Sonderfälle (Kinder, Frauen!) beachten.
\end{itemize}
}




\subsubsection{Gallenkolik}
\label{topic:gallenkolik}
\index{Kolik!Gallen-}
\index{Gallenkolik}

\Definition{Kolik-artige Schmerzen infolge einer Blockade
der Gallenwege durch Gallensteine}

\Z{\emph{Vgl. Kolik: \prettyref{topic:kolik}}}

\Layout{0}{\TxBeschreibung} 
{Gallensteine verursachen eine Verstopfung des
Gallengangs und folglich einen Rückstau von
Gallenflüssigkeit in die Leber.
} 
{
\begin{itemize}
  \item Verstopfung der Gallenwege (Stein)
  \item → Rückstau und Dehnung
\end{itemize}
}{}{}{}

\Layout{0}{\TxSymptome}
{%
Durch den Rückstau und der Dehnung des Gallengangs kommt es zu
starken, auf- und abschwellenden, krampfartigen
(\Es{kolikartigen}) \Es{Schmerzen} im \Es{rechten Oberbauch},
klassischerweise mit Ausstrahlung in die rechte Schulterregion. Begleitendes Fieber ist möglich.

Der rechte Oberbauch ist äußerst druckschmerzhaft, der
Patient ist sehr \Es{unruhig} und agitiert.
}{
\begin{itemize}
  \item Auf- und abschwellende, krampfartige
  (kolikartige) Schmerzen im re. Oberbauch
  \item Besonders nach fettreichen Mahlzeiten
  \item evtl. Gelbfärbung des Augenweiß (Gelbsucht (Ikterus))
  \item evtl. Übelkeit, Erbrechen
\end{itemize}
}{}{}{}

\MassnahmenNG{Spezielle Maßnahmen}{Gallenkolik}{}{
\begin{itemize}
  \item Allgemeine Maßnahmen bei Abdominalerkrankungen
  (\prettyref{m:am-abdominalerkrankungen}).
  \item Notarzt zur Schmerztherapie erwägen.
  \item Patient nüchtern lassen.
  \item Hospitalisierung: Abt. f. Chirurgie.
\end{itemize}
}


\subsubsection{Ösophagusvarizenblutung}

\Definition{Blutende "`Krampfadern"' in der
Speiseröhre.}\nopagebreak[4]

\Layout{0}{\TxBeschreibung}
{Bei Patienten mit einer fortgeschrittenen Lebererkrankung
(Leberzirrhose, \ref{topic:leberzirrhose}) kann es zur
Bildung von Krampfader-artigen Gefäßen in der Speiseröhre
kommen. Diese Gefäße können massive Blutungen verursachen
die tödlich enden können. 

Je nach Stärke  äußern sich die Blutungen durch
schwallartiges Erbrechen von Kaffeesatz-artigem
Mageninhalt\footnote{\T{Kaffeesatz-artiges Erbrechen}:
Durch die Magensäure kommt es zur Veränderung des Blutes,
es wird bräunlich, ähnlich dem Kaffeesatz im Kaffeefilter.}
(wenn das Blut von der Magensäure bereits verändert wurde) oder von rotem, frischen Blut.
}
{
\begin{itemize}
    \item Oft bei Lebererkrankungen
    \item Hoher Blutverlust möglich -- Lebensgefahr!
\end{itemize}
}{}{}{}

\Layout{0}{\TxSymptome}
{
\begin{itemize}
  \item Bluterbrechen im Schwall (Hämatemesis)
\end{itemize}
}{
\begin{itemize}
  \item Bluterbrechen im Schwall (Hämatemesis)
\end{itemize}
}{}{}{}



\Fazit{Der Patient ist aufgrund des Blutverlustes akut
vital bedroht, es besteht Schockgefahr.}

\MassnahmenNG{Spezielle Maßnahmen}{Oesophagusvarizenblutung}{}{

\begin{itemize}
  \item \TxMassVitBes
  \begin{itemize}
    \item \ce{O2} 10-15\LMin
    \item Lagerung:
    Situationsgerecht \AuthNotes{Lagerung? Beine hoch?
    OK hoch?}
  \end{itemize}
  \item Allgemeine Maßnahmen der Schockbekämpfung
  (\prettyref{m:am-schock})
  \item Spezielle Maßnahmen bei Hypovolämischem Schock
  (\prettyref{topic:schock-hypovolaemischer})
  \item \Z{Notarzt: evtl. Sengstaken-Blakemore-Sonde}
  \item Hospitalisation: Abteilung für Chirurgie, Voranmeldung
\end{itemize}
}






\subsubsection{Magen-Darm-Grippe -- Gastroenteritis}

\emph{\Lang{Syn.} Brechdurchfall, \Z{Gastroduodenitis}}

Die Magen-Darm-Grippe ist eine chronisch oder akut
verlaufende Entzündung der Magen- und der Dünndarmschleimhaut. 

\Layout{0}{\TxBeschreibung{} und Ursachen}
{
Die Gastroenteritis wird durch Erreger
wie Bakterien oder Viren hervorgerufen. Es gibt völlig
harmlose aber auch lebensbedrohliche Verlaufsformen, dies ist abhängig vom Erreger und
letztendlich auch vom Alter des Patienten. Durch den
Durchfall und das Erbrechen kann es zu einem maßgeblichen
Verlust von Wasser und Elektrolyten kommen. Wenn der Patient
nicht ausreichend trinken kann besteht die Gefahr der Austrocknung (Exsikkose).
}{
\begin{itemize}
  \item Vielfältig.
  \item Häufig Infektionen.
\end{itemize}
}{}{}{}


\Layout{2}{\TxSymptome}
{
\begin{itemize}
    \item Erbrechen (Emesis)
    \item Durchfall (Diarrhoe)
    \item Wasser-/Elektrolytverlust
\end{itemize}
}{
\begin{itemize}
    \item Erbrechen (Emesis)
    \item Durchfall (Diarrhoe)
    \item Wasser-/Elektrolytverlust
\end{itemize}
}{}{}{}

\Layout{0}{Besonders gefährdet}
{%
Besonders betroffen davon sind Kleinkinder, da sie
normalerweise einen höheren Wasseranteil am Körpergewicht
haben.  Die Symptome sind klassisch: Oberbauchkrämpfe,
die ziemlich heftig sein können, Übelkeit, Erbrechen, Durchfälle.
Exsikkose-Gefahr! 
}{
\begin{itemize}
    \item (Klein-)Kinder
    \item alte Menschen
\end{itemize}
}{}{}{}



\MassnahmenNG{Spezielle Maßnahmen}{Gastroenteritis}{}{
\begin{itemize}
  \item Allgemeine Maßnahmen bei Abdominalerkrankungen
  (\prettyref{m:am-abdominalerkrankungen})
  \item Patient nüchtern lassen.
  \item Hospitalisierung: Abt. f. Innere Medizin; bei
  Kindern Abt. f. Kinderheilkunde.
  \item Strikte Desinfektionsmaßnahmen!
\end{itemize}
}







\subsubsection{Flüssigkeitsmangel -- Exsikkose}
\ACh{WITM}{????-??-??}{NEU: Schlagworte}
\ACh{GABG}{2009-05-11}{NEU: Fließtexte}
\ACh{GABS}{2009-05-14}{Fließtexte stark modifiziert und angepasst. }

\emph{\Lang{Syn.} Dehydratation}

\Definition{„Austrocknung“: Es besteht ein hochgradiges Defizit von Flüssigkeit im Körper.}

Grundsätzlich wird (wurde) entweder \E{zu wenig Flüssigkeit zugeführt}
oder \E{zu viel Flüssigkeit an die Außenwelt abgegeben} (eine Kombination aus
beidem ist natürlich auch möglich). Dafür gibt es jeweils sehr verschiedene Ursachen: 

\Layout{0}{Ursachen}
{%
Ältere Menschen haben oft ein mangelndes Durstgefühl und daher
einen verminderten Trinkantrieb. Auch durch die Einnahme von
Drogen (Ecstacy) kann auf den Durst "`vergessen"' werden.  

Andererseits kann es durch \E{Erbrechen}, \E{Durchfall}
sowie \E{starkes Schwitzen} zu einem stark erhöhten
Flüssigkeitsverlust kommen. Bei \E{Fieber} wird auch
vermehrt Feuchtigkeit abgeatmet, und natürlichen können
auch \E{Medikamente} "`entwässern"'.
}{
\begin{itemize} %GABS: umgebaut
  \item Flüssigkeitszufuhr ${\downarrow}$
  \begin{itemize}
    \item Verminderter Trinkantrieb
  \end{itemize}
  \item Flüssigkeitsverlust ${\uparrow}$
  \begin{itemize}
    \item Erbrechen, Durchfall
    \item Starkes Schwitzen
    \item Fieber (Feuchtigkeit wird vermehrt abgeatmet)
    \item (Entwässerungs)medikamente
  \end{itemize}
\end{itemize}
}{}{}{}

Die genannten Ursachen sind natürlich nur ein kleiner Teil von sehr vielen
Möglichkeiten.

\Layout{0}{\TxSymptome}
{%
Der Patient wirkt allgemein \E{geschwächt bis apathisch},
oft \Es{verwirrt}, klagt über mehr oder minder heftiges
Durstgefühl, ev. über Wadenkrämpfe und Erbrechen. Die Haut
ist auffallend trocken, \Es{Falten sind abhebbar und
bleiben bestehen}. \Es{Mundhöhle und Zunge sind trocken},
evtl. borkig oder mit eingetrocknetem Schleim belegt. Die Augen sind charakteristischerweise tiefliegend, die Nase wirkt auffallend spitz. Der RR ist mitunter sehr nieder,
die HF eher hoch. \E{Der Patient ist bis zur Klärung der
Situation nüchtern zu belassen!}
}{
\begin{itemize}
    \item RR ${\downarrow}$
    \item stehende Hautfalten
    \item trockene, belegte Zunge
    \item Elektrolytentgleisung
    \item \Es{Verwirrtheit, Schwindel, Apathie}
\end{itemize}
}{}{}{}


% \par
% \hrule
% \par
% 
%
% \subsubsection{Flüssigkeitsmangel -- Exsikkose}
% \ACh{WITM}{????-??-??}{NEU: Schlagworte}
% \ACh{GABG}{2009-05-11}{NEU: Fließtexte}
% \ACh{GABS}{2009-05-14}{Fließtexte stark modifiziert und angepasst. }
% 
% \emph{\Lang{Syn.} Dehydratation}
% 
% \Definition{„Austrocknung“: Es besteht ein hochgradiges Defizit von Flüssigkeit im Körper.}
% 
% 
% \TagPar{Ursachen}Grundsätzlich wird (wurde) entweder
% \Es{zu wenig Flüssigkeit zugeführt} oder \Es{zu viel Flüssigkeit an die
% Außenwelt abgegeben} (eine Kombination aus beidem ist natürlich auch möglich). Dafür gibt es jeweils sehr verschiedene Ursachen: 
% 
% Ältere Menschen haben oft ein mangelndes Durstgefühl und daher
% einen verminderten Trinkantrieb. Auch durch die Einnahme von
% Drogen (Ecstacy) kann auf den Durst "`vergessen"' werden.  
% 
% Andererseits kann es durch \E{Erbrechen}, \E{Durchfall}
% sowie \E{starkes Schwitzen} zu einem stark erhöhtem
% Flüssigkeitsverlust kommen. Bei \E{Fieber} wird auch
% vermehrt Feuchtigkeit abgeatmet, und natürlichen können
% auch \E{Medikamente} "`entwässern"'.
% 
% 
% 
% Die genannten Ursachen sind natürlich nur ein kleiner Teil von sehr vielen
% Möglichkeiten.
% 
% \TagPar{\TxSymptome}Der Patient wirkt allgemein \E{geschwächt
% bis apathisch}, oft \Es{verwirrt}, klagt über mehr oder minder heftiges
% Durstgefühl, ev. über Wadenkrämpfe und Erbrechen. Die Haut
% ist auffallend trocken, \Es{Falten sind abhebbar und
% bleiben bestehen}. \Es{Mundhöhle und Zunge sind trocken},
% evtl. borkig oder mit eingetrocknetem Schleim belegt. Die Augen sind charakteristischerweise tiefliegend, die Nase wirkt auffallend spitz. Der RR ist mitunter sehr nieder,
% die HF eher hoch. \E{Der Patient ist bis zur Klärung der
% Situation nüchtern zu belassen!} 
% 
% 
% \subsubsection{Flüssigkeitsmangel -- Exsikkose}
% \ACh{WITM}{????-??-??}{NEU: Schlagworte}
% \ACh{GABG}{2009-05-11}{NEU: Fließtexte}
% \ACh{GABS}{2009-05-14}{Fließtexte stark modifiziert und angepasst. }
% 
% \emph{\Lang{Syn.} Dehydratation}
% \A{Version mit multicols}
% 
% \Definition{„Austrocknung“: Es besteht ein hochgradiges Defizit von Flüssigkeit im Körper.}
% 
% \begin{multicols}{2}
% \TagPar{Ursachen}Grundsätzlich wird (wurde) entweder
% \Es{zu wenig Flüssigkeit zugeführt} oder \Es{zu viel Flüssigkeit an die
% Außenwelt abgegeben} (eine Kombination aus beidem ist natürlich auch möglich). Dafür gibt es jeweils sehr verschiedene Ursachen: 
% 
% Ältere Menschen haben oft ein mangelndes Durstgefühl und daher
% einen verminderten Trinkantrieb. Auch durch die Einnahme von
% Drogen (Ecstacy) kann auf den Durst "`vergessen"' werden.  
% 
% Andererseits kann es durch \E{Erbrechen}, \E{Durchfall}
% sowie \E{starkes Schwitzen} zu einem stark erhöhtem
% Flüssigkeitsverlust kommen. Bei \E{Fieber} wird auch
% vermehrt Feuchtigkeit abgeatmet, und natürlichen können
% auch \E{Medikamente} "`entwässern"'.
% 
% 
% 
% Die genannten Ursachen sind natürlich nur ein kleiner Teil von sehr vielen
% Möglichkeiten.
% 
% \TagPar{\TxSymptome}Der Patient wirkt allgemein \E{geschwächt
% bis apathisch}, oft \Es{verwirrt}, klagt über mehr oder minder heftiges
% Durstgefühl, ev. über Wadenkrämpfe und Erbrechen. Die Haut
% ist auffallend trocken, \Es{Falten sind abhebbar und
% bleiben bestehen}. \Es{Mundhöhle und Zunge sind trocken},
% evtl. borkig oder mit eingetrocknetem Schleim belegt. Die Augen sind charakteristischerweise tiefliegend, die Nase wirkt auffallend spitz. Der RR ist mitunter sehr nieder,
% die HF eher hoch. \E{Der Patient ist bis zur Klärung der
% Situation nüchtern zu belassen!} 
% \end{multicols}
% 
% 
% \subsubsection{Flüssigkeitsmangel -- Exsikkose}
% \ACh{WITM}{????-??-??}{NEU: Schlagworte}
% \ACh{GABG}{2009-05-11}{NEU: Fließtexte}
% \ACh{GABS}{2009-05-14}{Fließtexte stark modifiziert und angepasst. }
% 
% \emph{\Lang{Syn.} Dehydratation}
% 
% \Definition{„Austrocknung“: Es besteht ein hochgradiges Defizit von Flüssigkeit im Körper.}
% 
% 
% \Married{
% \TagPar{Ursachen}Grundsätzlich wird (wurde) entweder
% \Es{zu wenig Flüssigkeit zugeführt} oder \Es{zu viel Flüssigkeit an die
% Außenwelt abgegeben} (eine Kombination aus beidem ist natürlich auch möglich). Dafür gibt es jeweils sehr verschiedene Ursachen: 
% 
% Ältere Menschen haben oft ein mangelndes Durstgefühl und daher
% einen verminderten Trinkantrieb. Auch durch die Einnahme von
% Drogen (Ecstacy) kann auf den Durst "`vergessen"' werden.  
% 
% Andererseits kann es durch \E{Erbrechen}, \E{Durchfall}
% sowie \E{starkes Schwitzen} zu einem stark erhöhtem
% Flüssigkeitsverlust kommen. Bei \E{Fieber} wird auch
% vermehrt Feuchtigkeit abgeatmet, und natürlichen können
% auch \E{Medikamente} "`entwässern"'.
% 
% 
% 
% Die genannten Ursachen sind natürlich nur ein kleiner Teil von sehr vielen
% Möglichkeiten.
% }{
% \TagPar{\TxSymptome}Der Patient wirkt allgemein \E{geschwächt
% bis apathisch}, oft \Es{verwirrt}, klagt über mehr oder minder heftiges
% Durstgefühl, ev. über Wadenkrämpfe und Erbrechen. Die Haut
% ist auffallend trocken, \Es{Falten sind abhebbar und
% bleiben bestehen}. \Es{Mundhöhle und Zunge sind trocken},
% evtl. borkig oder mit eingetrocknetem Schleim belegt. Die Augen sind charakteristischerweise tiefliegend, die Nase wirkt auffallend spitz. Der RR ist mitunter sehr nieder,
% die HF eher hoch. \E{Der Patient ist bis zur Klärung der
% Situation nüchtern zu belassen!} 
% }
% 




\subsection{Weniger häufige oder nicht-so-gut-erkennbare Erkrankungen}

\minisec{Querverweise}

\begin{itemize}
    \item Hepatitis: \prettyref{topic:hepatitis}
\end{itemize}




\subsubsection{Magen- und Zwölffingerdarmgeschwür}

\Z{\emph{Syn. Magengeschwür: Gastritis}}

\Definition{Magengeschwür: Entzündung der
Magenschleimhaut mit Substanzdefekt.}

\Layout{2}{Ursachen}
{
\begin{itemize}
    \item Bakteriell \Z{(Helicobacter pylori)} \cite{Malfertheiner2004}
    \item Medikamente\footnote{Nicht-steroidale Antirheumatika (NSAR), wie z.\,B.
    Voltaren{\texttrademark} (Diclofenac), Aspirin{\texttrademark}
    (Acetylsalicylsäure, ASS),etc., schädigen die Magenschleimhaut.}
    \item Chemisch-toxisch.
    \item \Z{Gallensäurenrückfluß (Reflux)}
    \item \Z{Autoimmunerkrankungen}
\end{itemize}
}{
\begin{itemize}
    \item Bakteriell \Z{(Helicobacter pylori)}
    \item Medikamente
    \item Chemisch-toxisch.
    \item \Z{Gallensäurenrückfluß (Reflux)}
    \item \Z{Autoimmunerkrankungen}
\end{itemize}
}{}{}{}

\Layout{0}{\TxSymptome}
{\AuthNotes{Fließtext fehlt: insbes. \underline{Wie} ist
der Schmerz, wann}
Es kommt zu Appetitlosigkeit, Völlegefühl, Übelkeit. Wenn
gleichzeitig eine Blutung besteht, dann sind Teerstühle
oder kaffeesatz-artiges Erbrechen zu beobachten.
}
{
\begin{itemize}
    \item Schmerzen (Oberbauch), Übelkeit, Erbrechen
    \item bei Blutung: Kaffesatz-Erbrechen,  Teerstuhl
\end{itemize}
}{}{}{}









\subsubsection{Bauchfellentzündung}
\label{topic:bauchfellentzuendung}
\label{topic:peritonitis}

\Z{\emph{\Lang{Syn.} Peritonitis}}

\Definition{Entzündung und Keimbesiedelung des 
Bauchfells}

\Layout{0}{\TxBeschreibung}{
Durch die Infektion und Entzündung des Bauchfelles kommt es
zu Symptomen eines Akuten Abdomens. Oft ist die Perforation
eines Hohlorganes  (Magendurchbruch, Darmperforation,
Blinddarmdurchbruch, \ldots) Ursache der Entzündung. 
}{

}{}{}{}

\Layout{0}{Ursache}{
\begin{itemize}
    \item Perforation eines Hohlorganes
    \item Infektion über Blut- / Lymphwege oder eingewandert
    über Fremdkörper
\end{itemize}
}{
\begin{itemize}
    \item Perforation eines Hohlorganes
    \item Infektion über Blut- / Lymphwege oder eingewandert
    über Fremdkörper
\end{itemize}
}{}{}{}

\Layout{0}{\TxSymptome}{
\begin{itemize}
    \item Symptome eines Akuten Abdomens
    \item Schmerzen (diffus)
    \item evtl. reflektorischer Darmverschluss
    \item evtl. Sepsis- und/oder Schockzeichen
\end{itemize}
}{
\begin{itemize}
    \item Symptome eines Akuten Abdomens
    \item Schmerzen (diffus)
    \item evtl. reflektorischer Darmverschluss
    \item evtl. Sepsis- und/oder Schockzeichen
\end{itemize}
}{}{}{}



% \Z{
% \subsection{Reserviert \A{Sonstige -- nicht akute --
% Erkrankungen}}
% \subsubsection{Reserviert \A{Chronisch entzündliche
% Darmerkrankungen}} }
% 



\section{Urologische Erkrankungen}

\Definition{Erkrankungen des harnableitenden Systems.}


\subsection{Nierenkolik}
\label{topic:nierenkolik}
\index{Kolik!Nieren-}
\index{Nierenkolik}

\Definition{Blockade des Harnleiters durch Nierensteine
mit massiven kolikartigen Schmerzen aufgrund des Harnstaus.}

\Z{\emph{Vgl. Kolik: \ref{topic:kolik}}}

\Layout{0}{Ursachen}{
Die Nierenkolik ist eine äußerst schmerzhafte
akute Erkrankung. Sie entsteht dadurch, dass Nierensteine
den Harnleiter blockieren und verstopfen. Dadurch wird der Harn
gestaut und der Harnleiter gedehnt. 
}{
\begin{itemize}
    \item Nierenstein(e)
\end{itemize}
}{}{}{}

\Layout{0}{\TxSymptome}{
Die Folge sind massive,
kolikartige Schmerzen, in der Flanke bzw. im Unterbauch der
jeweiligen Seite. Im Gegensatz zu anderen abdominellen
Erkrankungen ist der Patient extrem unruhig und agitiert.
Oft läuft er gestikulierend und wehklagend durch die
Gegend. Mitunter ist eine entsprechende Vorerkrankung
\E{(Nierensteine)} bereits bekannt und kann erfragt werden. 
}{
\begin{itemize}
    \item Rücken-/Flankenschmerzen
    \item Evtl. Unterbauchschmerzen bei  Stein-/Sandabgang
    \item "`Der Stein wandert, der Schmerz wandert, der Patient
    wandert"'
    \item Blut im Harn (Hämaturie)
\end{itemize}
}{}{}{}

\Layout{0}{Versorgung}{
Im Vordergrund bei der Versorgung
steht der rasche zügige Transport an eine Abteilung für Urologie. Wenn die
Schmerzen sehr massiv sind, muss eventuell eine
Schmerztherapie durch den Notarzt erfolgen.
}
{\SymbolText}{}{}{}


\MassnahmenNG{Spezielle Maßnahmen}{Nierenkolik}{}{
\begin{itemize}
  \item Lagerung nach Wunsch des Patienten
  \item schonender Transport
  \item Bei unerträglichen Schmerzen Notarzt zur
  Schmerztherapie beiziehen
  \item Hospitalisation: Abt. für Urologie
\end{itemize}
}


\subsection{Nierenbeckenentzündung}
\label{topic:nierenbeckenentzuendung}

{\TakeNotesLarge}
% TODO INHALT: fehlt komplett: Nierenbeckenentzündung 

\subsection{Harnwegsinfekt}
\label{topic:harnwegsinfekt}

{\TakeNotesLarge}
% TODO INHALT: fehlt komplett: Harnwegsinfekt 

\subsection{Akutes Harnverhalten}

\Definition{Abflussbehinderung des Harns aus der
Harnblase, verschiedene Ursachen sind möglich.}



Das akute Harnverhalten ist besonders bei Männern im "`besten"'
oder fortgeschrittenem Alter häufig anzutreffen. 

\Layout{0}{Ursachen}
{
Hauptursache  dafür ist eine
altersbedingte Vergrößerung der Prostata, die die Harnröhre abschnürt. Davon abgesehen
können auch alte Narben, Medikamente oder eine Störung der
Schließmuskel der Harnblase\footnote{zum Beispiel im Rahmen eines
Querschnitt-Syndroms nach Wirbelsäulenverletzungen} zu
Harnverhalten führen. 
}
{
\begin{itemize}
     \item Vergrößerung der \Es{Prostata}
    \item Lähmung (Querschnittlähmung)
    \item Schließmuskel-Krampf
    \item Verlegung der Harnröhre
    \item Medikamente
\end{itemize}
}{}{}{}

Je nach Füllungszustand der Harnblase kann das
Harnverhalten sehr schmerzhaft sein.

Die Maßnahmen beschränken sich auf einen zügigen Transport
an eine Abteilung für Urologie. Eine Notarzt-Nachforderung
ist kaum sinnvoll, da ein \Abk{NEF} oder \Abk{NAW} selten geeignetes
Material zur Harnableitung mitführt. 

\MassnahmenNG{Spezielle Maßnahmen}{Akutes Harnverhalten}{}
{
\begin{itemize}
    \item Zügiger Transport in eine Urologische Ambulanz
    \item \Z{vorsichtig ablassen über Harnkatheter oder
    Blasenpunktion (Spital)}
\end{itemize}
}




\subsection{Niereninsuffizienz und Nierenversagen}
\label{topic:cni-therapie-dialyse}
\label{topic:dialysepatienten-gefahren}

\AuthNotes{Abkz. CNI}

Meistens kommt es im Rahmen anderer Krankheiten
(Diabetes mellitus, Hypertonie, erhöhte Blutfette \ldots) zu einer massiven
Schädigung der kleinen Nierengefäße und in weiterer Folge zum Nierenversagen. Da
die zuvor erwähnten Krankheiten sehr häufig sind, ist in weiterer Folge die
Niereninsuffizienz ebenfalls häufig anzutreffen. 

Als Therapie steht der
Medizin anfänglich eine medikamentöse Behandlung, später jedoch nur die
Nierentransplantation oder die \E{Blutwäsche}
(\T{Dialyse}\index{Dialyse}) zur Verfügung. Bei der
klassischen Dialyse wird der Patient jeden zweiten oder
dritten Tag für einige Stunden an eine Maschine angeschlossen, welche sein
Blut reinigen kann. Die "`Dialysefahrten"' sind eine der
Hauptaufgaben des Krankentransportdienstes. 

\index{Dialyse-Patient!Gefahren}
\Es{Der Ausfall der normalen Nierenfunktion hat
Auswirkungen} auf den Wasserhaushalt, den
Elektrolythaushalt und den Säurebasenhaushalt. Diese Funktionen der Niere sind zum Überleben elementar (Vitalfunktionen zweiter
Ordnung). Dialysepatienten müssen daher sehr darauf achten
nicht zu viel zu trinken. Weiters können \E{Elektrolytstörungen} sogar
\E{lebensbedrohliche Notfälle} nach sich ziehen (bis hin
zum Herzstillstand).

Es ist daher wichtig auf den
Dialysepatienten -- auch im "`normalem"' Krankentransport --
immer ein wachsames Auge zu haben.







% \subsection{Hodentorsion}
% 
% Verdrehung eine Hodens im Hodensack  ={\textgreater}
% 
% Abschnürung von Samenstrang und  Blutgefäßen
% Oft nach Eintritt in Pubertät, ev. {\textgreater} 20 J.
% 
% \paragraph{Symptome:}
% 
% {\mdseries
% akuter heftiger Schmerz}
% 
% {\mdseries
% Schwellung}
% 
% \paragraph{Therapie:}
% 
% hochlagern
% 
% umgehende urologisch-chirurgische  Intervention (ad URO!!!)




% \clearpage





%\putbib
% Rea während Transport
\nocite{Eisenburger2008,Havel2007}
% General
\nocite{SiewertChirurgie8}
\nocite{AMLS3}
\nocite{Erc2005DeSec5}
\nocite{Erc2005Sec5}
\nocite{Erc2005DeSec6}
\nocite{Erc2005DeSec3}
\nocite{Erc2005Sec3}
\nocite{Juurlink2005}
\nocite{Kalla2006}

\nocite{Vlcek2008}
\nocite{Wells2000}
\nocite{BoehmerTaschRett6}
\nocite{Fauci2008}
\nocite{Lissauer2007}

\nocite{LPN3,Boecker2}
\nocite{Helfen2008}
\nocite{LippertAnatomie7}
\nocite{Nilsson1997}
\nocite{Patzelt2005}
\nocite{Sachsenweger2003}
\nocite{RedelsteinerHRNS1}
\nocite{ForMed2}
\nocite{StriebelAn7}
\nocite{Zeiler2001}
\nocite{KlinkePhysio3}
\nocite{DRAn3}
\nocite{ForthPharma8}
\nocite{LuellmannPharma16}
\nocite{ThierbachInterhospitaltransfer1}
\nocite{AnCompact3}
\nocite{RueckerBildatlas1}


\clearpage
\ChapterBibliography 

%\bibliographystyle{unsrtdin2}{\scriptsize \bibliography{Literatur-AAS}}
